% Lesson 28 — Speaking:Exchanges information about preventing communicable diseases in schools, workplaces, and the community — Anglais Tle A — Cours signé OnBuch+
% Programme officiel MINESEC. Compiler avec : tectonic EN28-speaking-exchanges-information-about-preventing-communi.tex
\newcommand{\DOCMATIERE}{Anglais}
\newcommand{\DOCNIVEAU}{Tle A}
\newcommand{\DOCMODULE}{Chapter 3 : Environment, Well-being and Health}
\newcommand{\DOCLECON}{EN28}
\newcommand{\DOCTITRE}{Lesson 28 — Speaking:Exchanges information about preventing communicable diseases in schools, workplaces, and the community}
% OnBuch+ — préambule « cours signés » ENRICHI (compilé avec tectonic / XeLaTeX)
% Système visuel « Pop » d'OnBuch+ : fond crème (jamais blanc), bordures encre
% épaisses, ombres « dures » décalées, titres Archivo Black en capitales.
% Version MATHÉMATIQUES : graphes (pgfplots/tikz), boîtes pédagogiques
% (définition, propriété, méthode, attention, le savais-tu, exercice résolu,
% exercice, TP, bilan), page de garde et signature OnBuch+ ancrée en bas.
%
% Macros attendues dans main.tex AVANT \input{preamble} :
%   \DOCMATIERE  \DOCNIVEAU  \DOCMODULE  \DOCLECON (numéro)  \DOCTITRE
\documentclass[11pt]{article}

\usepackage{fontspec}
\usepackage[english,french]{babel} % français = langue principale ; l'anglais via \eng{..} / textebox
\usepackage[a4paper,margin=2cm,top=3.4cm,bottom=2.8cm,headheight=1.4cm,footskip=1.3cm]{geometry}
\usepackage{amsmath,amssymb,amsfonts}
\usepackage{mathtools} % \xmapsto, \coloneqq... souvent utilisés par les modèles
\usepackage{cancel} % simplifications barrées (bilans de forces)
% \abs{z} (module, valeur absolue) et \norm{u} (norme), utilisés par les modèles
\providecommand{\abs}{}\let\abs\relax\DeclarePairedDelimiter{\abs}{\lvert}{\rvert}
\providecommand{\norm}{}\let\norm\relax\DeclarePairedDelimiter{\norm}{\lVert}{\rVert}
\usepackage[table]{xcolor} % [table] : fournit aussi \rowcolors (alternance de lignes)
\usepackage{pagecolor}
\usepackage{tikz}
\usetikzlibrary{arrows.meta,positioning,calc,shapes.geometric,shapes.misc,shapes.arrows,decorations.pathmorphing,decorations.pathreplacing,decorations.markings,patterns,shadows,mindmap,trees,fit,backgrounds,matrix,intersections,angles,quotes}
\usepackage{pgfplots}
\usepackage[version=4]{mhchem} % équations nucléaires \ce{^{235}_{92}U}
\usepackage[european,siunitx]{circuitikz} % schémas électriques
\pgfplotsset{compat=1.18}
\usepgfplotslibrary{fillbetween,groupplots} % aires entre courbes (calcul intégral)
\usepackage{enumitem}
\usepackage{fancyhdr}
% [most] charge toutes les bibliothèques tcolorbox (listings, minted, raster,
% documentation...) et gonfle inutilement la mémoire TeX sur un document long
% (~300 boîtes) : on ne charge que ce qui est réellement utilisé.
\usepackage[skins,breakable]{tcolorbox}
\usepackage{graphicx}
\usepackage{colortbl}
\usepackage{booktabs}
\usepackage{tabularx}
\usepackage{diagbox} % case d'en-tête diagonale des tableaux à double entrée
% Colonnes extensibles centrée / gauche / droite (C, L, R), souvent utilisées par les modèles
\newcolumntype{C}{>{\centering\arraybackslash}X}
\newcolumntype{L}{>{\raggedright\arraybackslash}X}
\newcolumntype{R}{>{\raggedleft\arraybackslash}X}
\usepackage{array}
\usepackage{multicol}
\usepackage{multirow}
\usepackage{siunitx}
% parse-numbers=false : accepte aussi \SI{2,5 \times 10^{-3}}{...}
\sisetup{locale=FR,output-decimal-marker={,},group-minimum-digits=4,parse-numbers=false}
\DeclareSIUnit\ohm{\ensuremath{\symup{\Omega}}} % U+2126 absent de la police
\DeclareSIUnit\division{div} % réglages d'oscilloscope (ms/div, V/div)
\DeclareSIUnit\tr{tr} % tours (vitesse de rotation, ex. \SI{1500}{\tr\per\minute})
\DeclareSIUnit\molar{M} % concentration molaire (ex. \SI{3}{\molar}, \SI{1}{\micro\molar})
\DeclareSIUnit\an{an} % année (le modèle confond parfois avec \year, un compteur TeX) — ex. \SI{5}{\kilo\meter\per\an}
\DeclareSIUnit\FCFA{FCFA} % franc CFA (ex. \SI{500}{\FCFA\per\kilo\gram})
\DeclareSIUnit\degreeBrix{\textdegree Brix} % teneur en sucre d'un sirop/jus (réfractométrie)
\DeclareSIUnit\month{mois} % le modèle utilise parfois \month, un compteur TeX, comme unité de durée
\usepackage{qrcode}
% \degree : commande fréquemment utilisée par les modèles pour un angle ou une
% température en dehors de \SI{}{} (non fournie par les paquets chargés).
\providecommand{\degree}{^{\circ}}
% \qed (fin de démonstration), utilisé par les modèles sans amsthm
\providecommand{\qed}{\hfill$\square$}
% \barre : double barre des valeurs interdites dans un tableau de variations (tkz-tab)
\providecommand{\barre}{\ensuremath{\|}}
% environnement proof (amsthm non chargé) utilisé par les modèles
\ifdefined\proof\else\newenvironment{proof}[1][Démonstration]{\par\noindent\textit{#1.}\ }{\qed\par}\fi
\usepackage{lastpage}
\usepackage{hyperref}
\hypersetup{hidelinks}

% ============================================================
% Palette « Pop » OnBuch+ — mêmes hex que la classe P (plus_ui.dart)
% ============================================================
\definecolor{poporange}{HTML}{F59321}
\definecolor{popdark}{HTML}{E07A0C}
\definecolor{popcream}{HTML}{FFF6E8}
\definecolor{popink}{HTML}{1C1714}
\definecolor{poppurple}{HTML}{7A5AE0}
\definecolor{popgreen}{HTML}{2E9E6A}
\definecolor{poppink}{HTML}{E0457B}
\definecolor{popblue}{HTML}{2F7DE1}
\definecolor{popgold}{HTML}{C98A12}
\definecolor{popmuted}{HTML}{8A7F76}
\definecolor{popline}{HTML}{EADFD2}
\definecolor{popgray}{HTML}{8A7F76}
\colorlet{popgrey}{popgray}
% Alias tolérants (le modèle nomme parfois les couleurs différemment)
\colorlet{popwhite}{white}
\colorlet{popblack}{popink}
\colorlet{popred}{poppink}
\colorlet{popyellow}{popgold}
% Teintes claires (fonds de tableaux/encadrés) — le modèle nomme parfois
% les couleurs avec un suffixe L (ex. popblueL) sans les avoir définies.
\colorlet{poporangeL}{poporange!20}
\colorlet{popdarkL}{popdark!20}
\colorlet{poppurpleL}{poppurple!20}
\colorlet{popgreenL}{popgreen!20}
\colorlet{poppinkL}{poppink!20}
\colorlet{popblueL}{popblue!20}
\colorlet{popgoldL}{popgold!20}
\colorlet{popredL}{popred!20}
\colorlet{popgrayL}{popgray!20}
% Teintes claires pour fonds de tableaux / zones
\colorlet{popblueL}{popblue!10!white}
\colorlet{poporangeL}{poporange!14!white}
\colorlet{popgreenL}{popgreen!12!white}
\colorlet{poppurpleL}{poppurple!12!white}
\colorlet{poppinkL}{poppink!12!white}
\colorlet{popgoldL}{popgold!14!white}

\pagecolor{popcream}

% ============================================================
% Typographie
% ============================================================
\defaultfontfeatures{Ligatures=TeX}
\setmainfont{PlusJakartaSans-Medium.ttf}[Path=../../../fonts/,BoldFont=PlusJakartaSans-Bold.ttf]
% Symboles, API et emojis absents de la police principale : repli sur DejaVu Sans (caractères actifs)
\newfontface\symfont{DejaVuSans.ttf}[Path=../../../fonts/]
\makeatletter
\newcommand\symactive[1]{\catcode#1=13 \begingroup\lccode`\~=#1 \lowercase{\endgroup\def~}{{\symfont\char#1}}}
\newcommand\symblank[1]{\catcode#1=13 \begingroup\lccode`\~=#1 \lowercase{\endgroup\def~}{}}
\newcommand\symhyph[1]{\catcode#1=13 \begingroup\lccode`\~=#1 \lowercase{\endgroup\def~}{\mbox{-}}}
\newcount\symc
\newcommand\symrange[2]{\symc=#1 \@whilenum\symc<#2 \do{\expandafter\symactive\expandafter{\the\symc}\advance\symc by 1 }}
\newcommand\symblankrange[2]{\symc=#1 \@whilenum\symc<#2 \do{\expandafter\symblank\expandafter{\the\symc}\advance\symc by 1 }}
\makeatother
\symrange{"0250}{"02FF}\symactive{"2713}\symactive{"2714}\symactive{"2717}\symactive{"2718}\symactive{"2610}\symactive{"2611}\symactive{"2612}
\symhyph{"2011}\symhyph{"2E1A}\symblankrange{"1F300}{"1FAFF}\symblank{"FE0F}
% Maths en sans-serif (Fira Math) avec chiffres et lettres droites de Plus
% Jakarta, pour que les formules se fondent dans le texte.
\usepackage[math-style=ISO,bold-style=ISO]{unicode-math}
\setmathfont{FiraMath-Regular.otf}[Path=../../../fonts/]
\setmathfont{PlusJakartaSans-Medium.ttf}[Path=../../../fonts/,range={up/num,up/{Latin,latin},\mathup}]
\usepackage{icomma} % virgule décimale sans espace en mode math
% Caractères mathématiques Unicode absents des polices de texte (Plus Jakarta,
% Archivo Black) : rendus par leur équivalent mathématique, en texte comme en
% formule — sinon ils s'affichent en carré vide (ex. « Arithmétique dans ℤ »).
\usepackage{newunicodechar}
\newunicodechar{ℝ}{\ensuremath{\mathbb{R}}}
\newunicodechar{ℤ}{\ensuremath{\mathbb{Z}}}
\newunicodechar{ℂ}{\ensuremath{\mathbb{C}}}
\newunicodechar{ℕ}{\ensuremath{\mathbb{N}}}
\newunicodechar{ℚ}{\ensuremath{\mathbb{Q}}}
\newunicodechar{𝔻}{\ensuremath{\mathbb{D}}}
\newunicodechar{α}{\ensuremath{\alpha}}
\newunicodechar{β}{\ensuremath{\beta}}
\newunicodechar{γ}{\ensuremath{\gamma}}
\newunicodechar{δ}{\ensuremath{\delta}}
\newunicodechar{ε}{\ensuremath{\varepsilon}}
\newunicodechar{θ}{\ensuremath{\theta}}
\newunicodechar{λ}{\ensuremath{\lambda}}
\newunicodechar{μ}{\ensuremath{\mu}}
\newunicodechar{σ}{\ensuremath{\sigma}}
\newunicodechar{τ}{\ensuremath{\tau}}
\newunicodechar{φ}{\ensuremath{\varphi}}
\newunicodechar{ψ}{\ensuremath{\psi}}
\newunicodechar{ω}{\ensuremath{\omega}}
\newunicodechar{Σ}{\ensuremath{\Sigma}}
% majuscules produites par \MakeUppercase dans les titres (α-aminés -> Α)
\newunicodechar{Α}{\ensuremath{\alpha}}
\newunicodechar{Β}{\ensuremath{\beta}}
\newunicodechar{Γ}{\ensuremath{\Gamma}}
\newunicodechar{Δ}{\ensuremath{\Delta}}
\newunicodechar{Ε}{\ensuremath{\varepsilon}}
\newunicodechar{Θ}{\ensuremath{\Theta}}
\newunicodechar{Λ}{\ensuremath{\Lambda}}
\newunicodechar{Μ}{\ensuremath{\mu}}
\newunicodechar{Τ}{\ensuremath{\tau}}
\newunicodechar{Φ}{\ensuremath{\Phi}}
\newunicodechar{Ψ}{\ensuremath{\Psi}}
\newunicodechar{Ω}{\ensuremath{\Omega}}
\newunicodechar{↦}{\ensuremath{\mapsto}}
\newunicodechar{∈}{\ensuremath{\in}}
\newunicodechar{∉}{\ensuremath{\notin}}
\newunicodechar{∧}{\ensuremath{\wedge}}
\newunicodechar{⇔}{\ensuremath{\Leftrightarrow}}
\newunicodechar{⟺}{\ensuremath{\Longleftrightarrow}}
\newunicodechar{⇒}{\ensuremath{\Rightarrow}}
\newunicodechar{⟹}{\ensuremath{\Longrightarrow}}
\newunicodechar{∘}{\ensuremath{\circ}}
\newunicodechar{∪}{\ensuremath{\cup}}
\newunicodechar{∩}{\ensuremath{\cap}}
\newunicodechar{≡}{\ensuremath{\equiv}}
\newunicodechar{⊂}{\ensuremath{\subset}}
\newunicodechar{⊥}{\ensuremath{\perp}}
\newunicodechar{∃}{\ensuremath{\exists}}
\newunicodechar{∀}{\ensuremath{\forall}}
\newunicodechar{∛}{\ensuremath{\sqrt[3]{\,}}}
\newunicodechar{∜}{\ensuremath{\sqrt[4]{\,}}}
\newunicodechar{⃗}{\ensuremath{{}^{\to}}}
\newunicodechar{ⁿ}{\ifmmode^{n}\else\textsuperscript{\ensuremath{n}}\fi}
\newunicodechar{ˣ}{\ifmmode^{x}\else\textsuperscript{\ensuremath{x}}\fi}
\newunicodechar{ᵇ}{\ifmmode^{b}\else\textsuperscript{\ensuremath{b}}\fi}
\newunicodechar{ʳ}{\ifmmode^{r}\else\textsuperscript{\ensuremath{r}}\fi}
\newunicodechar{ᵘ}{\ifmmode^{u}\else\textsuperscript{\ensuremath{u}}\fi}
\newunicodechar{ʸ}{\ifmmode^{y}\else\textsuperscript{\ensuremath{y}}\fi}
\newunicodechar{ᵃ}{\ifmmode^{a}\else\textsuperscript{\ensuremath{a}}\fi}
\newunicodechar{ᵉ}{\ifmmode^{e}\else\textsuperscript{\ensuremath{e}}\fi}
\newunicodechar{ᵏ}{\ifmmode^{k}\else\textsuperscript{\ensuremath{k}}\fi}
\newunicodechar{ᵖ}{\ifmmode^{p}\else\textsuperscript{\ensuremath{p}}\fi}
\newunicodechar{ᴺ}{\ifmmode^{N}\else\textsuperscript{\ensuremath{N}}\fi}
\newunicodechar{ᶜ}{\ifmmode^{c}\else\textsuperscript{\ensuremath{c}}\fi}
\newunicodechar{ʰ}{\ifmmode^{h}\else\textsuperscript{\ensuremath{h}}\fi}
\newunicodechar{ᵠ}{\ifmmode^{q}\else\textsuperscript{\ensuremath{q}}\fi}
\newunicodechar{ₙ}{\ifmmode_{n}\else\textsubscript{\ensuremath{n}}\fi}
\newunicodechar{ₐ}{\ifmmode_{a}\else\textsubscript{\ensuremath{a}}\fi}
\newunicodechar{ᵢ}{\ifmmode_{i}\else\textsubscript{\ensuremath{i}}\fi}
\newunicodechar{ᵣ}{\ifmmode_{r}\else\textsubscript{\ensuremath{r}}\fi}
\newunicodechar{ₖ}{\ifmmode_{k}\else\textsubscript{\ensuremath{k}}\fi}
\newunicodechar{ᵦ}{\ifmmode_{\beta}\else\textsubscript{\ensuremath{\beta}}\fi}
\newunicodechar{ₓ}{\ifmmode_{x}\else\textsubscript{\ensuremath{x}}\fi}
\newfontfamily\popdisplay{ArchivoBlack-Regular.ttf}[Path=../../../fonts/]
\newfontfamily\poplogo{SpaceGrotesk-Bold.ttf}[Path=../../../fonts/]
\newfontfamily\popsemi{PlusJakartaSans-SemiBold.ttf}[Path=../../../fonts/]
\newfontfamily\popxbold{PlusJakartaSans-ExtraBold.ttf}[Path=../../../fonts/]
\newfontfamily\popxboldls{PlusJakartaSans-ExtraBold.ttf}[Path=../../../fonts/,LetterSpace=3.0]

\setlist[enumerate]{leftmargin=1.7em,itemsep=5pt,topsep=4pt}
\setlist[itemize]{leftmargin=1.7em,itemsep=4pt,topsep=4pt,label={\color{poporange}\textbullet}}
\setlength{\parindent}{0pt}
\setlength{\parskip}{5pt}
\renewcommand{\arraystretch}{1.25}

% pgfplots : style Pop par défaut
\pgfplotsset{
  every axis/.append style={axis line style={popink,line width=1pt},tick style={popink},
    grid style={popline,line width=0.5pt},label style={font=\small},tick label style={font=\footnotesize}},
  popcurve/.style={poporange,line width=1.8pt,no markers,smooth},
}
% Même style disponible dans un \draw TikZ simple (hors pgfplots)
\tikzset{popcurve/.style={poporange,line width=1.8pt}}
\tikzset{popgrid/.style={popline,very thin}}
\colorlet{popcurve}{poporange} % le nom du style est parfois utilisé comme couleur

% ============================================================
% Pill / squiggle
% ============================================================
\newcommand{\poppill}[3]{%
  \tikz[baseline=-0.55ex]\node[draw=popink,line width=1.2pt,rounded corners=2.6pt,%
    fill=#2,inner xsep=6.5pt,inner ysep=3pt]{%
    \popxboldls\fontsize{7}{7}\selectfont\color{#3}\MakeUppercase{#1}};%
}
\newcommand{\popsquiggle}[1]{%
  \tikz[baseline=-0.4ex]{
    \draw[#1,line width=2.6pt,line cap=round]
      plot [smooth] coordinates {(0,0.09) (0.34,0.01) (0.68,0.21) (1.02,0.01) (1.36,0.10)};
  }%
}
% Mot-clé surligné (marqueur orange sous le texte)
\newcommand{\cle}[1]{{\popxbold\color{popdark}#1}}

% ============================================================
% En-tête / pied
% ============================================================
\pagestyle{fancy}
\fancyhf{}
\renewcommand{\headrulewidth}{0pt}
\renewcommand{\footrulewidth}{0pt}
\lhead{\raisebox{-0.15ex}{\poplogo\fontsize{13}{13}\selectfont\color{popink}OnBuch\color{poporange}+}}
\rhead{\poppill{\DOCMATIERE\ \textbullet\ \DOCNIVEAU\ \textbullet\ Leçon \DOCLECON}{white}{popink}}
\lfoot{\popxbold\fontsize{7.6}{7.6}\selectfont\color{popmuted}\MakeUppercase{Cours signé OnBuch+}\ \textbullet\ {\popsemi\DOCTITRE}}
\rfoot{\tikz[baseline=-0.6ex]\node[rounded corners=8.5pt,fill=popink,minimum height=17pt,minimum width=17pt,inner xsep=5pt,inner sep=0]{\popxbold\fontsize{8}{8}\selectfont\color{popcream}\thepage\,/\,\pageref*{LastPage}};}

% ============================================================
% Titres
% ============================================================
\newcommand{\chaptitre}[2]{%
  \begin{tcolorbox}[enhanced,colback=poporange,colframe=popink,boxrule=2.6pt,arc=11pt,
      drop shadow=popink,left=15pt,right=15pt,top=12pt,bottom=11pt,before skip=3pt,after skip=18pt]
    \poppill{#2}{popink}{popcream}\par\vspace{9pt}
    {\raggedright\hyphenpenalty=10000\exhyphenpenalty=10000\popdisplay\fontsize{22}{24}\selectfont\color{popcream}\MakeUppercase{#1}\par}
  \end{tcolorbox}
}
% Section numérotée : pastille orange avec le numéro + titre Archivo
\newcounter{popsec}
\newcommand{\coursec}[1]{%
  \stepcounter{popsec}\setcounter{popsub}{0}%
  \par\vspace{16pt}\noindent
  \begin{minipage}{\linewidth}
  \tikz[baseline=-0.7ex]\node[draw=popink,line width=1.6pt,fill=poporange,rounded corners=4pt,minimum size=22pt,inner sep=0]{\popdisplay\fontsize{12}{12}\selectfont\color{popcream}\thepopsec};%
  \hspace{8pt}{\popdisplay\fontsize{15}{17}\selectfont\color{popink}\MakeUppercase{#1}}\par
  \vspace{4pt}\noindent{\color{poporange}\rule{36pt}{3.6pt}}
  \end{minipage}\par\nopagebreak\vspace{8pt}
}
\newcounter{popsub}
\newcommand{\courssub}[1]{%
  \stepcounter{popsub}%
  \par\vspace{9pt}\noindent{\popxbold\fontsize{11.5}{13}\selectfont\color{popink}{\color{poporange}\thepopsec.\thepopsub}\hspace{6pt}#1}\par\nopagebreak\vspace{4pt}
}

% ============================================================
% Boîtes pédagogiques — même recette : fond blanc, bordure épaisse colorée,
% ombre dure encre, badge plein. Titre optionnel [#1].
% ============================================================
\tcbset{popbase/.style={breakable,enhanced,colback=white,boxrule=2.3pt,arc=9pt,
  shadow={3pt}{-3pt}{0pt}{popink},left=14pt,right=14pt,top=11pt,bottom=11pt,before skip=11pt,after skip=15pt,
  fontupper=\popsemi\fontsize{10.6}{14.5}\selectfont\color{popink},
  fontlower=\popsemi\fontsize{10.6}{14.5}\selectfont\color{popink}}}
% Coche dessinée (le glyphe ✓ n'existe pas dans les polices du document)
\newcommand{\popcheck}[1]{\tikz[baseline=-0.5ex]\draw[#1,line width=1.6pt,line cap=round,line join=round](-0.13,0.0)--(-0.04,-0.09)--(0.13,0.1);}
\providecommand{\coche}{\popcheck{popgreen}}
\providecommand{\resultat}[1]{\par\smallskip\noindent\fcolorbox{popgreen}{popgreenL}{\parbox{\dimexpr\linewidth-2\fboxsep-2\fboxrule\relax}{\textbf{Résultat.} #1}}\par\smallskip}
\newcommand{\popboxhead}[3]{\poppill{#1}{#2}{white}\ifx\relax#3\relax\else\hspace{6pt}{\popxbold\fontsize{10.6}{12}\selectfont\color{popink}#3}\fi\par\vspace{7pt}}

\newtcolorbox{aretenir}[1][]{popbase,colframe=popgreen,before upper={\popboxhead{À retenir}{popgreen}{#1}}}
% Texte en anglais (document de lecture, dialogue, extrait) : cadre
% d'étude. Un titre optionnel (source, auteur) s'affiche dans l'en-tête.
\newtcolorbox{textebox}[1][]{popbase,colframe=popblue,before upper={\popboxhead{Text}{popblue}{#1}\begin{otherlanguage}{english}},after upper={\end{otherlanguage}}}
% Marques ✓ / ✗ (forme correcte / fautive) souvent utilisées par les modèles
\providecommand{\cross}{\textcolor{poppink}{\ensuremath{\times}}}
\providecommand{\xmark}{\cross}\providecommand{\crossmark}{\cross}\providecommand{\cmark}{\ensuremath{\checkmark}}
% Ligne de réponse / justification à compléter (inventée par les modèles)
\providecommand{\justification}[1]{\par\noindent\textit{Justification~:} \dotfill\par}
% \textipa (package tipa non chargé) : la police ne couvre pas l'API -> simple passage
\providecommand{\textipa}[1]{#1}
% Soulignement (ulem non chargé) : \uline, \ul
\providecommand{\uline}[1]{\underline{#1}}\providecommand{\ul}[1]{\underline{#1}}
% \glossaire{\item ...} inventée par les modèles : liste de glossaire
\providecommand{\glossaire}[1]{\begin{itemize}#1\end{itemize}}
% \alert (beamer) inventée par les modèles : texte mis en évidence
\providecommand{\alert}[1]{\textbf{\textcolor{poppink}{#1}}}
% variantes ASCII d'ordinaux écrites en macro
\providecommand{\iere}{\textsuperscript{re}}\providecommand{\ieme}{\textsuperscript{e}}\providecommand{\ere}{\textsuperscript{re}}\providecommand{\eme}{\textsuperscript{e}}\providecommand{\er}{\textsuperscript{er}}\providecommand{\ie}{\textsuperscript{e}}
% Ordinaux français écrits en macro par les modèles : 1\ière, 3\ième
\newcommand{\ière}{\textsuperscript{re}}\newcommand{\ième}{\textsuperscript{e}}
% \exemple (inventée par les modèles) : étiquette « Exemple : »
\providecommand{\exemple}{\textbf{Exemple~:}\ }
% \square utilisable aussi hors mode math (cases à cocher)
\AtBeginDocument{\let\squareMath\square\renewcommand{\square}{\ensuremath{\squareMath}}}
% Mot ou phrase en anglais à l'intérieur d'un paragraphe français
\newcommand{\eng}[1]{\ifmmode\text{\textit{#1}}\else\begingroup\selectlanguage{english}\itshape #1\endgroup\fi}
\let\esp\eng % alias toléré
% flèches d'intonation inventées par les modèles
\providecommand{\textsearrow}{}\renewcommand{\textsearrow}{\ensuremath{\searrow}}\providecommand{\textnearrow}{}\renewcommand{\textnearrow}{\ensuremath{\nearrow}}
% \fr{..} : mot français (inventé par les modèles)
\providecommand{\fr}[1]{\textit{#1}}
\newtcolorbox{exemplebox}[1][]{popbase,colframe=poporange,before upper={\popboxhead{Exemple}{poporange}{#1}}}
\newtcolorbox{definition}[1][]{popbase,colframe=popblue,before upper={\popboxhead{Définition}{popblue}{#1}}}
\newtcolorbox{propriete}[1][]{popbase,colframe=poppurple,before upper={\popboxhead{Propriété}{poppurple}{#1}}}
\newtcolorbox{methode}[1][]{popbase,colframe=popgold,before upper={\popboxhead{Méthode}{popgold}{#1}}}
\newtcolorbox{attention}[1][]{popbase,colframe=poppink,before upper={\popboxhead{Attention}{poppink}{#1}}}
\newtcolorbox{experience}[1][]{popbase,colframe=popblue,colback=popblueL,before upper={\popboxhead{Expérience}{popblue}{#1}}}
% « Le savais-tu ? » : carte violette pleine (contexte, culture, Cameroun)
\newtcolorbox{savaistu}[1][]{popbase,colback=poppurple,colframe=popink,
  fontupper=\popsemi\fontsize{10.4}{14.2}\selectfont\color{white},
  before upper={\poppill{Le savais-tu\,?}{white}{poppurple}\ifx\relax#1\relax\else\hspace{6pt}{\popxbold\color{white}#1}\fi\par\vspace{7pt}}}
% Exercice résolu : énoncé en haut, \tcblower puis la solution sur fond crème
\newcounter{popexo}\newcounter{popexr}
\newtcolorbox{exoresolu}[1][]{popbase,colframe=poporange,colbacklower=popcream,
  segmentation style={popink,line width=1.2pt,dashed},
  before upper={\stepcounter{popexr}\popboxhead{Exercice résolu \thepopexr}{poporange}{#1}},
  before lower={\poppill{Solution}{popink}{popcream}\par\vspace{6pt}}}
% Exercice (non résolu en place) — la correction est dans la partie Corrigés
\newtcolorbox{exercice}[1][]{popbase,colframe=popink,
  before upper={\stepcounter{popexo}\popboxhead{Exercice \thepopexo}{popink}{#1}}}
% Corrigé
\newtcolorbox{corrige}[1][]{popbase,colframe=popgreen,colback=popgreenL,
  before upper={\popboxhead{Corrigé}{popgreen}{#1}}}
% Objectifs / prérequis / bilan
\newtcolorbox{objectifs}{popbase,colframe=popink,colback=poporangeL,
  before upper={\popboxhead{Objectifs de la leçon}{popink}{}}}
\newtcolorbox{prerequis}{popbase,colframe=popink,colback=popblueL,
  before upper={\popboxhead{Prérequis}{popblue}{}}}
\newtcolorbox{bilan}{popbase,colframe=popink,colback=white,boxrule=2.8pt,
  before upper={\popboxhead{Fiche bilan}{poporange}{L'essentiel à connaître pour l'examen}}}
% Figure encadrée avec légende
\newtcolorbox{popfigure}[1][]{enhanced,breakable=false,colback=white,colframe=popline,boxrule=1.4pt,arc=8pt,
  left=10pt,right=10pt,top=10pt,bottom=8pt,before skip=10pt,after skip=12pt,halign=center,
  lower separated=false,halign lower=center,
  fontlower=\popsemi\fontsize{9}{11.5}\selectfont\color{popmuted},
  #1}
\newcommand{\legende}[1]{\par\vspace{6pt}{\popsemi\fontsize{9}{11.5}\selectfont\color{popmuted}#1}}

% ============================================================
% Signature OnBuch+
% ============================================================
\newcommand{\poplogoinline}[1]{{\poplogo\fontsize{#1}{#1}\selectfont\color{popink}OnBuch\color{poporange}+}}
\newcommand{\signatureonbuch}{%
  \begin{tcolorbox}[enhanced,colback=popink,colframe=popink,boxrule=2.2pt,arc=10pt,
      drop shadow=poporange,left=14pt,right=14pt,top=10pt,bottom=10pt,before skip=0pt,after skip=0pt]
    \tikz[baseline=-0.7ex]\node[circle,fill=popgreen,draw=popcream,line width=1.3pt,minimum size=22pt,inner sep=0]{\popcheck{white}};%
    \hspace{9pt}\begin{minipage}[c]{0.62\linewidth}
      {\popxbold\fontsize{12}{13}\selectfont\color{popcream}Signé\ \textbullet\ {\poplogo OnBuch\color{poporange}+}}\par\vspace{2pt}
      {\raggedright\popsemi\fontsize{8}{10.5}\selectfont\color{popline}\MakeUppercase{Équipe pédagogique OnBuch+}\\\MakeUppercase{Conforme au programme officiel MINESEC}\par}
    \end{minipage}\hfill
    {\poplogo\fontsize{22}{22}\selectfont\color{popcream}OnBuch\color{poporange}+}
  \end{tcolorbox}%
}

% ============================================================
% Page de garde
% #1 titre · #2 matière · #3 niveau · #4 module · #5 n° leçon · #6 durée ·
% #7 description / objectifs (texte ou itemize)
% ============================================================
\newcommand{\pagegarde}[7]{%
  \thispagestyle{empty}
  \newgeometry{a4paper,margin=1.7cm,top=1.6cm,bottom=1.4cm}%
  \begin{tikzpicture}[remember picture,overlay]
    \fill[poporange!18] ([xshift=-3.2cm,yshift=-3.6cm]current page.north east) circle (4.6cm);
    \fill[poppurple!14] ([xshift=2.4cm,yshift=7.6cm]current page.south west) circle (3.8cm);
  \end{tikzpicture}%
  {\poplogo\fontsize{30}{30}\selectfont\color{popink}OnBuch\color{poporange}+}\hfill
  \raisebox{0.6ex}{\poppill{Cours signé \textbullet\ Premium}{popink}{popcream}}\par
  \vspace{1pt}\popsquiggle{poppurple}\par
  \vspace{8pt}
  \begin{tcolorbox}[enhanced,colback=poporange,colframe=popink,boxrule=2.8pt,arc=13pt,
      drop shadow=popink,left=18pt,right=18pt,top=12pt,bottom=13pt,after skip=10pt]
    \poppill{#2 \textbullet\ #3}{popink}{popcream}\hspace{5pt}\poppill{#4}{white}{popink}\par\vspace{8pt}
    {\popdisplay\fontsize{14}{14}\selectfont\color{popink}LEÇON #5}\par\vspace{4pt}
    {\raggedright\hyphenpenalty=10000\exhyphenpenalty=10000\popdisplay\fontsize{25}{27}\selectfont\color{popcream}\MakeUppercase{#1}\par}
    \vspace{8pt}
    {\popxbold\fontsize{9.5}{9.5}\selectfont\color{popink}\MakeUppercase{Durée conseillée : #6}}
  \end{tcolorbox}
  \begin{tcolorbox}[enhanced,colback=white,colframe=popink,boxrule=2.2pt,arc=10pt,
      drop shadow=popink,left=16pt,right=16pt,top=10pt,bottom=10pt,after skip=8pt]
    \poppill{Dans cette leçon}{poppurple}{white}\par\vspace{5pt}
    \popsemi\fontsize{9.3}{12.6}\selectfont\color{popink}#7
  \end{tcolorbox}
  \noindent\begin{minipage}[t]{0.487\linewidth}
    \begin{tcolorbox}[enhanced,colback=popgreen,colframe=popink,boxrule=2.2pt,arc=10pt,
        drop shadow=popink,left=11pt,right=11pt,top=10pt,bottom=10pt,height=3.1cm,valign=center]
      \tcbox[colback=white,colframe=popink,boxrule=1.3pt,arc=5pt,boxsep=0pt,nobeforeafter,
        left=3pt,right=3pt,top=3pt,bottom=3pt,valign=center]{\qrcode[height=1.9cm]{https://play.google.com/store/apps/details?id=com.luvvix.onbuch&hl=fr}}%
      \hspace{8pt}\begin{minipage}[c]{0.5\linewidth}
        {\popdisplay\fontsize{11}{12.5}\selectfont\color{white}\MakeUppercase{Télécharge OnBuch}}\par\vspace{4pt}
        {\raggedright\popsemi\fontsize{8.4}{11}\selectfont\color{white}Gratuit \textbullet\ Google Play\\Tuteur IA, annales, concours, résultats\par}
      \end{minipage}
    \end{tcolorbox}
  \end{minipage}\hfill
  \begin{minipage}[t]{0.487\linewidth}
    \begin{tcolorbox}[enhanced,colback=poppurple,colframe=popink,boxrule=2.2pt,arc=10pt,
        drop shadow=popink,left=11pt,right=11pt,top=10pt,bottom=10pt,height=3.1cm,valign=center]
      \tikz[baseline=-0.7ex]\node[circle,fill=white,draw=popink,line width=1.3pt,minimum size=1.6cm,inner sep=0]{\popdisplay\fontsize{24}{24}\selectfont\color{poppurple}+};%
      \hspace{8pt}\begin{minipage}[c]{0.6\linewidth}
        {\popdisplay\fontsize{11}{12.5}\selectfont\color{white}\MakeUppercase{Passe à OnBuch+}}\par\vspace{4pt}
        {\raggedright\popsemi\fontsize{8.4}{11}\selectfont\color{white}Tous les cours signés, Tuteur IA illimité, fiches PDF — onglet OnBuch+ de l'app\par}
      \end{minipage}
    \end{tcolorbox}
  \end{minipage}
  \vspace{8pt}
  \popcontenu
  \vfill
  \signatureonbuch
  \restoregeometry
  \newpage
}

% Rangée « Ce cours contient » (page de garde)
\newcommand{\poptile}[3]{% #1 couleur #2 gros texte #3 légende
  \begin{tcolorbox}[enhanced,colback=#1,colframe=popink,boxrule=2pt,arc=9pt,shadow={2.5pt}{-2.5pt}{0pt}{popink},
      left=6pt,right=6pt,top=9pt,bottom=9pt,halign=center,valign=center,height=1.75cm,width=\linewidth,nobeforeafter]
    {\popdisplay\fontsize{12.5}{13}\selectfont\color{white}\MakeUppercase{#2}}\par\vspace{3pt}
    {\centering\popsemi\fontsize{7.8}{9.5}\selectfont\color{white}#3\par}
  \end{tcolorbox}}
\newcommand{\popcontenu}{%
  \par\noindent{\popxboldls\fontsize{7.5}{7.5}\selectfont\color{popmuted}CE COURS SIGNÉ CONTIENT}\par\vspace{7pt}
  \noindent\begin{minipage}[t]{0.235\linewidth}\poptile{popblue}{Cours}{complet, illustré, pas à pas}\end{minipage}\hfill
  \begin{minipage}[t]{0.235\linewidth}\poptile{poporange}{Méthodes}{et exercices résolus}\end{minipage}\hfill
  \begin{minipage}[t]{0.235\linewidth}\poptile{poppink}{Exercices}{type examen + corrigés}\end{minipage}\hfill
  \begin{minipage}[t]{0.235\linewidth}\poptile{popgreen}{Fiche bilan}{l'essentiel à retenir}\end{minipage}\par}

% Sommaire
\newtcolorbox{sommaire}{enhanced,colback=white,colframe=popink,boxrule=2.2pt,arc=10pt,
  drop shadow=popink,left=18pt,right=18pt,top=13pt,bottom=13pt,before skip=6pt,after skip=20pt,
  before upper={\poppill{Sommaire}{poppurple}{white}\par\vspace{9pt}\popsemi\fontsize{10.6}{14.5}\selectfont\color{popink}}}

% Signature de fin de cours — poussée tout en bas de la dernière page
\newcommand{\findecours}{%
  \par\vfill\nopagebreak
  \begin{center}\popsquiggle{poporange}\end{center}\vspace{4pt}
  \signatureonbuch
}
\AtBeginDocument{\def\llbracket{\mathopen{[\mkern-3.5mu[}}\def\rrbracket{\mathclose{]\mkern-3.5mu]}}} % intervalles d'entiers (glyphe ⟦ absent de la police)
% Glyphes absents de Fira Math : redéfinis à partir de glyphes disponibles
\AtBeginDocument{%
  \def\setminus{\mathbin{\backslash}}\let\smallsetminus\setminus
  \def\vdots{\mathord{\vbox{\baselineskip3.2pt\lineskiplimit0pt\kern4pt\hbox{.}\hbox{.}\hbox{.}}}}%
  \def\ddots{\mathinner{\mkern1mu\raise6.5pt\hbox{.}\mkern2mu\raise3.6pt\hbox{.}\mkern2mu\raise0.7pt\hbox{.}\mkern1mu}}%
  \def\triangle{\mathord{\scalebox{0.9}{$\symup{\Delta}$}}}%
  \def\nabla{\mathord{\raisebox{\depth}{\scalebox{1}[-1]{$\symup{\Delta}$}}}}%
  \def\checkmark{\popcheck{popdark}}%
  \def\star{\mathbin{\ast}}\def\bigstar{\mathord{\ast}}%
  \def\diamond{\mathbin{\scalebox{0.6}{$\Diamond$}}}%
  \def\bigcup{\mathop{\mathchoice{\raisebox{-0.3em}{\scalebox{1.7}{$\cup$}}}{\scalebox{1.25}{$\cup$}}{\cup}{\cup}}\displaylimits}%
  \def\bigcap{\mathop{\mathchoice{\raisebox{-0.3em}{\scalebox{1.7}{$\cap$}}}{\scalebox{1.25}{$\cap$}}{\cap}{\cap}}\displaylimits}%
  \def\lesseqqgtr{\mathrel{\lessgtr}}\def\bigoplus{\mathop{\oplus}\displaylimits}%
}
\providecommand{\ssi}{si et seulement si} % abréviation fréquente
\AtBeginDocument{\providecommand{\wideparen}{\overparen}} % arc de cercle

\begin{document}

\pagegarde{Lesson 28 — Speaking:Exchanges information about preventing communicable diseases in schools, workplaces, and the community}{Anglais}{Tle A}{Chapter 3 : Environment, Well-being and Health}{EN28}{2 h}{Cette leçon d'expression orale apprend aux élèves de Terminale A à échanger des informations en anglais sur la prévention des maladies transmissibles (malaria, cholera, typhoid, COVID-19) dans les écoles, les lieux de travail et la communauté. À travers des dialogues modèles, des formules fonctionnelles et des jeux de rôle guidés, les élèves développent leur aisance à poser des questions, donner des informations, relancer et conclure une conversation.
\vspace{4pt}
\begin{multicols}{2}\small
\begin{itemize}
\item À la fin de cette leçon, je sais utiliser le vocabulaire essentiel des maladies transmissibles et de leur prévention dans un échange oral.
\item À la fin de cette leçon, je sais demander une information de façon polie et variée (Could you tell me...? / I'd like to know...).
\item À la fin de cette leçon, je sais donner une information claire et la structurer (First... Then... Finally...).
\item À la fin de cette leçon, je sais relancer une conversation qui ralentit (What about...? / Could you explain...?).
\item À la fin de cette leçon, je sais réagir et montrer mon intérêt (Really? / I see. / That's interesting!).
\item À la fin de cette leçon, je sais conclure poliment un échange (Thank you for your time. / Let's keep in touch.).
\end{itemize}\end{multicols}}

\begin{sommaire}
\begin{multicols}{2}
\begin{enumerate}
\item Warm-up : le contexte et le lexique de la santé
\item Dialogue modèle 1 : à l'infirmerie de l'école
\item Les formules fonctionnelles de l'échange d'informations
\item Dialogue modèle 2 : dans la communauté et au travail
\item Prononciation, accentuation et intonation
\item Jeux de rôle guidés et production orale
\item Activité d'intégration
\item Méthodes et exercices résolus
\item Exercices
\item Corrigés des exercices
\item Fiche bilan
\end{enumerate}
\end{multicols}
\end{sommaire}

\begin{objectifs}
\begin{itemize}
\item À la fin de cette leçon, je sais utiliser le vocabulaire essentiel des maladies transmissibles et de leur prévention dans un échange oral.
\item À la fin de cette leçon, je sais demander une information de façon polie et variée (Could you tell me...? / I'd like to know...).
\item À la fin de cette leçon, je sais donner une information claire et la structurer (First... Then... Finally...).
\item À la fin de cette leçon, je sais relancer une conversation qui ralentit (What about...? / Could you explain...?).
\item À la fin de cette leçon, je sais réagir et montrer mon intérêt (Really? / I see. / That's interesting!).
\item À la fin de cette leçon, je sais conclure poliment un échange (Thank you for your time. / Let's keep in touch.).
\item À la fin de cette leçon, je sais prononcer correctement le lexique de la santé avec la bonne accentuation et la bonne intonation.
\item À la fin de cette leçon, je sais mener un court dialogue de 8 à 10 répliques sur la prévention des maladies avec un partenaire.
\end{itemize}
\end{objectifs}

\begin{prerequis}
\begin{itemize}
\item Les questions en WH- et en Yes/No (Première) : What, How, Why, Do you...?
\item Le présent simple et les modaux should / must / can pour les conseils et obligations
\item Le vocabulaire de base de la santé vu en Première : disease, sick, hospital, medicine
\item Les salutations et formules de politesse (Hello, Excuse me, Thank you)
\item La leçon de vocabulary du chapitre 3 sur l'environnement et la santé
\end{itemize}
\end{prerequis}

\newpage

\coursec{Warm-up : le contexte et le lexique de la santé}

Imagine : tu es délégué de classe au Lycée Bilingue de Yaoundé. Après plusieurs cas de paludisme et de choléra signalés dans ton quartier, le proviseur te demande d'organiser une campagne de sensibilisation. Tu dois interroger l'infirmière scolaire, échanger avec tes camarades et convaincre les élèves d'adopter les bons gestes. Comment demander une information en anglais ? Comment relancer la conversation quand ton interlocuteur hésite ? Cette leçon te donne toutes les clés pour échanger des informations sur la prévention des maladies transmissibles, à l'école, au travail et dans la communauté.

\textbf{Problématique :} How can we exchange clear and accurate information about preventing communicable diseases in English ?

\textbf{Objectifs de la leçon :} à la fin de cette leçon, tu seras capable de (1) nommer les principales maladies transmissibles et leurs symptômes en anglais, (2) expliquer les modes de transmission et les moyens de prévention, (3) poser des questions et répondre de façon adaptée dans un court dialogue sur la santé.

\courssub{Brainstorming : les maladies transmissibles au Cameroun}

Avant de parler, il faut des mots. Ferme les yeux et réponds à cette question : \emph{Which communicable diseases do you know in Cameroon ?} (Quelles maladies transmissibles connais-tu au Cameroun ?) Voici les réponses qu'un bon élève de Terminale doit pouvoir donner spontanément :

\begin{center}
\begin{tabularx}{\linewidth}{|X|X|X|}
\hline
\rowcolor{popblueL} English & Français & Remarque \\
\hline
malaria & le paludisme & transmis par les moustiques \\
\hline
cholera & le choléra & prononce « ko-lé-ra » \\
\hline
typhoid fever & la fièvre typhoïde & eau et aliments contaminés \\
\hline
tuberculosis (TB) & la tuberculose & prononce « tiou-beur-kiou-lô-siss » \\
\hline
COVID-19 & la COVID-19 & maladie virale respiratoire \\
\hline
meningitis & la méningite & prononce « mé-nin-djaï-tiss » \\
\hline
measles & la rougeole & maladie évitable par la vaccination \\
\hline
\end{tabularx}
\end{center}

\begin{attention}[Prononciation et orthographe]
En anglais, les noms de maladies ne prennent \cle{pas d'article} quand on parle de la maladie en général : \eng{Malaria is common in Cameroon.} (Le paludisme est fréquent au Cameroun.) Ne dis pas \eng{the malaria} dans ce cas. Attention aussi à \eng{cholera} qui se prononce « ko-lé-ra » et non « tcho-lé-ra ». Enfin, \eng{measles} (la rougeole) se termine par un \textbf{-s} mais se construit au \cle{singulier} : \eng{Measles is a dangerous disease.} (La rougeole est une maladie dangereuse.)
\end{attention}

\courssub{Observer avant d'apprendre : un texte d'accroche}

Lis ce court article de presse scolaire. Souligne mentalement tous les mots liés à la santé : tu les retrouveras dans les tableaux de vocabulaire qui suivent.

\begin{textebox}[The Bilingual Herald — Yaoundé, school newspaper]
\textbf{School Launches Health Campaign After Cholera Alert}

Last week, the headmaster of Bilingual High School, Yaoundé, announced a major health campaign after two cases of cholera were reported in the neighbourhood. Cholera is a communicable disease that spreads through contaminated water and food. Its main symptoms are severe diarrhoea, vomiting and fever.

The school nurse, Mrs Etoundi, explained to students that prevention is simple but essential. “Wash your hands with soap and clean water before every meal,” she said. “Never drink water from unknown sources, and always wash fruits and vegetables.”

The campaign also focuses on malaria, which remains the most common disease in the school infirmary. Students are encouraged to sleep under mosquito nets and to clear stagnant water around their homes, because mosquitoes breed in standing water.

Community health workers will visit the school next Monday to talk about vaccination and sanitation. The headmaster declared: “A healthy student is a successful student. Prevention is better than cure.” Students will prepare posters in English and in French to inform the whole community about hygiene and disease prevention.
\end{textebox}

\textbf{Glossaire :}

\begin{center}
\begin{tabularx}{\linewidth}{|X|X|X|}
\hline
\rowcolor{popblueL} Mot & Nature & Traduction \\
\hline
to launch & verbe & lancer \\
\hline
neighbourhood & nom & quartier \\
\hline
severe & adjectif & grave, sévère \\
\hline
stagnant / standing water & nom & eau stagnante \\
\hline
to breed & verbe & se reproduire \\
\hline
to clear & verbe & dégager, éliminer \\
\hline
poster & nom & affiche \\
\hline
\end{tabularx}
\end{center}

\textbf{Comprehension questions :}
\begin{enumerate}
\item What disease was reported in the neighbourhood ?
\item What are the main symptoms of cholera mentioned in the text ?
\item Give two pieces of advice from the school nurse.
\item Why should students clear stagnant water ?
\item Who will visit the school next Monday, and why ?
\end{enumerate}

\begin{corrige}[Comprehension questions]
\begin{enumerate}
\item Two cases of \eng{cholera} were reported in the neighbourhood. (Deux cas de choléra ont été signalés dans le quartier.)
\item The main symptoms of cholera are \eng{severe diarrhoea, vomiting and fever}. (Diarrhée sévère, vomissements et fièvre.)
\item The nurse advises students to \eng{wash their hands with soap and clean water before every meal}, \eng{never to drink water from unknown sources} and \eng{always to wash fruits and vegetables}. (Deux de ces conseils suffisent.)
\item Because \eng{mosquitoes breed in standing water} ; clearing it helps to prevent malaria. (Parce que les moustiques se reproduisent dans l'eau stagnante.)
\item \eng{Community health workers} will visit the school to talk about \eng{vaccination and sanitation}. (Des agents de santé communautaire viendront parler de vaccination et d'assainissement.)
\end{enumerate}
\end{corrige}

\begin{textebox}[Proverbe anglais]
“Prevention is better than cure.” \\
« Mieux vaut prévenir que guérir. »
\end{textebox}

\courssub{Le lexique des maladies}

\begin{definition}[Communicable disease]
A \cle{communicable disease} (or \cle{infectious disease}) is a disease that can be passed from one person to another. (Une maladie transmissible d'une personne à une autre.) On dit aussi \eng{contagious disease} pour insister sur la contagion par contact direct.
\end{definition}

Voici le vocabulaire essentiel pour décrire une maladie et ses signes :

\begin{center}
\begin{tabularx}{\linewidth}{|X|X|X|}
\hline
\rowcolor{popblueL} English & Français & Exemple \\
\hline
an outbreak & un foyer épidémique, une flambée & \eng{There is a cholera outbreak in Douala.} \\
\hline
an epidemic & une épidémie & \eng{The epidemic spread quickly.} \\
\hline
a symptom & un symptôme & \eng{Fever is a common symptom.} \\
\hline
fever & la fièvre & \eng{The patient has a high fever.} \\
\hline
a cough / to cough & une toux / tousser & \eng{He coughs all night.} \\
\hline
diarrhoea & la diarrhée & \eng{Diarrhoea is a symptom of cholera.} \\
\hline
a headache & un mal de tête & \eng{She has a terrible headache.} \\
\hline
vomiting / to vomit & les vomissements / vomir & \eng{The child is vomiting.} \\
\hline
\end{tabularx}
\end{center}

\begin{propriete}[Orthographe et prononciation à risque]
\begin{itemize}
\item \eng{diarrhoea} : orthographe britannique avec \textbf{-oea} (américain : \eng{diarrhea}) ; prononce « daïa-ri-a ».
\item \eng{cough} : prononce « kof » ; le \textbf{-gh} se prononce « f ».
\item \eng{headache} : prononce « hé-dèik » ; le \textbf{-ch-} se prononce « k ».
\item \eng{symptom} : prononce « sin-ptom » ; le \textbf{p} se prononce.
\item \eng{fever} : prononce « fi-veur » ; on dit \eng{to have a fever} (avoir de la fièvre), jamais \eng{to have fever} avec \eng{a} oublié quand on précise : \eng{He has a high fever.}
\end{itemize}
\end{propriete}

\courssub{Le lexique de la transmission}

\begin{center}
\begin{tabularx}{\linewidth}{|X|X|X|}
\hline
\rowcolor{popblueL} English & Français & Exemple \\
\hline
to spread (spread, spread) & se propager, propager & \eng{Diseases spread fast in crowded places.} \\
\hline
to transmit & transmettre & \eng{Mosquitoes transmit malaria.} \\
\hline
to catch a disease & attraper une maladie & \eng{He caught malaria last year.} \\
\hline
contaminated water & eau contaminée & \eng{Never drink contaminated water.} \\
\hline
germs & les microbes & \eng{Soap kills germs.} \\
\hline
mosquitoes & les moustiques & \eng{Mosquitoes breed in stagnant water.} \\
\hline
\end{tabularx}
\end{center}

\begin{attention}[Faux amis et erreurs fréquentes]
Ne dis jamais \eng{to take a disease} : dis \cle{to catch a disease} (attraper une maladie) ; prétérit et participe passé irréguliers : \eng{catch — caught — caught}. De même, ne confonds pas :
\begin{itemize}
\item \eng{sick} : malade (courant) — \eng{My brother is sick today.}
\item \eng{ill} : malade (plus soutenu) — \eng{She fell ill during the holidays.}
\item \eng{injured} : blessé (accident, blessure physique) — \eng{Two players were injured during the match.}
\end{itemize}
Un élève qui a le paludisme est \eng{sick} ou \eng{ill}, jamais \eng{injured} ! Attention aussi au verbe \eng{to spread} qui est \cle{invariable} : \eng{spread — spread — spread}. Ne conjugue jamais \eng{spreaded}.
\end{attention}

\courssub{Le lexique de la prévention}

\begin{center}
\begin{tabularx}{\linewidth}{|X|X|X|}
\hline
\rowcolor{popgreenL} English & Français & Exemple \\
\hline
to prevent & prévenir, empêcher & \eng{We can prevent cholera.} \\
\hline
prevention & la prévention & \eng{Prevention is better than cure.} \\
\hline
hygiene & l'hygiène & \eng{Good hygiene saves lives.} \\
\hline
to wash one's hands & se laver les mains & \eng{Wash your hands with soap.} \\
\hline
clean water & eau propre, potable & \eng{Drink only clean water.} \\
\hline
a mosquito net & une moustiquaire & \eng{Sleep under a mosquito net.} \\
\hline
vaccination & la vaccination & \eng{Vaccination protects children.} \\
\hline
to vaccinate & vacciner & \eng{All babies are vaccinated.} \\
\hline
sanitation & l'assainissement & \eng{Poor sanitation spreads diseases.} \\
\hline
\end{tabularx}
\end{center}

\begin{exemplebox}[Exemples traduits à mémoriser]
\eng{Cholera spreads through contaminated water.} = Le choléra se propage par l'eau contaminée.

\eng{We sleep under mosquito nets to prevent malaria.} = Nous dormons sous des moustiquaires pour prévenir le paludisme.

\eng{Washing your hands with soap kills germs.} = Se laver les mains avec du savon tue les microbes.

\eng{Community health workers vaccinate children against measles.} = Les agents de santé communautaire vaccinent les enfants contre la rougeole.
\end{exemplebox}

\begin{propriete}[La famille du mot PREVENT]
\cle{to prevent} (verbe) $\rightarrow$ \cle{prevention} (nom) $\rightarrow$ \cle{preventive} (adjectif). Attention à la structure : \eng{to prevent something} ou \eng{to prevent somebody from doing something} : \eng{Nets prevent mosquitoes from biting us.} (Les moustiquaires empêchent les moustiques de nous piquer.) Ne dis pas \eng{prevent to do}. Même logique pour la vaccination : on vaccine \cle{against} une maladie : \eng{to vaccinate children against measles} (vacciner les enfants contre la rougeole).
\end{propriete}

\courssub{Les lieux et les acteurs de la santé}

\begin{center}
\begin{tabularx}{\linewidth}{|X|X|X|}
\hline
\rowcolor{popblueL} English & Français & Remarque \\
\hline
the school nurse & l'infirmière scolaire & prononce « skeul neurs » \\
\hline
the health centre & le centre de santé & CSI au Cameroun \\
\hline
a community health worker & un agent de santé communautaire & relais communautaire \\
\hline
the workplace & le lieu de travail & bureau, usine, marché \\
\hline
the headmaster & le proviseur & prononce « hèd-mas-teur » \\
\hline
the local authorities & les autorités locales & maire, préfet, chefs \\
\hline
\end{tabularx}
\end{center}

\courssub{Carte mentale : organiser le lexique}

\begin{popfigure}
\begin{center}
\begin{tikzpicture}[mindmap, grow cyclic, every node/.style={concept}, concept color=popblue!60, text=white,
level 1/.append style={level distance=3.4cm, sibling angle=90},
level 2/.append style={level distance=2.2cm, sibling angle=45}]
\node {COMMUNICABLE DISEASES}
  child[concept color=poporange!80] { node {Diseases}
    child { node[align=center] {malaria\\cholera\\typhoid} }
    child { node[align=center] {tuberculosis\\meningitis} }
  }
  child[concept color=poppurple!80] { node {Transmission}
    child { node[align=center] {to spread\\to transmit\\to catch} }
    child { node[align=center] {germs\\mosquitoes\\dirty water} }
  }
  child[concept color=popgreen!80] { node {Prevention}
    child { node[align=center] {wash hands\\clean water} }
    child { node[align=center] {mosquito nets\\vaccination} }
  }
  child[concept color=poppink!80] { node[align=center] {People \&\\Places}
    child { node[align=center] {school nurse\\health centre} }
    child { node[align=center] {health worker\\headmaster} }
  };
\end{tikzpicture}
\end{center}
\legende{Carte mentale du lexique : quatre familles de mots pour parler des maladies transmissibles.}
\end{popfigure}

\begin{savaistu}[Le paludisme au Cameroun]
Au Cameroun, le paludisme (\eng{malaria}) reste la première cause de consultation médicale dans les centres de santé. Le mot \eng{malaria} vient de l'italien \emph{mala aria}, qui signifie « mauvais air » : autrefois, on croyait que l'air vicié des marais rendait malade. On sait aujourd'hui que c'est la piqûre du moustique femelle \emph{anophèle} qui transmet la maladie — mais le mot est resté !
\end{savaistu}

\courssub{Exercice résolu : activer le lexique}

\begin{exoresolu}[Complète avec le mot qui convient]
Complète les phrases avec : \eng{outbreak — symptoms — mosquito net — contaminated — vaccinate — spread}.

1. Cholera can \ldots{} very quickly in crowded areas. \quad 2. Fever and diarrhoea are \ldots{} of cholera. \quad 3. There is a cholera \ldots{} in the neighbourhood. \quad 4. Never drink \ldots{} water. \quad 5. Sleep under a \ldots{} to avoid mosquito bites. \quad 6. Nurses \ldots{} children against many diseases.

\tcblower
\textbf{Solution détaillée :}

1. \eng{spread} — le sujet est une maladie, on parle de propagation : « Le choléra peut se propager très vite. » Après le modal \eng{can}, le verbe reste à la base verbale.

2. \eng{symptoms} — la fièvre et la diarrhée sont des signes de la maladie : « des symptômes du choléra ». Attention au pluriel : deux signes = \eng{symptoms}.

3. \eng{outbreak} — une flambée localisée dans un quartier : « un foyer de choléra ».

4. \eng{contaminated} — adjectif placé devant le nom \eng{water} : « eau contaminée ».

5. \eng{mosquito net} — « dors sous une moustiquaire pour éviter les piqûres de moustiques ».

6. \eng{vaccinate} — verbe conjugué au présent simple, sujet pluriel : « Les infirmières vaccinent les enfants. »
\end{exoresolu}

\courssub{Mise en situation orale rapide : l'interview}

Pour finir ce warm-up, place à l'oral. Le professeur joue le rôle d'un journaliste d'une radio communautaire (\eng{a community radio journalist}) et interviewe un élève sur le paludisme. Écoute le modèle, puis rejoue-le avec un camarade.

\begin{textebox}[Model interview — “Health Talk”, community radio]
\textbf{Journalist :} Good morning, dear listeners ! Today we are at Bilingual High School in Yaoundé. Excuse me, what is your name ?

\textbf{Student :} My name is Amina. I am a student in Terminale A.

\textbf{Journalist :} Amina, can you tell me what malaria is ?

\textbf{Student :} Yes, of course. Malaria is a communicable disease transmitted by mosquitoes.

\textbf{Journalist :} What are its main symptoms ?

\textbf{Student :} The main symptoms are high fever, headaches and vomiting.

\textbf{Journalist :} I see. And how can we prevent malaria in our community ?

\textbf{Student :} We should sleep under mosquito nets, clear stagnant water and keep our environment clean.

\textbf{Journalist :} Thank you very much, Amina. Remember, dear listeners : prevention is better than cure !
\end{textebox}

Observe la structure du dialogue : le journaliste \cle{ouvre} la conversation (\eng{Good morning, excuse me}), \cle{pose des questions} ouvertes (\eng{Can you tell me... ? What are... ? How can we... ?}), \cle{réagit} aux réponses (\eng{I see.}) et \cle{conclut} (\eng{Thank you very much}). C'est exactement cette architecture que tu devras reproduire dans tes propres dialogues.

\begin{experience}[À toi de jouer !]
Par deux : l'un est journaliste, l'autre est élève expert. Remplace le paludisme par le choléra ou la fièvre typhoïde, et utilise au moins cinq mots du lexique du jour (\eng{outbreak, symptoms, contaminated, prevention, hygiene, sanitation...}). Pose au moins trois questions : \eng{What is... ? What are the symptoms of... ? How can we prevent... ?} Puis échangez les rôles.
\end{experience}

\begin{aretenir}[L'essentiel du warm-up]
\begin{itemize}
\item Les noms de maladies en général : \cle{sans article} — \eng{Cholera is dangerous.}
\item Attraper une maladie : \cle{to catch a disease} (\eng{catch — caught — caught}).
\item Se propager : \cle{to spread} (\eng{spread — spread — spread}), invariable.
\item Prévenir : \cle{to prevent something} / \cle{to prevent somebody from doing something}.
\item Vacciner : \cle{to vaccinate somebody against a disease}.
\item Malade : \eng{sick} ou \eng{ill} ; \eng{injured} = blessé (accident).
\end{itemize}
\end{aretenir}

Tu disposes maintenant du contexte et du vocabulaire de base. Dans la prochaine étape, nous étudierons un premier dialogue modèle à l'infirmerie de l'école pour apprendre à demander et donner une information avec naturel.

\coursec{Dialogue modèle 1 : à l'infirmerie de l'école}

Dans la section précédente, nous avons découvert le contexte de la santé et le vocabulaire des maladies transmissibles : \eng{malaria}, \eng{cholera}, \eng{typhoid fever}, \eng{mosquito net}, \eng{stagnant water}. Ce lexique va maintenant prendre vie dans une situation de communication réelle. Notre objectif dans cette section : écouter, lire, comprendre et rejouer un premier dialogue modèle entre un élève et l'infirmière de son lycée. C'est le modèle sur lequel vous vous appuierez pour mener vos propres échanges d'informations, à l'école comme dans la communauté.

\courssub{Écoute et lecture du dialogue modèle}

Fermez d'abord votre manuel : l'enseignant lit le dialogue une première fois, et vous écoutez pour saisir l'impression générale (\emph{qui parle ? où ? de quoi ?}). Puis ouvrez-le et suivez la lecture à voix basse, en soulignant mentalement les mots que vous reconnaissez.

\begin{textebox}[At the school sick bay — À l'infirmerie de l'école]
\textbf{Abong:} Good morning, Mrs Etame. Could I ask you a few questions about malaria?

\textbf{Nurse:} Of course, Abong. Go ahead.

\textbf{Abong:} How is malaria transmitted?

\textbf{Nurse:} It is spread by mosquito bites, especially at night.

\textbf{Abong:} I see. And what can we do to prevent it at school?

\textbf{Nurse:} First, sleep under treated mosquito nets. Then, clear stagnant water around the school. Finally, go to the health centre if you have a fever.

\textbf{Abong:} That's very useful. Thank you for your time, Mrs Etame.

\textbf{Nurse:} You're welcome. Stay healthy!
\end{textebox}

\textbf{Glossaire} (mots et expressions à retenir) :

\begin{center}
\begin{tabularx}{\linewidth}{|l|l|X|}\hline
\rowcolor{popblueL} \textbf{English} & \textbf{Nature} & \textbf{Français} \\ \hline
sick bay & nom & infirmerie (d'une école) \\ \hline
nurse & nom & infirmier / infirmière \\ \hline
a few questions & expression & quelques questions \\ \hline
to transmit & verbe & transmettre \\ \hline
to spread (spread, spread) & verbe irrégulier & propager, transmettre \\ \hline
mosquito bite & nom & piqûre de moustique \\ \hline
to prevent & verbe & prévenir, empêcher \\ \hline
treated mosquito net & nom & moustiquaire imprégnée \\ \hline
stagnant water & nom & eau stagnante \\ \hline
health centre & nom & centre de santé \\ \hline
fever & nom & fièvre \\ \hline
\end{tabularx}
\end{center}

\begin{attention}[Prononciation des mots clés]
Pas d'alphabet phonétique : notez la prononciation en lettres françaises.
\begin{itemize}
\item \eng{nurse} se prononce « neu-seu » (un seul son long, ne dites pas « nour-ceu »).
\item \eng{fever} se prononce « fi-veu(r) ».
\item \eng{transmitted} : accent sur la deuxième syllabe — « trans-MI-tid ».
\item \eng{mosquito} se prononce « mo-SKI-to » (accent sur la deuxième syllabe).
\item \eng{stagnant} se prononce « sta-gnant » : le \eng{g} se prononce, contrairement au français \emph{stagnant}.
\item \eng{health} se prononce « helth » : le \eng{th} se dit avec la langue entre les dents, jamais « z » ni « s ».
\end{itemize}
\end{attention}

\courssub{Compréhension globale : qui parle ? où ? de quoi ?}

Avant de chercher le détail, on répond toujours aux trois questions fondamentales. C'est la démarche exigée à l'épreuve de compréhension du Baccalauréat : d'abord le sens général, ensuite les détails.

\begin{methode}[Comprendre un dialogue en quatre étapes]
\begin{enumerate}
\item \textbf{Identifier les locuteurs et le lieu} : qui parle, et où se déroule la scène ? Ici : \eng{Abong}, un élève, et \eng{Mrs Etame}, l'infirmière scolaire, à l'infirmerie (\eng{the school sick bay}).
\item \textbf{Repérer le sujet} : de quoi parlent-ils ? Ici : \eng{malaria} (le paludisme) et sa prévention à l'école.
\item \textbf{Noter les informations clés} : maladie, mode de transmission, mesures de prévention. Soulignez-les dans le texte.
\item \textbf{Relever les formules fonctionnelles} : comment salue-t-on ? comment demande-t-on une information ? comment remercie-t-on ? Ce sont ces formules que vous réutiliserez dans vos propres dialogues.
\end{enumerate}
\end{methode}

\begin{exoresolu}[Compréhension globale]
Répondez en anglais par des phrases complètes : 1) Who are the speakers? 2) Where does the dialogue take place? 3) What is the topic of the conversation?
\tcblower
1) \eng{The speakers are Abong, a student, and Mrs Etame, the school nurse.} (Les interlocuteurs sont Abong, un élève, et Mrs Etame, l'infirmière scolaire.)

2) \eng{The dialogue takes place at the school sick bay.} (Le dialogue se déroule à l'infirmerie de l'école.)

3) \eng{They are talking about malaria: how it is transmitted and how to prevent it at school.} (Ils parlent du paludisme : comment il se transmet et comment le prévenir à l'école.)
\end{exoresolu}

\courssub{Compréhension fine : maladies, transmission et prévention}

Lisons maintenant le dialogue en détail. Trois types d'informations doivent être repérés dans tout échange sur la santé : la \cle{maladie}, son mode de \cle{transmission} et les mesures de \cle{prévention}.

\begin{itemize}
\item \textbf{Maladie citée} : \eng{malaria} (le paludisme).
\item \textbf{Transmission} : \eng{It is spread by mosquito bites, especially at night.} « Elle est transmise par les piqûres de moustiques, surtout la nuit. »
\item \textbf{Prévention} : trois conseils, introduits par les marqueurs d'étapes \eng{First}, \eng{Then}, \eng{Finally} :
\begin{enumerate}
\item \eng{Sleep under treated mosquito nets.} « Dormez sous des moustiquaires imprégnées. »
\item \eng{Clear stagnant water around the school.} « Éliminez les eaux stagnantes autour de l'école. »
\item \eng{Go to the health centre if you have a fever.} « Allez au centre de santé si vous avez de la fièvre. »
\end{enumerate}
\end{itemize}

\begin{propriete}[Les marqueurs d'étapes : First, Then, Finally]
Pour énumérer des conseils ou des étapes, l'anglais utilise des \cle{connecteurs de séquence} : \eng{First} (d'abord), \eng{Then} / \eng{Next} (ensuite), \eng{Finally} (enfin). Ils se placent en tête de phrase, suivis d'une virgule. Ils sont indispensables à l'oral pour structurer une information et guider l'auditeur.

\eng{First, wash your hands. Then, boil the water. Finally, cover the food.} « D'abord, lavez-vous les mains. Ensuite, faites bouillir l'eau. Enfin, couvrez la nourriture. »

Remarque : les conseils sont donnés à l'\cle{impératif} (base verbale, sans sujet) : \eng{Sleep...}, \eng{Clear...}, \eng{Go...} C'est la forme normale des instructions et des recommandations en anglais.
\end{propriete}

\begin{propriete}[Le passif pour décrire la transmission : It is spread by...]
Pour dire comment une maladie se transmet, l'anglais emploie très souvent le \cle{passif} : sujet + \eng{be} (conjugué) + participe passé (+ \eng{by} + agent). On l'utilise parce que l'important est la maladie, pas celui qui transmet.

\eng{Malaria is spread by mosquito bites.} « Le paludisme est transmis par les piqûres de moustiques. »

\eng{Cholera is caused by dirty water.} « Le choléra est causé par l'eau sale. »

Attention au verbe irrégulier \eng{to spread} : ses trois formes sont identiques — \eng{spread, spread, spread}. Ne dites jamais \eng{*spreaded}.
\end{propriete}

\begin{center}
\begin{tabularx}{\linewidth}{|X|X|X|}\hline
\rowcolor{popblueL} \textbf{Français} & \textbf{English} & \textbf{Remarque} \\ \hline
Le paludisme est transmis par les moustiques. & Malaria is spread by mosquitoes. & passif : \eng{is spread by} \\ \hline
Le choléra se transmet par l'eau contaminée. & Cholera is transmitted through contaminated water. & \eng{through} = par le biais de \\ \hline
On peut le prévenir en dormant sous une moustiquaire. & It can be prevented by sleeping under a mosquito net. & \eng{by} + verbe en \eng{-ing} \\ \hline
\end{tabularx}
\end{center}

\courssub{Repérage des formules fonctionnelles}

Un échange d'informations suit une architecture régulière. Observons les formules du dialogue, fonction par fonction.

\begin{exemplebox}[Les formules du dialogue, traduites]
\begin{itemize}
\item \textbf{Saluer} : \eng{Good morning, Mrs Etame.} « Bonjour, Mrs Etame. »
\item \textbf{Demander la permission de poser des questions} : \eng{Could I ask you a few questions about malaria?} « Puis-je vous poser quelques questions sur le paludisme ? »
\item \textbf{Donner la parole} : \eng{Of course, Abong. Go ahead.} « Bien sûr, Abong. Je vous écoute. / Allez-y. »
\item \textbf{Demander une information} : \eng{How is malaria transmitted?} « Comment le paludisme est-il transmis ? »
\item \textbf{Donner une information} : \eng{It is spread by mosquito bites.} « Elle est transmise par les piqûres de moustiques. »
\item \textbf{Réagir / montrer qu'on suit} : \eng{I see.} « Je vois. » ; \eng{That's very useful.} « C'est très utile. »
\item \textbf{Remercier} : \eng{Thank you for your time, Mrs Etame.} « Merci pour votre temps, Mrs Etame. »
\item \textbf{Répondre à un remerciement et conclure} : \eng{You're welcome. Stay healthy!} « Je vous en prie. Restez en bonne santé ! »
\end{itemize}
\end{exemplebox}

\begin{aretenir}[La structure d'un échange d'informations]
Tout dialogue d'information suit cinq étapes : \cle{saluer} (\eng{greeting}) → \cle{demander une information} (\eng{asking for information}) → \cle{donner l'information} (\eng{giving information}) → \cle{réagir} (\eng{reacting}) → \cle{remercier et conclure} (\eng{concluding}). Gardez ce schéma en tête : c'est la colonne vertébrale de tous vos jeux de rôle.
\end{aretenir}

\begin{popfigure}
\begin{tikzpicture}[node distance=0.55cm and 0.9cm,
  etape/.style={rectangle, rounded corners=4pt, draw=popblue, fill=popblueL, thick, align=center, minimum width=3.1cm, minimum height=1cm, font=\small},
  fleche/.style={-{Stealth[length=3mm]}, thick, popdark}]
\node[etape, align=center] (g){\textbf{Greeting}\\ \eng{Good morning!}};
\node[etape, right=of g, fill=poppurpleL, draw=poppurple, align=center] (a){\textbf{Asking}\\ \eng{Could I ask you...?}};
\node[etape, right=of a, fill=popgreenL, draw=popgreen, align=center] (gi){\textbf{Giving information}\\ \eng{It is spread by...}};
\node[etape, below=of gi, fill=poporangeL, draw=poporange, align=center] (r){\textbf{Reacting}\\ \eng{I see. That's useful.}};
\node[etape, left=of r, fill=popgoldL, draw=popgold, align=center] (c){\textbf{Concluding}\\ \eng{Thank you. You're welcome.}};
\draw[fleche] (g) -- (a);
\draw[fleche] (a) -- (gi);
\draw[fleche] (gi) -- (r);
\draw[fleche] (r) -- (c);
\end{tikzpicture}
\legende{Organigramme du dialogue : les cinq étapes de l'échange d'informations.}
\end{popfigure}

\begin{attention}[Faux amis et pièges de traduction]
\begin{itemize}
\item \eng{Go ahead.} ne veut \textbf{pas} dire « va devant » : c'est une invitation à parler — « je t'écoute, vas-y, allez-y ».
\item \eng{You're welcome.} se dit \textbf{après un remerciement} (« je vous en prie »), jamais pour accueillir quelqu'un. Pour souhaiter la bienvenue, on dit \eng{Welcome to our school!}
\item \eng{Could I ask you a few questions?} : \eng{could} est ici un conditionnel de politesse, pas le passé de \eng{can}. Ne traduisez pas « est-ce que tu pouvais ».
\item Ne dites pas \eng{*How malaria is transmitted?} : en question directe, il faut l'inversion du sujet et de l'auxiliaire — \eng{How \textbf{is} malaria transmitted?}
\item \eng{a few questions} signifie « quelques questions » (sens positif), alors que \eng{few questions} signifie « peu de questions » (sens négatif). Ne confondez pas.
\end{itemize}
\end{attention}

\courssub{Lecture à voix haute en binôme : l'intonation}

La compréhension ne suffit pas : il faut \emph{dire} le dialogue. Travaillez en binôme : l'un joue Abong, l'autre Mrs Etame, puis échangez les rôles.

\begin{propriete}[Intonation : monter ou descendre ?]
\begin{itemize}
\item \textbf{Questions en \eng{wh-}} (\eng{How}, \eng{What}, \eng{Where}) : l'intonation \textbf{descend} à la fin. \eng{How is malaria transmitted?} ↘
\item \textbf{Questions fermées} (réponse \eng{yes/no}) : l'intonation \textbf{monte} à la fin. \eng{Could I ask you a few questions?} ↗
\item \textbf{Phrases affirmatives et remerciements} : intonation descendante. \eng{Thank you for your time.} ↘
\item \textbf{Accentuation} : frappez le mot qui porte l'information nouvelle — \eng{It is spread by MOsquito bites, especially at NIGHT.}
\end{itemize}
\end{propriete}

\begin{experience}[Lecture dramatisée du dialogue]
Par groupes de deux : 1) lisez le dialogue en respectant l'intonation (montée sur la question de politesse, descente sur les questions en \eng{wh-}) ; 2) relisez en remplaçant les gestes par le ton : l'infirmière doit avoir l'air professionnelle et rassurante, l'élève curieux et poli ; 3) un binôme volontaire joue le dialogue devant la classe ; la classe note sur une grille les formules entendues (\eng{greeting / asking / giving / reacting / concluding}).
\end{experience}

\courssub{Substitution : remplacer malaria par cholera}

Dernière étape avant la production libre : la \cle{substitution}. On garde la structure du dialogue et on ne change que le contenu — ici, la maladie et les conseils adaptés.

\begin{exemplebox}[Dialogue transformé : le choléra]
\textbf{Abong:} Good morning, Mrs Etame. Could I ask you a few questions about \textbf{cholera}?

\textbf{Nurse:} Of course, Abong. Go ahead.

\textbf{Abong:} How is cholera transmitted?

\textbf{Nurse:} It is spread through \textbf{contaminated water and food}.

\textbf{Abong:} I see. And what can we do to prevent it at school?

\textbf{Nurse:} First, \textbf{wash your hands with soap}. Then, \textbf{drink only clean, boiled or treated water}. Finally, \textbf{go to the health centre immediately if you have diarrhoea}.

\textbf{Abong:} That's very useful. Thank you for your time, Mrs Etame.

\textbf{Nurse:} You're welcome. Stay healthy!
\end{exemplebox}

\begin{center}
\begin{tabularx}{\linewidth}{|X|X|X|}\hline
\rowcolor{popblueL} \textbf{Maladie} & \textbf{Transmission} & \textbf{Prévention} \\ \hline
malaria & spread by mosquito bites & sleep under treated nets, clear stagnant water \\ \hline
cholera & spread through contaminated water and food & wash hands with soap, drink clean water \\ \hline
typhoid fever & spread through dirty food and water & cook food well, wash fruit and vegetables \\ \hline
tuberculosis & spread through the air (coughing) & cover your mouth, ventilate rooms \\ \hline
\end{tabularx}
\end{center}

\begin{exoresolu}[Substitution guidée]
En t'inspirant du dialogue modèle, écris la réplique de l'infirmière si Abong pose la question : \eng{How is typhoid fever transmitted? And what can we do to prevent it?} Utilise \eng{First}, \eng{Then}, \eng{Finally}.
\tcblower
Réponse possible : \eng{Typhoid fever is spread through dirty food and water. First, wash your hands before eating. Then, cook food well and wash fruit and vegetables. Finally, drink clean water and go to the health centre if you have a high fever.}

« La typhoïde se transmet par la nourriture et l'eau souillées. D'abord, lavez-vous les mains avant de manger. Ensuite, faites bien cuire les aliments et lavez les fruits et légumes. Enfin, buvez de l'eau propre et allez au centre de santé si vous avez une forte fièvre. »
\end{exoresolu}

\begin{exercice}[À vous de jouer]
En binôme, reconstruisez le dialogue complet en remplaçant \eng{malaria} par \eng{tuberculosis}. Respectez les cinq étapes (salutation, demande, information, réaction, conclusion) et utilisez au moins deux marqueurs de séquence. Présentez votre dialogue à la classe en soignant l'intonation.
\end{exercice}

\begin{corrige}[Exercice — éléments de correction]
Un dialogue réussi contient : 1) une salutation adaptée au moment de la journée (\eng{Good morning/afternoon}) ; 2) une demande polie avec \eng{Could I ask you...?} ; 3) une réponse au passif : \eng{Tuberculosis is spread through the air, especially when people cough.} ; 4) des conseils structurés : \eng{First, cover your mouth when you cough. Then, open the windows to ventilate classrooms. Finally, see a doctor if you cough for more than two weeks.} ; 5) un remerciement et une conclusion : \eng{Thank you for your time. — You're welcome. Stay healthy!} Vérifiez l'inversion dans les questions et l'intonation descendante des questions en \eng{wh-}.
\end{corrige}

Vous savez désormais écouter, comprendre et rejouer un échange d'informations à l'infirmerie. La section suivante systématisera les \emph{formules fonctionnelles} de l'échange — demander, donner une information, relancer, conclure — pour que vous puissiez les mobiliser dans n'importe quelle situation : à l'école, au travail ou dans la communauté.

\coursec{Les formules fonctionnelles de l'échange d'informations}

Dans la section précédente, nous avons écouté et lu le dialogue entre Sarah et l'infirmière scolaire à propos du paludisme. Reprenons un instant ce dialogue : comment Sarah obtient-elle ses informations ? Comment l'infirmière les donne-t-elle ? Comment Sarah montre-t-elle qu'elle écoute, et comment la conversation se termine-t-elle ? Derrière chaque échange réussi se cachent des \cle{formules fonctionnelles} (\eng{functional language}) : des expressions toutes faites, polies et naturelles, que tout anglophone utilise sans réfléchir. C'est exactement ce que nous allons apprendre ici, étape par étape, pour que vous puissiez à votre tour échanger des informations sur la prévention des maladies transmissibles — à l'école, au travail ou dans votre communauté.

\courssub{Observer avant d'apprendre : un mini-dialogue}

\begin{textebox}[At the community health centre in Bamenda]
\textbf{Student:} Excuse me, nurse. Could you tell me how cholera is transmitted?

\textbf{Nurse:} Well, cholera is mainly transmitted through contaminated water and food.

\textbf{Student:} I see. And what causes the contamination of water?

\textbf{Nurse:} In fact, the main cause is poor sanitation — when toilets are too close to wells or streams.

\textbf{Student:} Really? I didn't know that! How can we prevent cholera in our neighbourhood?

\textbf{Nurse:} Let me explain. First of all, always boil or treat drinking water. Then, wash your hands with soap before eating. Also, keep your environment clean. Finally, report any case of severe diarrhoea to the health centre immediately.

\textbf{Student:} You mentioned treating water. Could you say more?

\textbf{Nurse:} Of course. You can boil the water for at least one minute, or use water purification tablets.

\textbf{Student:} That's very useful. Thank you for your time.

\textbf{Nurse:} You're welcome. Let's work together to keep our community healthy.
\end{textebox}

\textbf{Glossaire :} \emph{sanitation} (n.) : assainissement ; \emph{well} (n.) : puits ; \emph{stream} (n.) : ruisseau ; \emph{severe} (adj.) : grave ; \emph{diarrhoea} (n.) : diarrhée ; \emph{purification tablets} (n.) : pastilles de purification ; \emph{to treat water} (v.) : traiter l'eau.

Observez bien : chaque réplique joue un \cle{rôle précis} — demander, donner, relancer, réagir, conclure. Repérez dans le dialogue les expressions \eng{Could you tell me...?}, \eng{Well, ...}, \eng{In fact, ...}, \eng{First of all ... Then ... Also ... Finally}, \eng{I see.}, \eng{Really?}, \eng{Thank you for your time.} C'est cette architecture que nous allons maintenant démonter, fonction par fonction.

\courssub{Demander une information : \eng{asking for information}}

En anglais, on évite la question trop directe quand on s'adresse à une personne que l'on connaît peu (une infirmière, un chef de quartier, un collègue). On préfère des \cle{amorces polies} (\eng{polite openers}).

\begin{center}
\begin{tabularx}{\linewidth}{|X|X|X|}
\hline
\rowcolor{popblueL} Formule anglaise & Traduction française & Emploi \\
\hline
\eng{Could you tell me how...?} & Pourriez-vous me dire comment... ? & très poli, formel \\
\hline
\eng{I'd like to know why...} & J'aimerais savoir pourquoi... & poli, neutre \\
\hline
\eng{What causes...?} & Qu'est-ce qui cause... ? & direct, entre pairs \\
\hline
\eng{How can we prevent...?} & Comment pouvons-nous prévenir... ? & direct, courant \\
\hline
\eng{Do you know if...?} & Savez-vous si... ? & question fermée (oui/non) \\
\hline
\end{tabularx}
\end{center}

\begin{propriete}[La règle d'or : l'ordre des mots après une amorce]
Après \eng{Could you tell me...}, \eng{I'd like to know...}, \eng{Do you know...}, la question devient une \cle{question indirecte} : on revient à l'ordre de la phrase \textbf{affirmative} (sujet + verbe), et on supprime l'inversion et l'auxiliaire \eng{do/does/did}.

\begin{center}
\begin{tabular}{|l|l|}
\hline
\rowcolor{popblueL} Question directe & Question indirecte (correcte) \\
\hline
\eng{How IS malaria transmitted?} & \eng{Could you tell me how malaria IS transmitted?} \\
\hline
\eng{Why IS hygiene important?} & \eng{I'd like to know why hygiene IS important.} \\
\hline
\eng{How CAN we prevent cholera?} & \eng{Do you know how we CAN prevent cholera?} \\
\hline
\end{tabular}
\end{center}

Seul l'ordre change ; le sens reste identique. Remarquez que le mot \eng{if} (« si ») remplace le mot interrogatif dans les questions fermées (réponse oui/non) : \eng{Do you know if the vaccine is free?} « Savez-vous si le vaccin est gratuit ? »
\end{propriete}

\begin{exemplebox}[Exemples traduits]
\eng{I'd like to know why hygiene is so important.} « J'aimerais savoir pourquoi l'hygiène est si importante. »

\eng{Could you tell me how tuberculosis spreads?} « Pourriez-vous me dire comment la tuberculose se propage ? »

\eng{Do you know if the health centre is open on Saturdays?} « Savez-vous si le centre de santé est ouvert le samedi ? »

\eng{What causes typhoid fever?} « Qu'est-ce qui cause la fièvre typhoïde ? »
\end{exemplebox}

\begin{attention}[Erreur fréquente des francophones]
Le français garde l'inversion dans les questions indirectes (« Pourriez-vous me dire comment peut-on prévenir le choléra ? »), ce qui pousse les élèves à écrire : \eng{Could you tell me how CAN WE prevent cholera?} — \textbf{FAUX}. En anglais, après l'amorce, on revient à l'ordre sujet + verbe : \eng{Could you tell me how WE CAN prevent cholera?} — \textbf{CORRECT}. Autre piège : ne dites jamais \eng{I'd like to know why is hygiene important} ; dites \eng{why hygiene is important}. Enfin, n'oubliez pas que \eng{information} est \textbf{indénombrable} en anglais : on dit \eng{some information}, jamais \eng{an information} ni \eng{informations}.
\end{attention}

\courssub{Donner une information : \eng{giving information}}

Pour répondre, l'anglophone ne balance jamais l'information brutalement : il l'\cle{introduit} par une petite formule qui prépare l'auditeur.

\begin{itemize}
\item \eng{Well, ...} « Eh bien, ... » — amorce naturelle, presque obligatoire à l'oral : \eng{Well, malaria is transmitted by mosquito bites.} « Eh bien, le paludisme est transmis par les piqûres de moustiques. »
\item \eng{In fact, ...} « En fait, ... » — pour préciser ou corriger : \eng{In fact, the main cause is contaminated water.} « En fait, la cause principale est l'eau contaminée. »
\item \eng{According to the nurse, ...} « Selon l'infirmière, ... » — pour citer une source fiable : \eng{According to the nurse, we should sleep under treated mosquito nets.} « Selon l'infirmière, nous devrions dormir sous des moustiquaires imprégnées. »
\item \eng{The main cause is...} « La cause principale est... » — pour structurer l'explication.
\item \eng{Let me explain.} « Laissez-moi vous expliquer. » — pour annoncer une réponse détaillée.
\end{itemize}

\courssub{Structurer son propos : les mots de liaison}

Quand l'information comporte plusieurs étapes (par exemple les gestes de prévention), on les ordonne grâce aux \cle{connecteurs de séquence} :

\begin{center}
\begin{tabularx}{\linewidth}{|X|X|X|}
\hline
\rowcolor{popgreenL} Connecteur & Français & Position \\
\hline
\eng{First / First of all} & D'abord / Tout d'abord & premier point \\
\hline
\eng{Then} & Ensuite & étape suivante \\
\hline
\eng{Next} & Puis & étape suivante \\
\hline
\eng{Also} & Aussi, de plus & ajout d'une idée \\
\hline
\eng{Finally} & Enfin & dernier point \\
\hline
\end{tabularx}
\end{center}

\eng{First of all, wash your hands regularly. Then, drink only safe water. Also, keep your surroundings clean. Finally, get vaccinated.} « Tout d'abord, lavez-vous les mains régulièrement. Ensuite, ne buvez que de l'eau saine. De plus, gardez votre environnement propre. Enfin, faites-vous vacciner. »

\courssub{Relancer la conversation : \eng{keeping the conversation going}}

Un bon échange ne s'arrête pas à la première réponse. Pour obtenir plus de détails :

\begin{itemize}
\item \eng{What about...?} « Et qu'en est-il de... ? » : \eng{What about vaccination?} « Et la vaccination ? »
\item \eng{And what else?} « Et quoi d'autre ? »
\item \eng{Could you explain that again?} « Pourriez-vous réexpliquer cela ? »
\item \eng{You mentioned... Could you say more?} « Vous avez mentionné... Pourriez-vous en dire plus ? »
\end{itemize}

\begin{exemplebox}[Exemple traduit]
\eng{You mentioned vaccination. Could you say more?} « Vous avez mentionné la vaccination. Pourriez-vous en dire plus ? »
\end{exemplebox}

\begin{methode}[Relancer sans être impoli]
\begin{enumerate}
\item Évitez la question directe sèche (\eng{Why? And then?}), qui sonne comme un interrogatoire en anglais.
\item Préférez une amorce en \eng{Could you...?} ou \eng{What about...?}.
\item Ajoutez \eng{please} pour adoucir : \eng{Could you repeat that, please?}
\item Montrez d'abord que vous avez écouté en reprenant un mot de votre interlocuteur : \eng{You mentioned clean water...}
\end{enumerate}
\end{methode}

\courssub{Réagir et montrer l'intérêt : \eng{reacting}}

En anglais, un auditeur silencieux est perçu comme indifférent. Il faut \cle{ponctuer} l'écoute de petites réactions :

\begin{center}
\begin{tabularx}{\linewidth}{|X|X|X|}
\hline
\rowcolor{poporangeL} Réaction & Français & Nuance \\
\hline
\eng{Really?} & Vraiment ? & surprise (intonation montante) \\
\hline
\eng{I see.} & Je vois. & compréhension neutre \\
\hline
\eng{That's interesting!} & C'est intéressant ! & intérêt sincère \\
\hline
\eng{I didn't know that!} & Je ne savais pas ça ! & étonnement \\
\hline
\eng{Oh, I understand now.} & Oh, je comprends maintenant. & clarification reçue \\
\hline
\end{tabularx}
\end{center}

\courssub{Conclure l'échange : \eng{closing the conversation}}

On ne quitte jamais une conversation anglophone sans formule de clôture :

\begin{itemize}
\item \eng{Thank you for your help. / Thank you for your time.} « Merci pour votre aide / votre temps. »
\item \eng{That's very useful.} « C'est très utile. »
\item \eng{It was nice talking to you.} « C'était un plaisir de parler avec vous. »
\item \eng{Let's work together on this.} « Travaillons ensemble sur ce sujet. »
\end{itemize}

\courssub{Vue d'ensemble : la carte des quatre fonctions}

\begin{popfigure}
\begin{center}
\begin{tikzpicture}[
  box/.style={draw=popdark, rounded corners, fill=#1, text width=3.6cm, align=center, font=\small, inner sep=5pt}
]
\node[box=popblue, align=center] (ask) at (0,0){\textcolor{white}{\textbf{ASKING}}\\[2pt]
\eng{Could you tell me how...?}\\
\eng{I'd like to know why...}\\
\eng{Do you know if...?}};
\node[box=popgreen, align=center] (give) at (5.2,0){\textcolor{white}{\textbf{GIVING}}\\[2pt]
\eng{Well, ...}\\
\eng{In fact, ...}\\
\eng{According to the nurse, ...}};
\node[box=poporange, align=center] (react) at (0,-3.4){\textcolor{white}{\textbf{REACTING}}\\[2pt]
\eng{Really?}\\
\eng{I see.}\\
\eng{That's interesting!}};
\node[box=poppurple, align=center] (close) at (5.2,-3.4){\textcolor{white}{\textbf{CONCLUDING}}\\[2pt]
\eng{Thank you for your time.}\\
\eng{That's very useful.}\\
\eng{It was nice talking to you.}};
\draw[-{Stealth}, thick, popdark] (ask) -- (give);
\draw[-{Stealth}, thick, popdark] (give) -- (react);
\draw[-{Stealth}, thick, popdark] (react) -- (close);
\end{tikzpicture}
\end{center}
\legende{Les quatre fonctions d'un échange d'informations, avec trois formules clés par fonction.}
\end{popfigure}

\begin{savaistu}[La politesse anglophone : ne jamais couper la parole]
Dans la culture anglophone, interrompre quelqu'un est très mal perçu, même pour une bonne raison. On attend une \cle{pause} naturelle, puis on intervient avec une formule d'excuse : \eng{Sorry, may I add something?} « Pardon, puis-je ajouter quelque chose ? » ou \eng{Excuse me, can I ask a question?} « Excusez-moi, puis-je poser une question ? » Cette règle vaut aussi en classe et lors des réunions communautaires : celui qui coupe la parole passe pour impoli, même si son idée est excellente.
\end{savaistu}

\begin{exoresolu}[Remettre le dialogue en ordre]
Énoncé — Voici les répliques d'un dialogue entre un élève et un agent de santé à Douala. Remettez-les dans l'ordre logique (1 à 5) et identifiez la fonction de chaque réplique.

a) \eng{Well, the main cause is stagnant water where mosquitoes breed.}

b) \eng{Thank you for your help. That's very useful.}

c) \eng{Could you tell me how we can prevent malaria in our neighbourhood?}

d) \eng{I see. And what else can we do?}

e) \eng{First of all, sleep under treated mosquito nets. Then, clear stagnant water around your houses.}

\tcblower
Solution — Ordre correct : \textbf{c -- a -- d -- e -- b}.

1. c : l'élève \textbf{demande} une information avec une question indirecte correcte (\eng{how we can prevent}, ordre affirmatif).

2. a : l'agent \textbf{donne} une première information, introduite par \eng{Well,} et structurée par \eng{the main cause is...}.

3. d : l'élève \textbf{réagit} (\eng{I see.}) puis \textbf{relance} (\eng{And what else can we do?}).

4. e : l'agent \textbf{développe} avec les connecteurs \eng{First of all} et \eng{Then}.

5. b : l'élève \textbf{conclut} poliment (\eng{Thank you for your help. That's very useful.}).

Ce dialogue suit parfaitement le schéma : Asking $\rightarrow$ Giving $\rightarrow$ Reacting $\rightarrow$ Concluding.
\end{exoresolu}

\begin{exercice}[À vous de jouer]
Complétez ce dialogue avec les formules fonctionnelles étudiées.

\textbf{A:} Excuse me. \_\_\_\_\_\_\_\_\_\_\_\_ tell me how typhoid is transmitted? (demande polie)

\textbf{B:} \_\_\_\_\_\_\_\_\_\_\_\_, it is transmitted through dirty food and water. (amorce)

\textbf{A:} \_\_\_\_\_\_\_\_\_\_\_\_? I didn't know that! (surprise)

\textbf{B:} \_\_\_\_\_\_\_\_\_\_\_\_, always wash your hands. \_\_\_\_\_\_\_\_\_\_\_\_, drink boiled water. \_\_\_\_\_\_\_\_\_\_\_\_, avoid street food during epidemics. (structure en trois étapes)

\textbf{A:} \_\_\_\_\_\_\_\_\_\_\_\_ for your time. It was nice talking to you. (conclusion)
\end{exercice}

\begin{corrige}[Exercice : À vous de jouer]
\textbf{A:} \eng{Could you} tell me how typhoid is transmitted? (attention : \eng{is transmitted}, ordre affirmatif, pas d'inversion)

\textbf{B:} \eng{Well}, it is transmitted through dirty food and water.

\textbf{A:} \eng{Really}? I didn't know that!

\textbf{B:} \eng{First of all}, always wash your hands. \eng{Then}, drink boiled water. \eng{Finally}, avoid street food during epidemics.

\textbf{A:} \eng{Thank you} for your time. It was nice talking to you.
\end{corrige}

\begin{aretenir}[L'essentiel de la section]
\begin{itemize}
\item Demander : amorce polie + \textbf{ordre affirmatif} : \eng{Could you tell me how malaria IS transmitted?}
\item Donner : introduire par \eng{Well, / In fact, / According to...} et structurer avec \eng{First, Then, Also, Finally}.
\item Relancer : \eng{What about...? / You mentioned... Could you say more?} — toujours avec douceur.
\item Réagir : \eng{Really? / I see. / That's interesting!} pour montrer qu'on écoute.
\item Conclure : \eng{Thank you for your time. / It was nice talking to you.}
\end{itemize}
\end{aretenir}

Maintenant que vous maîtrisez ces formules, nous allons les remettre en situation dans un deuxième dialogue modèle, cette fois dans la communauté et au travail.

\coursec{Dialogue modèle 2 : dans la communauté et au travail}

Dans la section précédente, nous avons découvert les formules fonctionnelles de l'échange d'informations : demander (\eng{Could you tell me...?}), donner une information (\eng{First... Then... Also...}), relancer (\eng{What about...?}) et conclure (\eng{Thank you for your time}). Mais ces formules ne s'emploient pas de la même façon avec tout le monde : on ne parle pas à un agent de santé comme à son voisin de classe ! Dans cette section, nous allons observer un dialogue formel entre une journaliste et un agent de santé communautaire, puis une variante semi-formelle entre collègues de travail, afin d'apprendre à \cle{adapter son registre de langue} au contexte et à vérifier sa compréhension par la \cle{reformulation}.

\courssub{Dialogue modèle : une interview sur le choléra dans le quartier}

\begin{textebox}[At the community health post — an interview about cholera]
Larissa: Excuse me, Mr Ndi. I'm a reporter for the school radio. Could you tell me how cholera spreads in our community?

Mr Ndi: Certainly. Cholera spreads mainly through contaminated water and food.

Larissa: I see. And what should people do to protect themselves?

Mr Ndi: First, drink boiled or treated water. Then, wash your hands with soap before eating. Also, cook food properly.

Larissa: So, you mean that dirty water is the main problem?

Mr Ndi: Exactly. If I understand correctly, your listeners are students, so tell them to avoid street food that is not covered.

Larissa: What about the markets and workplaces?

Mr Ndi: Sellers must keep their stalls clean, and employers should provide clean toilets and soap.

Larissa: That's very clear. Thank you very much for your time.

Mr Ndi: My pleasure. Spread the word, not the disease!
\end{textebox}

\textbf{Glossaire}\par
\begin{itemize}
\item \eng{to spread} (verbe irrégulier : \eng{spread – spread – spread}) : se propager, se répandre
\item \eng{contaminated} (adjectif) : contaminé, souillé
\item \eng{boiled water} (nom) : eau bouillie
\item \eng{treated water} (nom) : eau traitée
\item \eng{a stall} (nom) : un étal (au marché)
\item \eng{an employer} (nom) : un employeur
\item \eng{to provide} (verbe) : fournir
\item \eng{a listener} (nom) : un auditeur
\item \eng{street food} (nom) : nourriture de rue
\end{itemize}

\textbf{Questions de compréhension}\par
\begin{enumerate}
\item \eng{Who is Larissa and why is she talking to Mr Ndi?}
\item \eng{How does cholera spread, according to Mr Ndi?}
\item \eng{What three pieces of advice does Mr Ndi give to individuals?}
\item \eng{What should sellers and employers do?}
\item \eng{How does Larissa check that she has understood?}
\end{enumerate}

\begin{corrige}[Questions de compréhension]
\begin{enumerate}
\item \eng{She is a reporter for the school radio; she is interviewing him about cholera.} (C'est une journaliste de la radio scolaire ; elle l'interviewe au sujet du choléra.)
\item \eng{It spreads mainly through contaminated water and food.} (Il se propage surtout par l'eau et la nourriture contaminées.)
\item \eng{Drink boiled or treated water, wash your hands with soap before eating, and cook food properly.} (Boire de l'eau bouillie ou traitée, se laver les mains au savon avant de manger, et bien cuire les aliments.)
\item \eng{Sellers must keep their stalls clean, and employers should provide clean toilets and soap.} (Les vendeurs doivent garder leurs étals propres, et les employeurs devraient fournir des toilettes propres et du savon.)
\item \eng{She reformulates: “So, you mean that dirty water is the main problem?”} (Elle reformule : « Donc, vous voulez dire que l'eau sale est le problème principal ? »)
\end{enumerate}
\end{corrige}

\courssub{Observer : la reformulation pour vérifier sa compréhension}

Dans le dialogue, Larissa ne se contente pas d'écouter : elle \emph{reformule} ce qu'elle a compris. Observons :

\eng{So, you mean that dirty water is the main problem?} = « Donc, vous voulez dire que l'eau sale est le problème principal ? »

\eng{If I understand correctly, your listeners are students.} = « Si je comprends bien, vos auditeurs sont des élèves. »

\begin{propriete}[Les formules de reformulation]
Pour vérifier qu'on a bien compris une information, on reformule avec l'une de ces structures :
\begin{itemize}
\item \eng{So, you mean that + phrase ?} = « Donc, vous voulez dire que... ? » (semi-formel)
\item \eng{If I understand correctly, + phrase.} = « Si je comprends bien, ... » (formel)
\item \eng{In other words, + phrase ?} = « En d'autres termes, ... ? » (formel)
\item \eng{Do you mean + nom / -ing ?} = « Voulez-vous dire... ? » (neutre)
\end{itemize}
La reformulation se fait souvent avec une \cle{intonation montante} à la fin, car c'est une question de vérification. Remarquez qu'après \eng{Do you mean}, on emploie un nom ou un verbe en \eng{-ing} : \eng{Do you mean washing hands?} et non \eng{Do you mean wash hands?}.
\end{propriete}

\begin{exemplebox}[Exemples traduits]
\eng{So, you mean that dirty water is the main problem?} = « Donc, vous voulez dire que l'eau sale est le problème principal ? »

\eng{Employers should provide clean toilets.} = « Les employeurs devraient fournir des toilettes propres. »

\eng{If I understand correctly, we must boil the water before drinking it.} = « Si je comprends bien, nous devons faire bouillir l'eau avant de la boire. »

\eng{Do you mean washing hands five times a day?} = « Vous voulez dire se laver les mains cinq fois par jour ? »

\eng{In other words, the market needs more dustbins?} = « En d'autres termes, le marché a besoin de plus de poubelles ? »
\end{exemplebox}

\begin{methode}[Vérifier sa compréhension en trois étapes]
\begin{enumerate}
\item \textbf{Écouter} l'information sans interrompre, en montrant qu'on suit : \eng{I see. / Right. / I understand.}
\item \textbf{Reformuler} avec ses propres mots : \eng{So, you mean that...?} ou \eng{If I understand correctly, ...}.
\item \textbf{Attendre la confirmation} : l'interlocuteur répond \eng{Exactly. / That's right.} ou corrige : \eng{Not exactly. I said...}.
\end{enumerate}
Cette stratégie est très valorisée à l'oral du Baccalauréat : elle montre que vous savez \emph{interagir}, pas seulement réciter.
\end{methode}

\courssub{Variante : un échange semi-formel entre collègues de travail}

Le même contenu (prévention des maladies transmissibles) peut être échangé dans un registre plus détendu. Écoutons deux collègues d'une entreprise de Douala.

\begin{textebox}[At the office — talking about hygiene at work]
Amina: Hey, Paul, do you know why there are hand sanitisers at every door now?

Paul: Yeah! The boss wants to stop the flu and COVID-19 from spreading at work.

Amina: Really? So, you mean that we have to sanitise our hands every time we come in?

Paul: That's right. And if you feel sick, you should stay at home. Don't come to the office with a fever!

Amina: OK. And what about the meetings? The meeting room is so small.

Paul: Good question. We should open the windows and keep some distance. If I understand correctly, the manager will also buy masks for everyone.

Amina: Great idea. Thanks for the information!

Paul: No problem. Stay healthy!
\end{textebox}

\textbf{Glossaire}\par
\begin{itemize}
\item \eng{hand sanitiser} (nom) : gel hydroalcoolique
\item \eng{the flu} (nom) : la grippe
\item \eng{to feel sick} (expression) : se sentir malade
\item \eng{a fever} (nom) : une fièvre
\item \eng{to keep some distance} (expression) : garder ses distances
\item \eng{the boss / the manager} (nom) : le patron / le directeur
\end{itemize}

\courssub{Comparaison des registres : formel, semi-formel, informel}

Comparons maintenant les deux dialogues. Le contenu est proche (prévention des maladies), mais la \emph{forme} change selon l'interlocuteur et le lieu.

\begin{center}
\begin{tabularx}{\linewidth}{|X|X|X|}
\hline
\rowcolor{popblueL} Formal (nurse, authorities, journalist) & Semi-formal / informal (classmates, colleagues, friends) & Fonction \\
\hline
\eng{Excuse me, Sir / Mr Ndi.} & \eng{Hey, Paul!} & S'adresser à quelqu'un \\
\hline
\eng{Could you tell me how cholera spreads?} & \eng{Do you know why there are sanitisers?} & Demander une information \\
\hline
\eng{What should people do to protect themselves?} & \eng{What about the meetings?} & Relancer \\
\hline
\eng{That's very clear. Thank you very much for your time.} & \eng{Great idea. Thanks!} & Remercier \\
\hline
\eng{My pleasure.} & \eng{No problem.} & Répondre au remerciement \\
\hline
\end{tabularx}
\end{center}

\begin{propriete}[Adapter les formules au contexte]
\begin{itemize}
\item \textbf{Registre formel} (infirmière, autorités, journaliste, inconnu âgé) : formules complètes et polies — \eng{Excuse me, Sir...}, \eng{Could you tell me...?}, \eng{Would you mind explaining...?}, \eng{Thank you very much for your time.}
\item \textbf{Registre semi-formel} (camarades de classe, collègues) : formules directes mais correctes — \eng{Do you know...?}, \eng{Can you tell me...?}, \eng{Thanks a lot.}
\item \textbf{Registre informel} (amis proches) : raccourcis et interjections — \eng{Hey, do you know...?}, \eng{Yeah!}, \eng{No problem.}
\end{itemize}
En anglais, le choix du registre est marqué par la \cle{formule d'ouverture}, les \cle{contractions} (\eng{I'm, don't, yeah}) et la longueur des phrases. Notez aussi la nuance entre \eng{must} (obligation forte, règle) et \eng{should} (conseil) : \eng{Sellers must keep their stalls clean} (c'est une obligation) ; \eng{Employers should provide clean toilets} (c'est un devoir moral, un conseil).
\end{propriete}

\begin{popfigure}
\begin{center}
\begin{tikzpicture}
\draw[-{Stealth}, line width=2pt, popmuted] (0,0) -- (11,0);
\node[draw=popgreen, fill=popgreenL, rounded corners, align=center, minimum width=3cm, minimum height=1.4cm] at (1.8,1.3) {\textbf{Informal}\\ \eng{friends}\\ \eng{“Hey, do you know...?”}};
\node[draw=poporange, fill=poporangeL, rounded corners, align=center, minimum width=3cm, minimum height=1.4cm] at (5.5,1.3) {\textbf{Semi-formal}\\ \eng{classmates, colleagues}\\ \eng{“Can you tell me...?”}};
\node[draw=popblue, fill=popblueL, rounded corners, align=center, minimum width=3cm, minimum height=1.4cm] at (9.2,1.3) {\textbf{Formal}\\ \eng{nurse, authorities}\\ \eng{“Excuse me, Sir...”}};
\fill[popgreen] (1.8,0) circle (3pt);
\fill[poporange] (5.5,0) circle (3pt);
\fill[popblue] (9.2,0) circle (3pt);
\node[below, popmuted] at (1.8,-0.1) {\emph{proche}};
\node[below, popmuted] at (9.2,-0.1) {\emph{distant}};
\end{tikzpicture}
\end{center}
\legende{L'échelle du registre de langue : plus l'interlocuteur est distant (statut, âge, inconnus), plus le langage est formel.}
\end{popfigure}

\begin{exoresolu}[Choisir le bon registre]
Énoncé : pour chaque situation, choisissez la formule adaptée (a ou b).

1. You interview a doctor at the hospital. \quad a) \eng{Hey, what's cholera?} \quad b) \eng{Excuse me, Doctor. Could you tell me what cholera is?}

2. You ask your deskmate about the homework. \quad a) \eng{Hey, do you know what the homework is?} \quad b) \eng{Excuse me, Sir. Would you mind telling me the homework?}

3. You thank a community health worker after an interview. \quad a) \eng{Thank you very much for your time.} \quad b) \eng{Yeah, thanks, bye!}

4. You check your understanding with a classmate. \quad a) \eng{So, you mean that the meeting is cancelled?} \quad b) \eng{I would be grateful if you could confirm your previous statement.}
\tcblower
Solution détaillée :
\begin{enumerate}
\item \textbf{b)} — un médecin est une autorité : registre formel (\eng{Excuse me}, \eng{Could you tell me...?}).
\item \textbf{a)} — un voisin de classe est un pair : registre informel ou semi-formel ; \eng{Sir} serait ridicule ici.
\item \textbf{a)} — après une interview, on remercie formellement pour le temps accordé.
\item \textbf{a)} — entre camarades, \eng{So, you mean that...?} est naturel ; la formule b) est lourde et artificielle entre pairs.
\end{enumerate}
\end{exoresolu}

\begin{attention}[Faux amis et traductions mot à mot]
\begin{itemize}
\item \eng{Spread the word!} = « Fais passer le message ! » — ne traduisez jamais « répands le mot ». \eng{Word} signifie ici « nouvelle, information ».
\item \eng{To catch a disease} = « attraper une maladie » : \eng{catch} ne signifie pas ici « attraper un objet ». On dit \eng{catch a cold / catch the flu}, mais pour les maladies graves on préfère \eng{contract a disease} (registre formel).
\item \eng{Actually} = « en fait », et non « actuellement » (\eng{currently}).
\item Attention à \eng{sensible} / \eng{sensitive} : \eng{sensitive} = « sensible » ; « sensé » se dit \eng{sensible}. \eng{It is sensible to boil water.} = « Il est sensé de faire bouillir l'eau. »
\end{itemize}
\end{attention}

\begin{savaistu}[“Spread the word, not the disease!”]
Ce slogan est typique des campagnes de santé publique anglophones. Les slogans publicitaires et sanitaires anglais utilisent souvent la \cle{rime} ou le \cle{parallélisme} pour être mémorables : ici, la structure \eng{Spread the X, not the Y} oppose deux noms (\eng{word} / \eng{disease}) autour du même verbe. Au Cameroun anglophone comme dans de nombreux pays, les radios communautaires jouent un rôle central dans la prévention des épidémies : un message simple, répété et bien formulé sauve des vies. À vous d'inventer votre propre slogan en classe !
\end{savaistu}

\begin{exercice}[À vous de jouer]
\begin{enumerate}
\item Reformulez chaque information avec \eng{So, you mean that...?} :
\begin{enumerate}
\item \eng{People must wash their hands with soap before eating.}
\item \eng{Employers should provide clean toilets at work.}
\end{enumerate}
\item Réécrivez ce dialogue informel en registre formel (journaliste / agent de santé) : \eng{— Hey, why is everyone washing their hands? — Because of cholera, man!}
\item Préparez avec un camarade un mini-dialogue : l'un est journaliste de la radio scolaire, l'autre est le délégué d'hygiène du lycée. Utilisez au moins une formule de reformulation et une conclusion polie.
\end{enumerate}
\end{exercice}

\begin{corrige}[Exercice — À vous de jouer]
\begin{enumerate}
\item a) \eng{So, you mean that washing hands with soap is necessary before every meal?} \quad b) \eng{So, you mean that it is the employers' duty to provide clean toilets?}
\item \eng{— Excuse me, Sir. Could you tell me why everyone is washing their hands? — Certainly. It is a measure to prevent cholera.}
\item Modèle possible : \eng{— Excuse me, could you tell me how our school prevents diseases? — Of course. We clean the toilets every day and we provide soap. — So, you mean that hygiene is everyone's business? — Exactly! — Thank you for your time. — You're welcome.}
\end{enumerate}
\end{corrige}

\begin{aretenir}[L'essentiel de la section]
\begin{itemize}
\item On adapte son \cle{registre} à l'interlocuteur : \eng{Excuse me, Sir, could you tell me...?} (formel) vs \eng{Hey, do you know...?} (informel).
\item Pour vérifier sa compréhension, on \cle{reformule} : \eng{So, you mean that...?} / \eng{If I understand correctly, ...}.
\item On conclut un échange formel par \eng{Thank you very much for your time.} et l'interlocuteur répond \eng{My pleasure.} ou \eng{You're welcome.}
\item \eng{Spread the word} = « faire passer le message » ; \eng{catch a disease} = « attraper une maladie » : attention aux traductions mot à mot !
\end{itemize}
\end{aretenir}

\coursec{Prononciation, accentuation et intonation}

Nous venons de travailler sur le deuxième dialogue, celui qui se déroule dans la communauté et au travail. Vous connaissez maintenant les formules pour demander et donner une information sur la prévention des maladies transmissibles. Mais un dialogue réussi ne dépend pas seulement des mots choisis : il dépend aussi de la manière dont on les prononce. En anglais, un mot mal accentué ou une question posée avec la mauvaise intonation peut rendre votre message incompréhensible, même si la grammaire est parfaite. Cette section va donc entraîner votre oreille et votre voix sur trois points précis : l'\cle{accentuation} des mots (word stress), la prononciation des mots difficiles du champ lexical de la santé, et l'\cle{intonation} des questions et des énumérations.

\courssub{Observer d'abord : écouter la musique de la langue}

Lisez à voix haute les deux répliques suivantes, extraites de nos dialogues modèles. Les flèches indiquent le mouvement de la voix.

\eng{Do you sleep under a mosquito net? ↗} « Dort-on sous une moustiquaire ? »

\eng{How is cholera transmitted? ↘} « Comment le choléra se transmet-il ? »

Vous remarquez deux choses. D'abord, à l'intérieur de chaque mot, une syllabe est prononcée plus fort que les autres : mos\cle{QUI}to, \cle{CHO}lera, trans\cle{MI}tted. Ensuite, la voix ne reste pas plate : elle \emph{monte} à la fin de la première question et \emph{descend} à la fin de la seconde. Ce n'est pas un hasard : c'est une règle générale de l'anglais parlé, et nous allons l'étudier méthodiquement.

\begin{methode}[Bien accentuer un mot anglais]
En anglais, chaque mot de plus d'une syllabe possède une \cle{syllabe accentuée} (stressed syllable), prononcée plus fort, plus longtemps et plus haut que les autres. Dans cette leçon, nous l'écrivons en MAJUSCULES.
\begin{enumerate}
\item Découpez le mot en syllabes : pre-ven-tion.
\item Repérez la syllabe accentuée dans le dictionnaire ou dans la leçon : pre-VEN-tion.
\item Frappez cette syllabe en la disant plus fort et plus longuement ; les autres syllabes sont rapides et faibles.
\item Vérifiez en frappant dans vos mains : un coup fort sur la syllabe accentuée, des coups légers sur les autres.
\end{enumerate}
Mal placer l'accent peut rendre le mot incompréhensible pour un anglophone : dites \eng{preVENtion} et non « pré-ven-TION » à la française.
\end{methode}

\courssub{L'accentuation des mots clés de la leçon}

Voici les huit mots essentiels de nos dialogues, avec leur accentuation. Entraînez-vous à les dire à voix haute, d'abord en chœur avec la classe, puis individuellement.

\begin{center}
\begin{tabularx}{\linewidth}{|X|X|X|}\hline
\rowcolor{popblueL} Mot anglais & Accentuation & Traduction \\ \hline
disease & diSEASE (« di-ZIZ ») & maladie \\ \hline
prevention & preVENtion (« pri-VEN-cheun ») & prévention \\ \hline
hygiene & HYgiene (« HAÏ-djine ») & hygiène \\ \hline
vaccination & vacciNAtion (« va-ksi-NEI-cheun ») & vaccination \\ \hline
mosquito & mosQUIto (« moss-KI-to ») & moustique \\ \hline
contaminated & conTAMinated (« kon-TA-mi-né-tid ») & contaminé \\ \hline
community & comMUnity (« ko-MIOU-ni-ti ») & communauté \\ \hline
epidemic & epiDEmic (« é-pi-DÉ-mik ») & épidémie \\ \hline
\end{tabularx}
\end{center}

\begin{attention}[L'accent anglais n'est pas l'accent français]
En français, l'accent tombe presque toujours sur la \emph{dernière} syllabe du groupe de mots : « la prévenTION », « la commuNAUTÉ ». En anglais, l'accent tombe souvent sur la \emph{première} syllabe ou sur une syllabe du milieu, et il est beaucoup plus marqué. Les erreurs typiques des francophones :
\begin{itemize}
\item « pré-ven-TION » au lieu de \eng{preVENtion} ;
\item « commu-ni-TÉ » au lieu de \eng{comMUnity} ;
\item « é-pi-dé-MIK » régulier au lieu de \eng{epiDEmic}.
\end{itemize}
Astuce : quand vous apprenez un mot nouveau, apprenez toujours son accent en même temps que son orthographe, comme s'il faisait partie du mot.
\end{attention}

\courssub{Les mots pièges du vocabulaire de la santé}

Certains mots de ce chapitre ont une prononciation qui surprend les francophones, soit à cause de l'orthographe, soit à cause d'un faux ami avec le français.

\begin{center}
\begin{tabularx}{\linewidth}{|X|X|X|}\hline
\rowcolor{popblueL} Mot & Prononciation (lettres françaises) & Piège à éviter \\ \hline
cholera & « KO-lé-ra » & ne pas dire « cho-lé-ra » : \eng{ch} se dit « k » ici \\ \hline
typhoid & « TAÏ-foïd » & ne pas dire « ti-foïd » à la française \\ \hline
diarrhoea & « daïo-RI-a » & orthographe britannique en \eng{-oea}, accent sur « RI » \\ \hline
cough & « KOF » & \eng{-ough} se prononce « of », jamais « oug » \\ \hline
fever & « FI-veur » & accent sur la première syllabe, \eng{v} bien sonore \\ \hline
\end{tabularx}
\end{center}

\begin{exemplebox}[Phrases d'entraînement]
Répétez chaque phrase en insistant sur la syllabe en majuscules.
\begin{itemize}
\item \eng{CHOlera is caused by conTAMinated water.} « Le choléra est causé par l'eau contaminée. »
\item \eng{He has a high FIver and a bad COUGH.} « Il a une forte fièvre et une mauvaise toux. »
\item \eng{VacciNAtion is the best preVENtion.} « La vaccination est la meilleure prévention. »
\item \eng{An epiDEmic can spread quickly in the comMUnity.} « Une épidémie peut se propager vite dans la communauté. »
\end{itemize}
\end{exemplebox}

\begin{attention}[Le « h » et le « th » anglais]
Deux sons n'existent pas en français et demandent un entraînement spécial.
\begin{itemize}
\item Le \eng{h} anglais \emph{se prononce toujours}, comme un souffle léger : \eng{hygiene} (« HAÏ-djine »), \eng{hospital} (« HOSS-pi-tol »), \eng{hands} (« hendz »). Ne dites jamais « ygiène » ou « ospital » comme en français.
\item Le \eng{th} de \eng{health} (« HELSS ») se prononce avec la pointe de la langue entre les dents, en soufflant. Ne le remplacez ni par « s » ni par « z » : dites \eng{health}, \eng{three}, \eng{method} avec la langue entre les dents.
\end{itemize}
\end{attention}

\begin{savaistu}[Le mystérieux groupe « -ough »]
Le mot \eng{cough} se prononce « KOF », mais \eng{though} se prononce « ZO », \eng{through} se dit « SROU », \eng{enough} se dit « i-NEUF », \eng{bought} se dit « BOTE » et \eng{plough} se dit « PLAO ». Le groupe \eng{-ough} possède plus de six prononciations différentes en anglais ! C'est un héritage de l'histoire de la langue : l'orthographe anglaise a été fixée avant que la prononciation ne cesse d'évoluer. Moralité : vérifiez toujours la prononciation d'un mot nouveau dans un dictionnaire.
\end{savaistu}

\courssub{L'intonation des questions : monter ou descendre ?}

Observez ces deux questions, tirées de nos dialogues.

\eng{Do you wash your hands with soap? ↗} « Te laves-tu les mains avec du savon ? » (la voix monte)

\eng{Why is hygiene important? ↘} « Pourquoi l'hygiène est-elle importante ? » (la voix descend)

Pourquoi cette différence ? Parce que la première question attend une réponse \emph{oui} ou \emph{non}, tandis que la seconde demande une information précise. L'anglais code cette différence par la mélodie de la phrase.

\begin{propriete}[La règle d'intonation des questions]
\begin{itemize}
\item \cle{Questions fermées} (réponse yes/no, commençant par un auxiliaire : do, does, is, are, can, have...) : la voix \textbf{monte} sur le dernier mot accentué ↗. Exemple : \eng{Is malaria dangerous? ↗} « Le paludisme est-il dangereux ? »
\item \cle{Questions ouvertes} (commençant par un mot en WH- : what, how, why, where, when, who) : la voix \textbf{descend} sur le dernier mot accentué ↘. Exemple : \eng{What causes malaria? ↘} « Qu'est-ce qui cause le paludisme ? »
\end{itemize}
Exception : une question en WH- peut monter si l'on demande poliment de répéter une information : \eng{When did you say the vaccination starts? ↗} « Quand avez-vous dit que la vaccination commence, déjà ? »
\end{propriete}

\begin{popfigure}
\begin{center}
\begin{tikzpicture}[scale=1]
% Yes/No question : montante
\node[align=center, font=\small\bfseries, color=popdark] at (3,2.6) {Do you sleep under a mosquito net?};
\draw[->, very thick, popblue] (0.6,1.2) .. controls (2.2,1.3) and (3.6,1.6) .. (5.2,2.3);
\node[align=center, font=\footnotesize, color=popblue] at (3,0.7) {Question Yes/No : la voix monte ↗};
% WH- question : descendante
\node[align=center, font=\small\bfseries, color=popdark] at (9.4,2.6) {How is cholera transmitted?};
\draw[->, very thick, poporange] (7.2,2.2) .. controls (8.6,1.8) and (10.2,1.4) .. (11.6,0.9);
\node[align=center, font=\footnotesize, color=poporange] at (9.4,0.4) {Question en WH- : la voix descend ↘};
\end{tikzpicture}
\end{center}
\legende{Les deux mélodies fondamentales des questions anglaises : montante pour les questions fermées, descendante pour les questions ouvertes.}
\end{popfigure}

\begin{center}
\begin{tabularx}{\linewidth}{|X|X|X|}\hline
\rowcolor{popblueL} Type de question & Exemple anglais & Intonation \\ \hline
Fermée (yes/no) & \eng{Do you boil the water?} & montante ↗ \\ \hline
Fermée (yes/no) & \eng{Is the clinic open today?} & montante ↗ \\ \hline
Ouverte (WH-) & \eng{Where is the health centre?} & descendante ↘ \\ \hline
Ouverte (WH-) & \eng{How can we prevent typhoid?} & descendante ↘ \\ \hline
\end{tabularx}
\end{center}

Comparaison avec le français : en français, la voix monte presque toujours à la fin de \emph{toutes} les questions (« Tu viens ? », « Où vas-tu ? »). Le grand piège du francophone est donc de faire monter la voix sur les questions en WH- anglaises, ce qui sonne étrange ou hésitant pour un anglophone. Entraînez-vous à \emph{laisser tomber} la voix sur le dernier mot : \eng{What causes maLAria? ↘}.

\courssub{L'intonation de l'énumération}

Quand on donne une liste de conseils ou d'étapes, comme dans nos dialogues de prévention, l'anglais utilise une mélodie très régulière.

\begin{propriete}[Règle de l'énumération]
Dans une liste, la voix \textbf{monte} sur chaque élément, sauf le dernier, où elle \textbf{descend}. La chute finale signale à l'auditeur : « la liste est terminée ».

\eng{First, wash your hands ↗ ; then, boil the water ↗ ; finally, cook the food well ↘.} « D'abord, lavez-vous les mains ; ensuite, faites bouillir l'eau ; enfin, faites bien cuire la nourriture. »
\end{propriete}

\begin{exemplebox}[Autres énumérations à pratiquer]
\begin{itemize}
\item \eng{We need soap ↗, clean water ↗, and mosquito nets ↘.} « Il nous faut du savon, de l'eau propre et des moustiquaires. »
\item \eng{The nurse checks your temperature ↗, your blood pressure ↗, and your weight ↘.} « L'infirmière vérifie ta température, ta tension et ton poids. »
\end{itemize}
\end{exemplebox}

\courssub{Entraînement guidé : de la chorale au binôme}

\begin{experience}[Choral drilling puis travail en binôme]
\begin{enumerate}
\item \textbf{En chœur} : la classe répète après l'enseignant les répliques clés des deux dialogues, en tapant l'accent des mots dans les mains : \eng{What are the SYMptoms of CHOlera? ↘} — \eng{Do you SLEEP under a mosQUIto net? ↗}
\item \textbf{Par rangées} : une rangée pose la question fermée (voix montante), l'autre répond ; puis on inverse avec une question en WH- (voix descendante).
\item \textbf{En binôme} : l'élève A lit l'énumération des conseils d'hygiène, l'élève B ferme les yeux et lève la main quand il entend la voix descendre (signal de fin de liste). Puis on échange les rôles.
\item \textbf{Contrôle qualité} : chaque binôme choisit sa meilleure question et la dit devant la classe ; les autres élèves doivent dire si la voix montait ou descendait, et pourquoi.
\end{enumerate}
\end{experience}

\begin{exoresolu}[Identifier l'intonation correcte]
Énoncé. Pour chaque question, indiquez si la voix monte ↗ ou descend ↘, puis traduisez-la. 1) \eng{Is the water contaminated?} 2) \eng{What are the symptoms of typhoid?} 3) \eng{Have you been vaccinated?} 4) \eng{When does the vaccination campaign start?} 5) \eng{Can malaria be prevented?}
\tcblower
Solution détaillée.
\begin{enumerate}
\item \eng{Is the water contaminated? ↗} — question fermée (auxiliaire \eng{is}, réponse yes/no) : la voix monte. « L'eau est-elle contaminée ? »
\item \eng{What are the symptoms of typhoid? ↘} — question ouverte en WH- : la voix descend sur \eng{TYphoid}. « Quels sont les symptômes de la typhoïde ? »
\item \eng{Have you been vaccinated? ↗} — question fermée (auxiliaire \eng{have}) : la voix monte. « As-tu été vacciné ? »
\item \eng{When does the vaccination campaign start? ↘} — question ouverte (\eng{when}) : la voix descend. « Quand commence la campagne de vaccination ? »
\item \eng{Can malaria be prevented? ↗} — question fermée (auxiliaire \eng{can}) : la voix monte. « Le paludisme peut-il être prévenu ? »
\end{enumerate}
Méthode : repérez le premier mot de la question. Si c'est un auxiliaire (is, do, have, can...), la voix monte ; si c'est un mot en WH- (what, how, when, why, where, who), la voix descend.
\end{exoresolu}

\begin{exercice}[À vous de jouer]
1) Placez l'accent principal (écrivez la syllabe accentuée en majuscules) : prevention, mosquito, epidemic, community, contaminated. 2) Prononcez à voix haute : cholera, cough, diarrhoea, hygiene, health. 3) Dites les questions suivantes avec la bonne intonation et justifiez votre choix : \eng{Do you wash your hands before eating?} — \eng{Why is clean water important?} 4) Construisez une énumération de trois conseils de prévention et marquez les flèches d'intonation.
\end{exercice}

\begin{corrige}[Exercice : Prononciation et intonation]
1) preVENtion ; mosQUIto ; epiDEmic ; comMUnity ; conTAMinated. 2) « KO-lé-ra », « KOF », « daïo-RI-a », « HAÏ-djine » (h soufflé), « HELSS » (langue entre les dents). 3) \eng{Do you wash your hands before eating? ↗} : question fermée commençant par l'auxiliaire \eng{do}, donc voix montante. \eng{Why is clean water important? ↘} : question ouverte en WH-, donc voix descendante sur \eng{imPORtant}. 4) Exemple de réponse : \eng{First, wash your hands with soap ↗ ; then, drink only boiled water ↗ ; finally, sleep under a mosquito net ↘.} La voix monte sur les deux premiers conseils et descend sur le dernier pour marquer la fin de la liste.
\end{corrige}

\begin{aretenir}[L'essentiel de la section]
\begin{itemize}
\item Chaque mot anglais a une syllabe accentuée, souvent différente du français : \eng{preVENtion}, \eng{comMUnity}, \eng{epiDEmic}.
\item Le \eng{h} anglais se prononce (\eng{hygiene}, \eng{hospital}) ; le \eng{th} se dit la langue entre les dents (\eng{health}).
\item Question fermée (yes/no) = voix qui monte ↗ ; question ouverte (WH-) = voix qui descend ↘.
\item Dans une énumération : voix montante sur chaque élément, descendante sur le dernier.
\item L'orthographe anglaise ne garantit pas la prononciation : \eng{cough} = « KOF », mais \eng{though} = « ZO ».
\end{itemize}
\end{aretenir}

Votre voix est maintenant prête : vous savez frapper les bonnes syllabes, prononcer les mots pièges et faire danser la mélodie des questions. Il ne reste plus qu'à mettre tout cela en action dans la dernière section : les jeux de rôle guidés et la production orale libre, où vous incarnerez infirmiers, élèves et membres de la communauté dans de vrais échanges en anglais.

\coursec{Jeux de rôle guidés et production orale}

Après avoir travaillé la prononciation, l'accentuation et l'intonation dans la section précédente, il est temps de mettre tout cela en pratique. Savoir dire \eng{“Could you tell me how to prevent malaria?”} avec la bonne intonation ne suffit pas : il faut maintenant être capable de mener un échange complet, du salut à la conclusion, en réagissant naturellement à ce que dit votre partenaire. C'est exactement l'objectif des jeux de rôle (\eng{role plays}) : simuler une situation réelle de communication pour s'entraîner à échanger des informations sur la prévention des maladies transmissibles, à l'école, au travail et dans la communauté.

\courssub{Qu'est-ce qu'un jeu de rôle réussi ?}

Dans un jeu de rôle, chaque élève reçoit une \cle{fiche de rôle} (\eng{role card}) qui décrit son personnage, son objectif et les informations qu'il doit obtenir ou transmettre. Contrairement à un dialogue récité, le jeu de rôle exige d'\emph{improviser} à partir d'une trame : on ne connaît pas à l'avance les mots exacts de son partenaire.

Un échange d'informations réussi suit presque toujours les mêmes six étapes, quel que soit le sujet (paludisme, choléra, grippe) :

\begin{popfigure}
\begin{tikzpicture}[
  node distance=0.55cm,
  etape/.style={rectangle, rounded corners=3pt, draw=popblue, fill=popblueL, text width=4.6cm, align=center, minimum height=0.9cm, font=\small},
  fleche/.style={-{Stealth[length=2.5mm]}, thick, poporange}
]
\node[etape, align=center] (e1){\textbf{1. Greet}\\ “Good morning, Nurse.”};
\node[etape, below=of e1, align=center] (e2){\textbf{2. Introduce the topic}\\ “I'd like to ask about malaria.”};
\node[etape, below=of e2, align=center] (e3){\textbf{3. Ask questions}\\ “How is it transmitted?”};
\node[etape, below=of e3, align=center] (e4){\textbf{4. Give structured information}\\ “First… Then… Finally…”};
\node[etape, below=of e4, align=center] (e5){\textbf{5. React and check understanding}\\ “I see. So, if I understand correctly…”};
\node[etape, below=of e5, align=center] (e6){\textbf{6. Conclude politely}\\ “Thank you very much. Goodbye.”};
\draw[fleche] (e1) -- (e2);
\draw[fleche] (e2) -- (e3);
\draw[fleche] (e3) -- (e4);
\draw[fleche] (e4) -- (e5);
\draw[fleche] (e5) -- (e6);
\end{tikzpicture}
\legende{Les six étapes d'un jeu de rôle réussi : de la salutation à la conclusion polie.}
\end{popfigure}

\begin{methode}[Réussir un jeu de rôle en quatre temps]
\begin{enumerate}
\item \textbf{Lire attentivement sa fiche de rôle} : qui suis-je ? qui est mon partenaire ? quel est mon objectif (obtenir une information, donner des conseils, convaincre) ?
\item \textbf{Préparer mentalement trois questions et trois informations} avant de commencer. Notez des mots-clés, pas des phrases entières.
\item \textbf{Pendant l'échange, utiliser au moins deux formules de relance} : \eng{“Could you explain that again?”}, \eng{“What do you mean exactly?”}, \eng{“Anything else?”}
\item \textbf{Conclure sans s'arrêter brusquement} : remerciez, résumez en une phrase, saluez. Ne dites jamais simplement \eng{“Finished!”}
\end{enumerate}
\end{methode}

\courssub{Jeu de rôle 1 (guidé) : à l'infirmerie — la prévention du paludisme}

\textbf{Situation.} L'élève A est délégué(e) de classe (\eng{class prefect}) ; l'élève B est l'infirmier(ère) scolaire (\eng{school nurse}). Le sujet : la prévention du \cle{paludisme} (\eng{malaria}) dans l'école.

\begin{textebox}[Trame du jeu de rôle 1 — At the school infirmary]
\textbf{A:} Greet the nurse and ask about malaria.\\
\textbf{B:} Answer and mention mosquito bites.\\
\textbf{A:} Ask how to prevent it at school.\\
\textbf{B:} Give three pieces of advice (nets, stagnant water, health centre).\\
\textbf{A:} React, check understanding, thank and conclude.
\end{textebox}

Voici comment cette trame peut devenir un dialogue réel. Observez d'abord les formules en gras, puis lisez l'exemple complet :

\begin{exemplebox}[Un exemple de réalisation du jeu de rôle 1]
\eng{\textbf{A:} Good morning, Nurse. I'm the class prefect of Terminale A. Could I ask you a few questions about malaria?}\\
« Bonjour, Madame l'infirmière. Je suis le délégué de la Terminale A. Puis-je vous poser quelques questions sur le paludisme ? »

\eng{\textbf{B:} Of course. Malaria is a serious disease. It is transmitted by mosquito bites, especially at night.}\\
« Bien sûr. Le paludisme est une maladie grave. Il est transmis par les piqûres de moustiques, surtout la nuit. »

\eng{\textbf{A:} I see. And how can we prevent it here at school?}\\
« Je vois. Et comment pouvons-nous le prévenir ici à l'école ? »

\eng{\textbf{B:} Well, first, students should sleep under treated mosquito nets. Then, we must clear stagnant water around the classrooms, because mosquitoes breed there. Finally, anyone with a fever should go to the health centre immediately.}\\
« Eh bien, d'abord, les élèves devraient dormir sous des moustiquaires imprégnées. Ensuite, nous devons éliminer les eaux stagnantes autour des salles de classe, car les moustiques s'y reproduisent. Enfin, toute personne ayant de la fièvre devrait se rendre au centre de santé immédiatement. »

\eng{\textbf{A:} If I understand correctly, we must clear stagnant water and use mosquito nets. Thank you very much for this information, Nurse. Goodbye!}\\
« Si je comprends bien, nous devons éliminer les eaux stagnantes et utiliser des moustiquaires. Merci beaucoup pour ces informations, Madame. Au revoir ! »
\end{exemplebox}

Remarquez les deux phrases-clés de la fin :

\begin{exemplebox}[Deux formules à réutiliser dans tous vos jeux de rôle]
\eng{If I understand correctly, we must clear stagnant water.} = « Si je comprends bien, nous devons éliminer les eaux stagnantes. » — formule de \cle{vérification de la compréhension}.

\eng{Let's work together to keep our school clean.} = « Travaillons ensemble pour garder notre école propre. » — formule de \cle{conclusion mobilisatrice}, idéale pour clore un échange sur la santé communautaire.
\end{exemplebox}

\courssub{Jeu de rôle 2 (semi-guidé) : à la radio scolaire — une épidémie de choléra}

\textbf{Situation.} L'élève A est journaliste à la radio scolaire (\eng{school radio journalist}) ; l'élève B est agent de santé communautaire (\eng{community health worker}). Le sujet : une épidémie de \cle{choléra} (\eng{cholera outbreak}) dans le quartier. Cette fois, seule la fiche de rôle est fournie : à vous de construire les questions.

\begin{experience}[Fiches de rôle — The cholera interview]
\textbf{Role card A — Journalist:} You work for your school radio. A cholera outbreak has been reported in your neighbourhood. Interview the health worker. You must find out: (1) what cholera is and how it spreads, (2) the symptoms, (3) what people should do to protect themselves, (4) where to go for treatment. Prepare at least four questions. React to the answers and conclude your interview professionally.

\textbf{Role card B — Community health worker:} A journalist interviews you about cholera. Give clear, structured information: cholera spreads through dirty water and unwashed food; symptoms include diarrhoea and vomiting; advice: boil drinking water, wash hands with soap, wash fruit and vegetables, go to the health centre quickly. Use \eng{first}, \eng{then}, \eng{finally}.
\end{experience}

\begin{exoresolu}[Construire les questions du journaliste]
Énoncé : à partir de la fiche de rôle A, formulez en anglais les quatre questions que le journaliste doit poser.
\tcblower
Solution détaillée : on transforme chaque objectif de la fiche en question ouverte ou fermée correcte (inversion du sujet et de l'auxiliaire, ou \eng{do/does} pour le présent simple) :
\begin{enumerate}
\item \eng{What is cholera, and how does it spread?} (question ouverte avec \eng{what} ; \eng{does} pour la troisième personne)
\item \eng{What are the symptoms of this disease?}
\item \eng{What should people do to protect themselves?} (\eng{should} pour demander un conseil)
\item \eng{Where should sick people go for treatment?}
\end{enumerate}
Le journaliste peut conclure ainsi : \eng{Thank you for this important information. Dear listeners, remember: boil your water and wash your hands. This was your school radio.} — « Merci pour ces informations importantes. Chers auditeurs, rappelez-vous : faites bouillir votre eau et lavez-vous les mains. C'était votre radio scolaire. »
\end{exoresolu}

\courssub{Jeu de rôle 3 (libre) : au travail à Douala — l'hygiène après un cas de grippe}

\textbf{Situation.} Deux collègues (\eng{colleagues}) travaillent dans une entreprise de Douala. Un employé est absent à cause de la \cle{grippe} (\eng{flu}). Les deux collègues échangent sur l'hygiène au travail (\eng{workplace hygiene}). Aucune trame n'est fournie : c'est la production libre.

\begin{experience}[Free role play — At the office in Douala]
En binôme, préparez en cinq minutes un dialogue d'environ dix répliques. Contraintes obligatoires :
\begin{itemize}
\item saluer et introduire le sujet (\eng{“Have you heard? Peter is sick with the flu.”}) ;
\item poser au moins trois questions ;
\item donner des informations structurées avec \eng{first}, \eng{then}, \eng{finally} (se laver les mains, aérer le bureau, rester chez soi en cas de fièvre, se faire vacciner) ;
\item utiliser au moins deux réactions naturelles (\eng{“Really?”}, \eng{“I see.”}, \eng{“That's a good point.”}) ;
\item conclure poliment, par exemple : \eng{Let's work together to keep our office healthy.}
\end{itemize}
Puis jouez la scène devant la classe, sans lire vos notes.
\end{experience}

\begin{attention}[Ne récite pas ton dialogue par cœur !]
L'erreur classique des élèves francophones est d'apprendre le dialogue mot à mot et de le débiter sans regarder le partenaire. À l'oral, l'examinateur valorise au contraire :
\begin{itemize}
\item la \cle{spontanéité} : adaptez-vous à ce que dit réellement votre partenaire, même s'il sort de la trame ;
\item le \cle{contact visuel} (\eng{eye contact}) : regardez votre interlocuteur, pas votre feuille ;
\item les \cle{réactions naturelles} : \eng{“I see.”} « Je vois. », \eng{“Really?”} « Vraiment ? », \eng{“Oh, I didn't know that!”} « Oh, je ne savais pas ! » ;
\item la capacité à \cle{relancer} quand il y a un silence : \eng{“Could you tell me more about that?”}
\end{itemize}
Autre piège : ne traduisez pas « se laver les mains » par un possessif mal accordé — dites \eng{wash your hands} (en anglais, le possessif s'accorde avec le sujet : \eng{I wash my hands}, \eng{you wash your hands}, \eng{they wash their hands}), contrairement au français « se laver \emph{les} mains » où l'article défini suffit.
\end{attention}

\courssub{Grille d'auto-évaluation et feedback}

Après chaque jeu de rôle, évaluez-vous honnêtement avec la grille suivante, puis échangez vos remarques en binôme.

\begin{center}
\begin{tabularx}{\linewidth}{|X|l|l|}
\hline
\rowcolor{popblueL} \textbf{Critère} & \textbf{Oui} & \textbf{À améliorer} \\
\hline
Did I greet my partner? (Ai-je salué ?) & & \\
\hline
Did I introduce the topic clearly? (Ai-je introduit le sujet ?) & & \\
\hline
Did I ask at least three questions? (Ai-je posé au moins 3 questions ?) & & \\
\hline
Did I give structured information (first, then, finally)? & & \\
\hline
Did I react and check understanding? (Ai-je réagi et vérifié ma compréhension ?) & & \\
\hline
Did I use at least two follow-up phrases? (Ai-je relancé ?) & & \\
\hline
Did I conclude politely? (Ai-je conclu poliment ?) & & \\
\hline
\end{tabularx}
\end{center}

Le \cle{feedback} se fait ensuite à deux niveaux :
\begin{itemize}
\item \textbf{Les pairs} (\eng{peer feedback}) : chaque binôme qui écoute note un point fort (\eng{strength}) et un axe d'amélioration (\eng{area for improvement}), par exemple : \eng{“Your questions were very clear, but try to look at your partner more.”} « Tes questions étaient très claires, mais essaie de regarder davantage ton partenaire. »
\item \textbf{Le professeur} : il corrige les erreurs de langue les plus fréquentes (questions sans inversion, oubli du \eng{-s} à la troisième personne, prononciation de \eng{health} « hèlth » et non « hilt ») et félicite les réussites.
\end{itemize}

\begin{exercice}[À vous de jouer]
Reprenez le jeu de rôle 2 (choléra) avec un nouveau partenaire, mais inversez les rôles. Cette fois, le journaliste doit poser une question inattendue, par exemple : \eng{“Is it safe to buy food from street vendors during the outbreak?”} Remplissez ensuite la grille d'auto-évaluation.
\end{exercice}

\begin{corrige}[Exercice — indications méthodologiques]
Il n'existe pas de réponse unique, mais un bon échange contient : un salut (\eng{Good afternoon, thank you for joining us}), au moins trois questions correctement formées (inversion ou auxiliaire \eng{do/does/is}), des réponses structurées avec \eng{first / then / finally}, au moins deux réactions (\eng{I see} ; \eng{That's worrying!}), une vérification de compréhension (\eng{So, if I understand correctly, we should boil all drinking water?}) et une conclusion (\eng{Thank you for your time. Stay safe, everyone!}). Vérifiez chaque ligne de la grille : si une case est « à améliorer », rejouez la scène en vous concentrant sur ce point.
\end{corrige}

\courssub{Complément Bac : briller à l'épreuve orale}

\begin{savaistu}[L'épreuve d'expression orale au Baccalauréat A]
Au Cameroun, à l'épreuve orale du Baccalauréat A, l'interaction avec l'examinateur ou avec un pair est évaluée sur trois critères principaux : la \cle{pertinence des questions} posées, la \cle{correction de la langue} (grammaire, vocabulaire, prononciation) et l'\cle{aisance} (\eng{fluency}) — c'est-à-dire la capacité à parler sans longs silences et à réagir naturellement. Un élève qui écoute vraiment son partenaire et rebondit sur ses propos obtient toujours une meilleure note qu'un élève qui récite parfaitement un texte appris par cœur.
\end{savaistu}

\begin{aretenir}[Les trois atouts d'un candidat excellent]
À l'épreuve orale, l'examinateur apprécie particulièrement :
\begin{enumerate}
\item les \cle{reformulations} : \eng{“So, if I understand correctly, you mean that…”} — elles prouvent que vous écoutez et comprenez réellement ;
\item les \cle{connecteurs} : \eng{first}, \eng{then}, \eng{finally}, mais aussi \eng{moreover} « de plus », \eng{however} « cependant » — ils structurent votre propos ;
\item une \cle{conclusion claire} : jamais d'arrêt brutal ; terminez par un résumé et un remerciement, par exemple \eng{“To sum up, prevention is better than cure. Thank you.”} « Pour résumer, mieux vaut prévenir que guérir. Merci. »
\end{enumerate}
\end{aretenir}

En maîtrisant ces trois jeux de rôle — guidé, semi-guidé puis libre — vous avez progressé de la simple répétition vers la communication authentique. Vous êtes désormais capable d'échanger des informations sur la prévention des maladies transmissibles dans n'importe quel contexte : l'école, le quartier ou l'entreprise. Gardez la grille d'auto-évaluation à portée de main : elle vous servira pour toutes les futures leçons d'expression orale de l'année.

\newpage

\coursec{Activité d'intégration}

\courssub{Objectifs de l'activité}

À la fin de cette activité d'intégration, l'élève sera capable de :

\begin{itemize}
\item mobiliser dans une tâche authentique l'ensemble des formules fonctionnelles étudiées dans la leçon : demander une information (\eng{Could you tell me...?}), donner une information (\eng{The main way to prevent it is...}), relancer (\eng{What about...?}, \eng{Could you give us an example?}), réagir (\eng{Really? I didn't know that!}) et conclure (\eng{Thank you very much for your time.}) ;
\item mener un échange oral de deux minutes en anglais, de façon fluide et naturelle, sur la prévention des maladies transmissibles à l'école, au travail et dans la communauté ;
\item employer correctement le lexique de la santé et de l'hygiène (\eng{prevention, symptoms, transmission, hygiene, vaccination, safe water, mosquito net}...) ;
\item soigner sa prononciation, son accentuation et son intonation (intonation montante pour les questions fermées, descendante pour les questions ouvertes) ;
\item écouter activement les productions des camarades et relever des informations précises dans une fiche d'écoute (compétence \cle{listening}).
\end{itemize}

\courssub{Situation}

\begin{textebox}[School Health Radio — the context]
Your school, Government High School Biyem-Assi, Yaoundé, has just launched its own radio programme: \textbf{“School Health Radio”}. Every Friday, students broadcast a short programme in English about health and well-being. This week, the Ministry of Public Health has asked schools to raise awareness about \textbf{communicable diseases} after cases of cholera were reported in some districts of the Centre region.

The radio team has been selected: you! In groups of three, you will prepare and perform a live two-minute interview. One student is the \textbf{journalist}, one is a \textbf{health expert} (school nurse, community health worker or doctor), and one is a \textbf{citizen} (a student, a shopkeeper or an employee) who asks an extra question. Your interview must inform the listeners about one disease and how to prevent it in one specific place. The best interview will be recorded and broadcast during the school assembly on Monday morning!
\end{textebox}

\begin{definition}[Rappel de vocabulaire : communicable diseases]
\eng{communicable disease} (nom) : maladie transmissible (qui passe d'une personne à une autre). \eng{outbreak} : épidémie, flambée. \eng{awareness campaign} : campagne de sensibilisation. \eng{to raise awareness} : sensibiliser. \eng{health worker} : agent de santé. \eng{listener} : auditeur. \eng{to broadcast} : diffuser (à la radio) — verbe irrégulier : \eng{broadcast – broadcast – broadcast}.
\end{definition}

\courssub{Consignes}

\begin{enumerate}
\item \textbf{Tirage au sort.} Chaque groupe tire au sort une maladie (\eng{malaria}, \eng{cholera}, \eng{typhoid fever} ou \eng{COVID-19}) et un lieu (\eng{school}, \eng{workplace}, \eng{market} ou \eng{community}).

\item \textbf{Préparation (10 minutes).} Dans votre groupe, rédigez la trame de votre interview en suivant la structure fonctionnelle : salutation et présentation de l'invité $\rightarrow$ questions du journaliste $\rightarrow$ informations structurées de l'expert (causes, symptômes, au moins trois mesures de prévention adaptées au lieu) $\rightarrow$ relances (\eng{Could you explain why...?}, \eng{What should people do if...?}) $\rightarrow$ question supplémentaire du citoyen $\rightarrow$ conclusion et remerciements. Répartissez les rôles et répétez.

\item \textbf{Vérification linguistique.} Relisez votre script : questions bien formées (inversion ou \eng{do/does/did}), vocabulaire de la santé correct, formules de politesse (\eng{Could you...?}, \eng{Thank you for...}).

\item \textbf{Performance (2 minutes par groupe).} Jouez votre interview devant la classe, sans lire intégralement votre texte. Soignez la prononciation et l'intonation.

\item \textbf{Écoute active.} Pendant les interviews des autres groupes, remplissez votre fiche d'écoute : nom de la maladie, une cause, trois mesures de prévention citées.

\item \textbf{Vote final.} Votez pour la meilleure interview selon la grille : formules fonctionnelles utilisées, prononciation, intonation, naturel de l'échange.
\end{enumerate}

\begin{experience}[La fiche d'écoute — listening grid]
Pendant chaque interview, complétez la grille suivante :

\begin{center}
\begin{tabularx}{\linewidth}{|X|X|X|X|}
\hline
\rowcolor{popblueL} Disease & One cause & Prevention 1 & Prevention 2 \& 3 \\
\hline
 & & & \\
\hline
 & & & \\
\hline
 & & & \\
\hline
\end{tabularx}
\end{center}
\end{experience}

\courssub{Ressources à mobiliser}

\begin{aretenir}[La trame fonctionnelle de l'interview]
\begin{itemize}
\item \textbf{Ouvrir :} \eng{Good morning, dear listeners! Welcome to School Health Radio. Today, our guest is Dr Ngo Bassa, a community health worker.}
\item \textbf{Demander :} \eng{Could you tell us what cholera is?} / \eng{How is it transmitted?} / \eng{What are the main symptoms?}
\item \textbf{Donner une information :} \eng{Cholera is caused by contaminated water and food.} / \eng{The best way to prevent it is to drink safe, boiled or treated water.}
\item \textbf{Relancer :} \eng{Could you give us a concrete example?} / \eng{And what about prevention at school?}
\item \textbf{Réagir :} \eng{Really? That's very interesting!} / \eng{I see, thank you.}
\item \textbf{Conclure :} \eng{Thank you very much for your time and your advice.} / \eng{That was Dr Ngo Bassa on School Health Radio. Stay healthy!}
\end{itemize}
\end{aretenir}

\begin{propriete}[Rappel : former les questions]
\begin{itemize}
\item Question ouverte avec mot interrogatif : \textbf{Wh-word + auxiliaire + sujet + verbe} : \eng{How can we prevent malaria?} (« Comment pouvons-nous prévenir le paludisme ? ») — intonation descendante.
\item Question fermée (oui/non) : \textbf{auxiliaire + sujet + verbe} : \eng{Is cholera dangerous for children?} (« Le choléra est-il dangereux pour les enfants ? ») — intonation montante.
\item Question de politesse (question indirecte) : \eng{Could you tell me + sujet + verbe} (sans inversion après \eng{tell me}) : \eng{Could you tell me what the symptoms are?} et non \emph{Could you tell me what are the symptoms?}
\end{itemize}
\end{propriete}

\begin{propriete}[Rappel : exprimer le conseil et l'obligation]
Pour donner des mesures de prévention, l'expert utilise :
\begin{itemize}
\item \eng{should + base verbale} (conseil) : \eng{Students should sleep under mosquito nets.} (« Les élèves devraient dormir sous des moustiquaires. »)
\item \eng{must + base verbale} (obligation forte) : \eng{We must drink safe water.} (« Nous devons boire de l'eau saine. »)
\item \eng{the best way to... is to...} : \eng{The best way to prevent cholera is to wash your hands with soap.} (« La meilleure façon de prévenir le choléra est de se laver les mains au savon. »)
\end{itemize}
\end{propriete}

\begin{attention}[Pièges fréquents des francophones]
\begin{itemize}
\item \eng{advice} est indénombrable : \eng{some advice}, jamais \emph{advices}.
\item \eng{information} est indénombrable : \eng{useful information}, jamais \emph{informations}.
\item Ne traduisez pas « sensibiliser » par \emph{sensibilize} : dites \eng{to raise awareness} ou \eng{to educate people}.
\item \eng{to prevent} se construit avec \eng{from + V-ing} : \eng{Washing hands prevents people from getting sick.} (« Se laver les mains empêche les gens de tomber malades. »)
\item Prononciation : \eng{disease} se prononce « di-ziz » ; \eng{health} se prononce « èl-th » (le « th » est sourd : placez la langue entre les dents et soufflez, sans faire vibrer les cordes vocales) ; \eng{cholera} s'accentue sur la première syllabe : « KO-lé-ra » ; \eng{stomach} se prononce « steu-mèk ».
\end{itemize}
\end{attention}

\begin{savaistu}[La radio scolaire, un outil de santé publique]
Dans de nombreux pays anglophones d'Afrique (Ghana, Kenya, Nigeria), comme dans les régions anglophones du Cameroun, la radio reste le média le plus écouté : les campagnes de santé publique y diffusent leurs messages en anglais simple et dans les langues locales. Les « school health clubs » et les radios scolaires jouent un rôle réel dans la prévention : pendant l'épidémie d'Ebola en Afrique de l'Ouest, puis lors de la pandémie de COVID-19, des messages radiophoniques enregistrés par des élèves ont contribué à diffuser les gestes barrières dans les communautés. Votre interview d'aujourd'hui reproduit donc une pratique professionnelle authentique !
\end{savaistu}

\courssub{Proposition de production et corrigé commenté}

\begin{exoresolu}[Modèle d'interview — malaria at school]
Tirage au sort : \eng{malaria} + \eng{school}. Rédigez et jouez l'interview complète.

\tcblower

\textbf{Journalist :} Good morning, dear listeners, and welcome to School Health Radio! Today we are talking about malaria, a very common disease in our country. Our guest is Mrs Etoundi, our school nurse. Good morning, Mrs Etoundi, thank you for coming.

\textbf{Expert :} Good morning, and thank you for inviting me.

\textbf{Journalist :} Mrs Etoundi, could you tell us what malaria is and how it is transmitted?

\textbf{Expert :} Of course. Malaria is a serious disease caused by a parasite. It is transmitted to people through the bites of infected mosquitoes, usually at night.

\textbf{Journalist :} I see. And what are the main symptoms?

\textbf{Expert :} The main symptoms are high fever, headaches, chills and vomiting. If a student feels these symptoms, he or she should go to the health centre immediately for a test.

\textbf{Journalist :} Really? That's important to know. Now, how can we prevent malaria here at school?

\textbf{Expert :} Well, there are three simple measures. First, students should sleep under insecticide-treated mosquito nets at home. Second, we must clear the bushes and remove stagnant water around the school, because mosquitoes breed there. Third, everyone should wear long sleeves in the evening.

\textbf{Journalist :} Could you give us an example of what the school has already done?

\textbf{Expert :} Yes, of course. Last month, the students cleaned the school compound and emptied all the containers with stagnant water.

\textbf{Citizen (a student) :} Excuse me, Mrs Etoundi. I have a question. What should we do if a classmate has a fever during lessons?

\textbf{Expert :} Good question! You should inform the teacher immediately, and the student should be taken to the sick bay or to the nearest health centre. Never give medicine without a test.

\textbf{Journalist :} Thank you very much, Mrs Etoundi, for your time and your useful advice. And thank you for your question, Amina. That was School Health Radio. Remember: sleep under a mosquito net and keep your environment clean. Stay healthy, everyone!

\textbf{Commentaire des choix de langue :}
\begin{itemize}
\item La trame fonctionnelle est respectée : salutation (\eng{Good morning, dear listeners}), présentation de l'invité, questions ouvertes à intonation descendante (\eng{How can we prevent malaria here at school?}), relances (\eng{Could you give us an example...?}), réactions (\eng{Really? That's important to know.}), question du citoyen, conclusion avec remerciements.
\item Questions de politesse correctement formées : \eng{could you tell us what malaria is} — pas d'inversion après \eng{tell us}.
\item Indénombrables respectés : \eng{useful advice} (et non \emph{advices}).
\item Lexique précis : \eng{parasite} (parasite), \eng{mosquito net} (moustiquaire), \eng{stagnant water} (eau stagnante), \eng{symptoms} (symptômes), \eng{health centre} (centre de santé), \eng{sick bay} (infirmerie).
\item Structures de conseil variées : \eng{should + base verbale}, \eng{we must...}, \eng{the best way to prevent it is...}.
\item Contexte camerounais authentique : infirmière scolaire, brousse autour de l'école, eau stagnante, moustiquaires imprégnées d'insecticide.
\end{itemize}
\end{exoresolu}

\courssub{Bilan de l'activité}

\begin{aretenir}[Ce que j'ai appris à faire]
Grâce à cette activité, je sais désormais :
\begin{itemize}
\item conduire une interview courte en anglais, de la salutation à la conclusion, en respectant une trame fonctionnelle claire ;
\item poser des questions ouvertes et fermées correctement formées, avec la bonne intonation ;
\item donner des informations structurées sur une maladie transmissible (cause, transmission, symptômes, prévention) en adaptant mon discours au lieu (école, lieu de travail, marché, communauté) ;
\item relancer mon interlocuteur et réagir naturellement à ses réponses ;
\item écouter activement un échange oral pour en relever les informations essentielles ;
\item éviter les erreurs typiques des francophones (\emph{advices}, \emph{informations}, inversion après \eng{Could you tell me...}).
\end{itemize}
\end{aretenir}

\begin{methode}[Pour aller plus loin : l'auto-évaluation]
Après votre performance, évaluez-vous honnêtement :
\begin{enumerate}
\item Ai-je utilisé au moins quatre formules fonctionnelles différentes (demander, donner, relancer, conclure) ?
\item Mes questions étaient-elles grammaticalement correctes ?
\item Mon intonation distinguait-elle les questions ouvertes (descendante) et fermées (montante) ?
\item Ai-je parlé de façon naturelle, sans lire mon texte ?
\end{enumerate}
Si vous répondez « oui » à trois questions au moins, votre objectif est atteint. Sinon, répétez l'interview à la maison avec un camarade ou en vous enregistrant : c'est en parlant qu'on apprend à parler !
\end{methode}

\newpage

\coursec{Méthodes et exercices résolus}

\coursec{Lesson 28 — Speaking: Exchanges information about preventing communicable diseases in schools, workplaces, and the community}

\begin{textebox}[Warm-up : Contexte et lexique de la santé — Health Vocabulary]
\eng{Communicable diseases such as cholera, malaria, and typhoid fever remain major public health challenges in our schools, workplaces, and communities. Prevention relies on simple but effective actions: washing hands with soap, drinking safe water, sleeping under treated mosquito nets, keeping the environment clean, and getting vaccinated.}
« Les maladies transmissibles telles que le choléra, le paludisme et la fièvre typhoïde restent des défis majeurs de santé publique dans nos écoles, lieux de travail et communautés. La prévention repose sur des actions simples mais efficaces : se laver les mains au savon, boire de l'eau potable, dormir sous des moustiquaires imprégnées, maintenir l'environnement propre et se faire vacciner. »

\begin{center}
\begin{tabularx}{\linewidth}{|X|X|X|}
\hline
\rowcolor{popblueL} \textbf{English} & \textbf{Français} & \textbf{Classe / Note} \\
\hline
\eng{communicable disease} & maladie transmissible & nom composé \\
\eng{cholera} & choléra & nom \\
\eng{malaria} & paludisme & nom \\
\eng{typhoid fever} & fièvre typhoïde & nom composé \\
\eng{public health} & santé publique & nom composé \\
\eng{prevention} & prévention & nom \\
\eng{to wash hands with soap} & se laver les mains au savon & verbe + compléments \\
\eng{safe water / treated water} & eau potable / eau traitée & nom + adjectif \\
\eng{treated mosquito net} & moustiquaire imprégnée & nom composé \\
\eng{to keep the environment clean} & maintenir l'environnement propre & verbe + objet + adjectif \\
\eng{to get vaccinated (against)} & se faire vacciner (contre) & verbe + préposition \\
\eng{symptoms} & symptômes & nom (pluriel) \\
\eng{to boil water} & faire bouillir l'eau & verbe + objet \\
\eng{to cover food} & couvrir la nourriture & verbe + objet \\
\eng{health worker / nurse} & agent de santé / infirmier(ère) & nom \\
\hline
\end{tabularx}
\end{center}

\textbf{Glossaire complémentaire :} \eng{diarrhoea} (diarrhée), \eng{dehydration} (déshydratation), \eng{outbreak} (épidémie, flambée), \eng{sanitation} (assainissement), \eng{contaminated} (contaminé), \eng{to spread} (se propager), \eng{immune system} (système immunitaire).
\end{textebox}

\begin{textebox}[Support de lecture — Health Advisory for Schools and Communities]
\eng{MINISTRY OF PUBLIC HEALTH — REGIONAL DELEGATION, CENTRE REGION}

\eng{Preventing Communicable Diseases: A Shared Responsibility}

\eng{YAOUNDÉ, October 2024 — With the rainy season approaching, the risk of cholera and malaria outbreaks increases in schools and markets. The Ministry reminds everyone of the essential barrier gestures:}

\eng{1. Wash your hands with clean water and soap before eating, after using the toilet, and after handling waste.}
\eng{2. Drink only boiled or treated water. Avoid buying drinks in unclean containers.}
\eng{3. Sleep under a long-lasting insecticidal net (LLIN) every night.}
\eng{4. Keep your surroundings clean: dispose of refuse properly, clear stagnant water, and cut tall grass.}
\eng{5. Cover food to protect it from flies and dust.}
\eng{6. Ensure children complete their vaccination schedule (measles, yellow fever, polio, etc.).}
\eng{7. If a student or colleague shows symptoms (high fever, severe diarrhoea, vomiting), do not self-medicate. Take them to the nearest health centre immediately.}

\eng{“Prevention is better than cure. Let us work together for a healthier Cameroon.”}
« MINISTÈRE DE LA SANTÉ PUBLIQUE — DÉLÉGATION RÉGIONALE, RÉGION DU CENTRE. Prévention des maladies transmissibles : une responsabilité partagée. Yaoundé, octobre 2024 — Avec l'approche de la saison des pluies, le risque d'épidémies de choléra et de paludisme augmente dans les écoles et les marchés. Le Ministère rappelle à tous les gestes barrières essentiels : 1. Lavez-vous les mains à l'eau propre et au savon avant de manger, après les toilettes et après avoir manipulé des déchets. 2. Ne buvez que de l'eau bouillie ou traitée. Évitez d'acheter des boissons dans des contenants sales. 3. Dormez sous une moustiquaire imprégnée à longue durée d'action (MILD) chaque nuit. 4. Maintenez votre environnement propre : jetez les ordures correctement, éliminez l'eau stagnante, coupez l'herbe haute. 5. Couvrez la nourriture pour la protéger des mouches et de la poussière. 6. Assurez-vous que les enfants terminent leur calendrier de vaccination (rougeole, fièvre jaune, polio, etc.). 7. Si un élève ou un collègue présente des symptômes (forte fièvre, diarrhée sévère, vomissements), ne vous automédicamentez pas. Conduisez-les au centre de santé le plus proche immédiatement. "Mieux vaut prévenir que guérir. Travaillons ensemble pour un Cameroun en meilleure santé." »
\end{textebox}

\begin{propriete}[Points de grammaire clés : Conseils, obligation, questions indirectes, voix passive]
\begin{enumerate}
\item \textbf{Modaux de conseil et d'obligation} : \cle{should} (conseil), \cle{must} (obligation forte, interne), \cle{have to} (obligation externe, règle).
\begin{itemize}
\item \eng{You should wash your hands regularly.} (Conseil) — « Tu devrais / Il faut que tu te laves les mains régulièrement. »
\item \eng{Students must wear clean uniforms.} (Règle scolaire) — « Les élèves doivent porter des uniformes propres. »
\item \eng{Workers have to follow safety protocols.} (Règle de l'entreprise) — « Les travailleurs sont tenus de suivre les protocoles de sécurité. »
\end{itemize}
\emph{Forme négative :} \eng{should not / shouldn't} (déconseil), \eng{must not / mustn't} (interdiction), \eng{don't have to} (absence d'obligation = « pas besoin de »).

\item \textbf{Questions indirectes (politesse)} : Après \eng{Could you tell me...}, \eng{I would like to know...}, \eng{Do you know...}, l'ordre est \textbf{sujet + verbe} (pas d'inversion, pas d'auxiliaire \eng{do/does/did}).
\begin{itemize}
\item Direct : \eng{How does cholera spread?} → Indirect : \eng{Could you tell me how cholera spreads?}
\item Direct : \eng{When do you wash your hands?} → Indirect : \eng{I would like to know when you wash your hands.}
\end{itemize}

\item \textbf{Voix passive (présent simple) :} \eng{be (is/are) + participe passé}. Très fréquent pour les règles sanitaires.
\begin{itemize}
\item \eng{Water is boiled before drinking.} (L'eau est bouillie avant d'être bue.)
\item \eng{Children are vaccinated against measles at school.} (Les enfants sont vaccinés contre la rougeole à l'école.)
\item \eng{Waste must be disposed of properly.} (Les déchets doivent être éliminés correctement.) — Passif modal.
\end{itemize}

\item \textbf{Prépositions verbes / adjectifs du thème} :
\begin{center}
\begin{tabular}{|l|l|l|}
\hline
\rowcolor{popblueL} \textbf{Anglais} & \textbf{Français} & \textbf{Exemple} \\
\hline
\eng{vaccinated against} & vacciné contre & \eng{He is vaccinated against yellow fever.} \\
\eng{protect ... from} & protéger ... de & \eng{Protect food from flies.} \\
\eng{suffer from} & souffrir de & \eng{Many people suffer from malaria.} \\
\eng{listen to} & écouter & \eng{Listen to the nurse's advice.} \\
\eng{responsible for} & responsable de & \eng{We are responsible for our health.} \\
\hline
\end{tabular}
\end{center}
\end{enumerate}
\end{propriete}

\courssub{Fiche méthode 1 — Échanger des informations sur la prévention des maladies}

\begin{methode}[Échanger des informations dans la situation proposée]
Pour mener un échange d'informations sur la prévention des maladies transmissibles (à l'école, au travail, dans la communauté), suis ces étapes :

\begin{enumerate}
\item \textbf{Identifier le cadre et les rôles} : qui parle (infirmier/élève \eng{nurse/student}, agent de santé/habitant \eng{health worker/resident}, employeur/employé \eng{employer/employee}) et où (\eng{school infirmary}, \eng{health centre}, \eng{workplace}, \eng{market}). Le choix du vocabulaire dépend du lieu : \eng{wash your hands with soap} « lavez-vous les mains au savon », \eng{sleep under a treated mosquito net} « dormez sous une moustiquaire imprégnée », \eng{boil drinking water} « faites bouillir l'eau de boisson », \eng{get vaccinated against...} « faites-vous vacciner contre... ».
\item \textbf{Ouvrir l'échange poliment} : \eng{Good morning, nurse. May I ask you a question?} « Bonjour infirmière. Puis-je vous poser une question ? » ou \eng{Excuse me, could you give me some information about...?} « Excusez-moi, pourriez-vous me donner des informations sur... ? »
\item \textbf{Poser des questions précises} avec les bons mots interrogatifs : \eng{How can we prevent cholera in our school?} « Comment pouvons-nous prévenir le choléra dans notre école ? » ; \eng{What should we do if a classmate has a fever?} « Que devons-nous faire si un camarade a de la fièvre ? » ; \eng{When should we wash our hands?} « Quand devons-nous nous laver les mains ? »
\item \textbf{Donner l'information avec des structures de conseil} : \eng{You should...} « Tu devrais / Il faut que tu... » ; \eng{We must...} « Nous devons... » ; \eng{It is important to...} « Il est important de... » ; \eng{The best way to prevent malaria is to...} « La meilleure façon de prévenir le paludisme est de... »
\item \textbf{Relancer et réagir} pour montrer qu'on écoute (backchanneling) : \eng{I see.} « Je vois / Je comprends. » ; \eng{Really?} « Vraiment ? » ; \eng{Could you explain that again, please?} « Pourriez-vous expliquer cela à nouveau, s'il vous plaît ? » ; \eng{And what about...?} « Et qu'en est-il de... ? »
\item \textbf{Conclure et remercier} : \eng{Thank you very much for this information. I will share it with my classmates.} « Merci beaucoup pour ces informations. Je les partagerai avec mes camarades. »
\end{enumerate}

\textbf{Réflexe} : un bon échange = question + réponse + réaction + nouvelle question. Ne récite pas un monologue : dialogue !
\end{methode}

\begin{exoresolu}[Échange d'informations — type examen]
\textbf{Énoncé} — Complete the following dialogue between a student and the school nurse with appropriate questions or answers.

\eng{Student: Good morning, nurse. (1) ........................ about preventing malaria in our school?}

\eng{Nurse: Of course. The best way is to sleep under a treated mosquito net.}

\eng{Student: I see. (2) ........................ keep our classrooms clean?}

\eng{Nurse: You should sweep them every morning and empty the dustbins.}

\eng{Student: (3) ........................ . Thank you very much, nurse.}
\tcblower
\textbf{Solution détaillée}

\textbf{(1)} L'élève ouvre l'échange et demande la permission de poser une question. La réponse de l'infirmière (\eng{Of course}) montre qu'on attend une demande polie. Réponse : \eng{May I ask you a question} ou \eng{Could you give me some information}. La structure \eng{May I + base verbale} exprime la demande de permission ; \eng{Could you + base verbale} exprime la demande polie.

\textbf{(2)} La réponse \eng{You should sweep them every morning...} est un conseil sur la manière de faire : on attend donc une question en \eng{How}. Réponse : \eng{How should we} ou \eng{How can we}. Le mot interrogatif \eng{how} porte sur le moyen, ce qui correspond exactement au contenu de la réponse.

\textbf{(3)} L'élève réagit à l'information avant de remercier : c'est la phase de clôture. Réponse : \eng{That is very useful information} ou \eng{I understand now}. Une simple réaction positive (\eng{I see}, \eng{Thank you}) convient aussi.

\textbf{Dialogue complet rédigé :}

\eng{Student: Good morning, nurse. May I ask you a question about preventing malaria in our school?}

\eng{Nurse: Of course. The best way is to sleep under a treated mosquito net.}

\eng{Student: I see. How should we keep our classrooms clean?}

\eng{Nurse: You should sweep them every morning and empty the dustbins.}

\eng{Student: That is very useful information. Thank you very much, nurse.}

\textbf{Conclusion} : toujours faire correspondre le mot interrogatif (\eng{how}, \eng{what}, \eng{when}, \eng{where}) au contenu de la réponse donnée.
\end{exoresolu}

\courssub{Fiche méthode 2 — Poser des questions et répondre de façon adaptée}

\begin{methode}[Poser des questions et répondre de façon adaptée]
\begin{enumerate}
\item \textbf{Choisir le bon type de question} :
\begin{itemize}
\item Question fermée (réponse yes/no) : auxiliaire + sujet + verbe : \eng{Do you wash your hands before eating?} « Te laves-tu les mains avant de manger ? » ; \eng{Is cholera a communicable disease?} « Le choléra est-il une maladie transmissible ? »
\item Question ouverte (information) : mot interrogatif + auxiliaire + sujet + verbe : \eng{What are the symptoms of typhoid fever?} « Quels sont les symptômes de la fièvre typhoïde ? » ; \eng{How does malaria spread?} « Comment le paludisme se propage-t-il ? »
\end{itemize}
\item \textbf{Respecter l'inversion sujet–auxiliaire} : en anglais, on ne pose jamais une question comme en français avec l'intonation seule. On dit \eng{Do you know...?} et jamais \eng{You know...?} avec une simple intonation montante.
\item \textbf{Utiliser le bon auxiliaire} : \eng{do/does} au présent simple, \eng{did} au prétérit, \eng{is/are} avec \eng{be} ou le présent continu : \eng{Does a mosquito net protect us?} « Une moustiquaire nous protège-t-elle ? » ; \eng{Did the health worker visit your quarter?} « L'agent de santé a-t-il visité votre quartier ? »
\item \textbf{Répondre de façon complète et adaptée} : à une question fermée, réponds \eng{Yes, we do. / No, it isn't.} puis ajoute une information ; à une question ouverte, donne une phrase complète : \eng{Malaria spreads through the bites of infected mosquitoes.} « Le paludisme se propage par les piqûres de moustiques infectés. »
\item \textbf{Adapter le registre} : avec une autorité (nurse, doctor, health officer), reste poli : \eng{Could you tell me...?}, \eng{I would like to know...}
\end{enumerate}
\end{methode}

\begin{exoresolu}[Questions et réponses — type examen]
\textbf{Énoncé} — A. Put the words in the correct order to make questions. B. Answer each question with a complete sentence about disease prevention.

1. \eng{cholera / how / spread / does / ?}

2. \eng{wash / you / do / hands / your / when / ?}

3. \eng{vaccinated / children / are / against measles / your school / in / ?}
\tcblower
\textbf{Solution détaillée}

\textbf{A. Ordre des mots}

\textbf{1.} Mot interrogatif (\eng{how}) + auxiliaire (\eng{does}, car le sujet \eng{cholera} est à la 3\textsuperscript{e} personne du singulier) + sujet + base verbale : \eng{How does cholera spread?} Attention : avec \eng{does}, le verbe redevient \eng{spread} sans -s.

\textbf{2.} \eng{When do you wash your hands?} Le mot interrogatif \eng{when} ouvre la phrase, suivi de l'auxiliaire \eng{do} (sujet \eng{you}), puis du sujet et du verbe. L'ordre des compléments : \eng{your hands} après le verbe.

\textbf{3.} Question fermée avec le verbe \eng{be} au passif : \eng{Are children vaccinated against measles in your school?} L'auxiliaire \eng{are} passe devant le sujet \eng{children} ; le participe passé \eng{vaccinated} reste après le sujet ; la préposition correcte est \eng{against} (et non \eng{for} ou \eng{from}).

\textbf{B. Réponses rédigées}

1. \eng{Cholera spreads through contaminated water and food, so we must drink boiled or treated water.} « Le choléra se propage par l'eau et les aliments contaminés, donc nous devons boire de l'eau bouillie ou traitée. »

2. \eng{I wash my hands with soap before eating, after using the toilet and after playing.} « Je me lave les mains au savon avant de manger, après avoir utilisé les toilettes et après avoir joué. »

3. \eng{Yes, they are. All the children in our school are vaccinated against measles during the vaccination campaign.} « Oui, c'est le cas. Tous les enfants de notre école sont vaccinés contre la rougeole pendant la campagne de vaccination. »

\textbf{Conclusion} : une question anglaise correcte = mot interrogatif (si besoin) + auxiliaire + sujet + verbe. Une bonne réponse reprend l'auxiliaire de la question (\eng{Yes, they are}) puis développe l'information.
\end{exoresolu}

\courssub{Fiche méthode 3 — Utiliser les formules fonctionnelles de l'échange}

\begin{methode}[Les formules fonctionnelles : demander, donner, relancer, conclure]
Mémorise un répertoire de formules prêtes à l'emploi, classées par fonction (traduction entre guillemets) :

\begin{enumerate}
\item \textbf{Demander une information} : \eng{Could you tell me how to prevent typhoid?} « Pourriez-vous me dire comment prévenir la fièvre typhoïde ? » ; \eng{I would like to know what the symptoms of cholera are.} « Je voudrais savoir quels sont les symptômes du choléra. » ; \eng{What should we do to keep our workplace clean?} « Que devrions-nous faire pour garder notre lieu de travail propre ? »
\item \textbf{Donner une information / un conseil} : \eng{You should always wash your hands with soap.} « Vous devriez toujours vous laver les mains au savon. » ; \eng{The main cause of cholera is dirty water.} « La cause principale du choléra est l'eau sale. » ; \eng{It is important to cover food to protect it from flies.} « Il est important de couvrir la nourriture pour la protéger des mouches. »
\item \textbf{Relancer l'échange} : \eng{And what about vaccination?} « Et qu'en est-il de la vaccination ? » ; \eng{Could you give me an example?} « Pourriez-vous me donner un exemple ? » ; \eng{What do you mean by “communicable”?} « Que voulez-vous dire par "transmissible" ? »
\item \textbf{Montrer qu'on comprend (réagir)} : \eng{I see.} « Je vois / Je comprends. » ; \eng{That is interesting.} « C'est intéressant. » ; \eng{You are right.} « Vous avez raison. » ; \eng{Exactly.} « Exactement. »
\item \textbf{Demander de répéter ou de préciser} : \eng{Sorry, could you repeat that, please?} « Désolé, pourriez-vous répéter, s'il vous plaît ? » ; \eng{Could you speak more slowly, please?} « Pourriez-vous parler plus lentement, s'il vous plaît ? »
\item \textbf{Conclure} : \eng{Thank you for your advice.} « Merci pour vos conseils. » ; \eng{I will tell my family about it.} « Je le dirai à ma famille. » ; \eng{It was very helpful, thank you.} « C'était très utile, merci. »
\end{enumerate}

\textbf{Réflexe} : dans \eng{Could you tell me how to prevent typhoid?}, il n'y a \textbf{pas d'inversion} après \eng{tell me} : c'est une question indirecte à l'infinitif (\eng{how to prevent}). Avec un sujet, l'ordre est sujet + verbe : \eng{Could you tell me how we can prevent typhoid?}
\end{methode}

\begin{exoresolu}[Formules fonctionnelles — type examen]
\textbf{Énoncé} — Match each sentence (1–5) with its function (a–e), then rewrite sentence 4 as an indirect question beginning with \eng{Could you tell me...}.

Functions : a) asking for information — b) giving information — c) reacting — d) concluding — e) asking for repetition.

1. \eng{Sorry, could you repeat that, please?}

2. \eng{The best way to prevent malaria is to sleep under a mosquito net.}

3. \eng{I see. That is very interesting.}

4. \eng{How can we prevent cholera at school?}

5. \eng{Thank you very much for your advice, Madam.}
\tcblower
\textbf{Solution détaillée}

\textbf{Appariement :}

1 → \textbf{e} : \eng{Could you repeat that, please?} sert à demander de répéter une information mal entendue.

2 → \textbf{b} : la phrase donne un conseil de prévention (\eng{The best way to prevent malaria is to...}) : c'est donner une information.

3 → \textbf{c} : \eng{I see. That is very interesting.} exprime une réaction qui montre qu'on écoute.

4 → \textbf{a} : question ouverte en \eng{how} qui demande une information sur les moyens de prévention.

5 → \textbf{d} : le remerciement final clôt l'échange : c'est la conclusion.

\textbf{Transformation en question indirecte :}

Question directe : \eng{How can we prevent cholera at school?}

Question indirecte : \eng{Could you tell me how we can prevent cholera at school?}

\textbf{Justification grammaticale} : dans une question indirecte introduite par \eng{Could you tell me...}, on retrouve l'ordre sujet + verbe (comme dans une phrase affirmative) : \eng{how we can prevent}, et non \eng{how can we prevent}. Il n'y a qu'une seule inversion, celle de l'introduction (\eng{Could you}). C'est une erreur classique de garder l'inversion dans la seconde partie.

\textbf{Conclusion} : chaque formule a une fonction précise dans l'échange ; repérer la fonction aide à choisir la bonne formule à l'oral comme à l'écrit.
\end{exoresolu}

\courssub{Fiche méthode 4 — Mener un court dialogue (jeu de rôle guidé)}

\begin{methode}[Mener un court dialogue en anglais]
\begin{enumerate}
\item \textbf{Lire la consigne et répartir les rôles} : par exemple, Student A = \eng{community health worker}, Student B = \eng{shopkeeper at Mokolo market}. Note qui tu es, où tu es, ce que tu dois demander ou expliquer.
\item \textbf{Préparer 3 à 4 questions et 3 à 4 réponses} avec le lexique du thème : \cle{prevent}, \cle{symptoms}, \cle{wash hands}, \cle{boil water}, \cle{vaccination}, \cle{mosquito net}, \cle{keep the environment clean}.
\item \textbf{Structurer le dialogue en 4 temps} : salutations et ouverture → questions/réponses (le cœur de l'échange) → relances et réactions → remerciements et conclusion.
\item \textbf{Soigner la prononciation et l'intonation} : voix montante ↗ pour les questions fermées (\eng{Do you boil your drinking water? ↗}), voix descendante ↘ pour les questions en \eng{wh-} (\eng{How can we prevent cholera? ↘}) et pour les affirmations.
\item \textbf{Improviser sans paniquer} : si tu ne comprends pas, utilise \eng{Sorry, could you repeat that, please?} ; si tu cherches tes mots, utilise \eng{Well...}, \eng{Let me think...}
\item \textbf{Viser 8 à 10 répliques} au minimum pour un dialogue complet et équilibré : chaque locuteur parle autant que l'autre.
\end{enumerate}
\end{methode}

\begin{exoresolu}[Jeu de rôle — production d'un dialogue type examen]
\textbf{Énoncé} — Role-play: You are a health worker visiting a secondary school in Yaoundé. A student asks you how to prevent communicable diseases at school, at work and in the community. Write a dialogue of at least ten lines.
\tcblower
\textbf{Solution détaillée — démarche}

\textbf{Étape 1 : rôles et cadre.} Student A = \eng{health worker} ; Student B = \eng{student}. Lieu : \eng{a school in Yaoundé}. Thème : \eng{prevention of communicable diseases} (cholera, malaria, typhoid).

\textbf{Étape 2 : plan en 4 temps.} Salutations → questions de l'élève et conseils de l'agent → relance sur la communauté et le travail → conclusion et remerciements.

\textbf{Étape 3 : dialogue rédigé.}

\eng{Student: Good morning, Madam. Welcome to our school. May I ask you a few questions?}

\eng{Health worker: Good morning. Of course, go ahead.}

\eng{Student: How can we prevent cholera in our school?}

\eng{Health worker: You should drink only boiled or treated water and wash your hands with soap before eating.}

\eng{Student: I see. And what about malaria?}

\eng{Health worker: The best protection is to sleep under a treated mosquito net and to cut the grass around the classrooms.}

\eng{Student: Could you tell me what we should do in our community?}

\eng{Health worker: You should keep your quarter clean, cover food to protect it from flies and take sick people to the health centre quickly.}

\eng{Student: That is very useful. Should workers also protect themselves at their workplaces?}

\eng{Health worker: Yes, they should. They must keep their offices and workshops clean and wash their hands regularly.}

\eng{Student: Thank you very much for your advice, Madam. I will share it with my classmates.}

\eng{Health worker: You are welcome. Goodbye!}

\textbf{Justification} : le dialogue compte 12 répliques équilibrées ; il utilise les formules fonctionnelles (\eng{May I ask...?}, \eng{Could you tell me...?}, \eng{I see}, \eng{Thank you for your advice}), les structures de conseil (\eng{should}, \eng{must}, \eng{the best protection is to...}) et le lexique du thème. Les questions fermées (\eng{Should workers also...?}) appellent bien la réponse \eng{Yes, they should.}

\textbf{Conclusion} : un bon dialogue d'examen est équilibré, poli, structuré en 4 temps et riche en lexique thématique correct.
\end{exoresolu}

\courssub{Fiche méthode 5 — Prononciation, accentuation et intonation}

\begin{methode}[Prononcer correctement le lexique de la santé et soigner l'intonation]
\begin{enumerate}
\item \textbf{Accentuer le bon mot (accent de phrase)} : dans une phrase d'information, l'accent tonique tombe sur le(s) mot(s) nouveau(x) et important(s) (mots de contenu : noms, verbes, adjectifs, adverbes). Les mots-outils (auxiliaires, pronoms, articles, prépositions) sont faibles (forme faible / schwa).
\eng{You should WASH your HANDS with SOAP.} « Tu devrais te laver les mains avec du savon. »
\item \textbf{Choisir la bonne intonation} :
\begin{itemize}
\item Voix montante ↗ : questions fermées (yes/no), politesse, listes non terminées.
\item Voix descendante ↘ : questions en \eng{wh-} (\eng{who, what, where, when, why, how}), affirmations, conseils, ordres, formules de clôture.
\end{itemize}
\item \textbf{Soigner les mots difficiles du thème} (prononciation approximative en lettres françaises, accent tonique souligné) :
\begin{itemize}
\item \eng{disease} /dɪˈziːz/ : « di-\textbf{ZIIZ} » (pas « di-siz »)
\item \eng{stomach} /ˈstʌmək/ : « \textbf{STA}-mak »
\item \eng{cough} /kɒf/ : « \textbf{KOF} »
\item \eng{fever} /ˈfiːvə/ : « \textbf{FII}-veur »
\item \eng{hygiene} /ˈhaɪdʒiːn/ : « \textbf{HAÏ}-djin »
\item \eng{vaccination} /ˌvæksɪˈneɪʃn/ : « vak-si-\textbf{NEI}-cheun »
\item \eng{cholera} /ˈkɒlərə/ : « \textbf{KO}-le-ra »
\item \eng{mosquito} /məˈskiːtəʊ/ : « mos-\textbf{KII}-too »
\item \eng{diarrhoea} /ˌdaɪəˈriːə/ : « da-ya-\textbf{RII}-a »
\item \eng{environment} /ɪnˈvaɪrənmənt/ : « in-\textbf{VAÏ}-ron-meunt »
\end{itemize}
\item \textbf{Marquer une pause} après chaque groupe de sens (chunking) : \eng{To prevent cholera / you should drink boiled water / and wash your hands with soap.}
\item \textbf{S'entraîner en écho (shadowing)} : répéter chaque réplique du dialogue modèle après le professeur (ou l'audio), en imitant l'intonation et le rythme avant de jouer la scène.
\end{enumerate}
\end{methode}

\begin{exoresolu}[Prononciation et intonation — type examen]
\textbf{Énoncé} — 1. Indicate with an arrow (rising ↗ or falling ↘) the intonation of each sentence. 2. Underline the stressed word(s) in sentence (c). 3. Give the pronunciation in French letters of: \eng{cough}, \eng{disease}, \eng{fever}.

a) \eng{Do you sleep under a mosquito net?}

b) \eng{What are the symptoms of typhoid fever?}

c) \eng{You must boil your drinking water.}

d) \eng{Thank you for your advice, nurse.}
\tcblower
\textbf{Solution détaillée}

\textbf{1. Intonation}

a) \eng{Do you sleep under a mosquito net? ↗} — question fermée (réponse yes/no) : la voix monte à la fin.

b) \eng{What are the symptoms of typhoid fever? ↘} — question ouverte en \eng{what} : la voix descend à la fin.

c) \eng{You must boil your drinking water. ↘} — affirmation/conseil : intonation descendante.

d) \eng{Thank you for your advice, nurse. ↘} — formule de clôture : voix descendante, avec une légère montée possible sur \eng{advice} avant la chute finale.

\textbf{2. Accentuation}

Dans \eng{You must BOIL your DRINKING WATER}, les mots accentués sont les mots porteurs de sens (mots de contenu) : le verbe \eng{boil} et le groupe nominal \eng{drinking water}. Les mots-outils (\eng{you}, \eng{must} — ici semi-auxiliaire souvent faible, \eng{your}) sont prononcés faiblement (forme faible). En anglais, ce sont les mots de contenu qui portent le rythme de la phrase (langue à accent d'intensité).

\textbf{3. Prononciation en lettres françaises}

\eng{cough} : « \textbf{KOF} » ; \eng{disease} : « di-\textbf{ZIIZ} » ; \eng{fever} : « \textbf{FII}-veur ».

\textbf{Conclusion} : à l'oral, l'examinateur évalue autant le contenu que la clarté : une intonation montante pour les questions fermées, descendante pour les questions en \eng{wh-} et les affirmations, et un accent sur les mots importants rendent ton échange naturel et compréhensible.
\end{exoresolu}

\begin{experience}[Activité de classe — Simulation "Journée Santé à l'École"]
\textbf{Objectif} : Réinvestir le lexique, les formules fonctionnelles et la prononciation dans une mise en situation réaliste.

\textbf{Déroulement (20 min)} :
\begin{enumerate}
\item \textbf{Préparation (5 min)} : Par binômes. Tirer au sort une carte "Rôle A" et "Rôle B".
\begin{itemize}
\item Carte 1 : A = \eng{School nurse}, B = \eng{Class prefect} (Thème : choléra / hygiène cantine).
\item Carte 2 : A = \eng{Community health worker}, B = \eng{Market vendor} (Thème : paludisme / eau stagnante).
\item Carte 3 : A = \eng{Company safety officer}, B = \eng{Factory worker} (Thème : hygiène lieu de travail / vaccination).
\end{itemize}
\item \textbf{Répétition guidée (5 min)} : Chaque binôme rédige un squelette (4 questions / 4 réponses) sur une fiche. Le professeur circule, corrige la grammaire (auxiliaires, ordre des mots) et la prononciation (mots-clés).
\item \textbf{Représentation / Enregistrement (10 min)} : Les binômes jouent la scène devant la classe ou s'enregistrent sur téléphone (application OnBuch+ ou dictaphone). Les autres élèves notent sur une grille d'évaluation par les pairs : \eng{Content / Vocabulary} (5 pts), \eng{Grammar / Structures} (5 pts), \eng{Pronunciation / Intonation} (5 pts), \eng{Interaction / Fluency} (5 pts).
\item \textbf{Retour collectif} : Le professeur commente 2-3 points forts et 2 points à améliorer relevés chez plusieurs groupes.
\end{enumerate}
\end{experience}

\begin{savaistu}[Culture \& Contexte : Santé publique au Cameroun et dans le monde anglophone]
\begin{itemize}
\item \textbf{Programme Élargi de Vaccination (PEV) au Cameroun} : Gratuit pour les enfants de 0 à 11 mois (BCG, Polio, Penta, Rougeole, Fièvre jaune, ROTA, Pneumo). En Terminale, vous pouvez être concernés par la vaccination contre le \eng{HPV} (Papillomavirus) pour les filles, ou le tétanos.
\item \textbf{Paludisme} : Première cause de consultation et d'hospitalisation. La distribution gratuite de \eng{LLINs} (Long-Lasting Insecticidal Nets / MILD) a lieu tous les 3 ans. \eng{Indoor Residual Spraying} (Pulvérisation intra-domiciliaire) dans certaines zones.
\item \textbf{Choléra} : Épidémies récurrentes (Douala, Yaoundé, Nord). Lié à l'accès à l'eau potable et à l'assainissement. \eng{Oral Rehydration Salts (ORS)} / \eng{SRO} (Sel de Réhydratation Orale) : geste qui sauve en attendant le centre de santé.
\item \textbf{Anglophonie} : Au Nigeria, au Ghana, en Afrique du Sud, les \eng{Community Health Workers} (Agents de santé communautaires) jouent un rôle clé dans l'éducation pour la santé (\eng{Health Education}) au village. Le concept de \eng{WASH} (Water, Sanitation and Hygiene) est central dans les programmes ONU / OMS.
\item \textbf{Faux ami culturel} : En anglais, \eng{"dispensary"} est vieillot ; on dit \eng{health centre} ou \eng{clinic}. \eng{"Hospital"} = hôpital (grande structure), pas simple infirmerie.
\end{itemize}
\end{savaistu}

\begin{exercice}[Entraînement personnel — À faire chez soi ou en classe]
\textbf{Exercice 1 — Vocabulaire :} Complétez les phrases avec le mot approprié de la liste : \eng{vaccinated, stagnant, boiled, symptoms, nets, outbreak, sanitation, contaminate}.
1. During the rainy season, \eng{............} water becomes a breeding ground for mosquitoes.
2. The \eng{............} of typhoid include high fever, headache, and stomach pain.
3. All children should be \eng{............} against measles and yellow fever.
4. Poor \eng{............} in the neighbourhood led to a cholera \eng{............} last year.
5. Drink only \eng{............} or treated water to avoid diseases.
6. Sleep under treated mosquito \eng{............} every night.
7. Flies can \eng{............} food if it is left uncovered.

\textbf{Exercice 2 — Grammaire :} Mettez les verbes entre parenthèses à la forme correcte (passif, modal, question indirecte).
1. Water \eng{(must / boil)} before drinking.
2. The patients \eng{(treat)} with antibiotics at the health centre yesterday.
3. Could you tell me where the nearest clinic \eng{(be)}?
4. Food \eng{(should / cover)} to protect it from flies.
5. I would like to know how malaria \eng{(spread)}.

\textbf{Exercice 3 — Expression orale / écrite :} Préparez un dialogue de 10 lignes minimum.
\textit{Situation :} Vous êtes un élève (Student A). Votre camarade (Student B) vient d'arriver de la zone rurale et ne connaît pas les mesures de prévention du choléra à l'école. Expliquez-lui les 3 règles d'or (eau, mains, latrines) en utilisant \eng{should / must / have to} et les formules de politesse.
\end{exercice}

\begin{corrige}[Exercice 1]
1. \eng{stagnant} — 2. \eng{symptoms} — 3. \eng{vaccinated} — 4. \eng{sanitation / outbreak} — 5. \eng{boiled} — 6. \eng{nets} — 7. \eng{contaminate}.
\end{corrige}

\begin{corrige}[Exercice 2]
1. \eng{must be boiled} (passif modal obligation).
2. \eng{were treated} (passif prétérit — action passée terminée).
3. \eng{is} (question indirecte : ordre sujet + verbe, pas d'auxiliaire \eng{do}).
4. \eng{should be covered} (passif modal conseil).
5. \eng{spreads} (question indirecte : sujet \eng{malaria} 3\textsuperscript{e} pers. sg → verbe + s, pas d'auxiliaire \eng{does}).
\end{corrige}

\begin{corrige}[Exercice 3 — Proposition de dialogue]
\eng{Student A: Hello! Welcome to our school. Do you know how we prevent cholera here?}
\eng{Student B: Hi. No, I don't. In my village, we just drink water from the river.}
\eng{Student A: That is risky. You must drink only boiled or treated water.}
\eng{Student B: I see. And what else?}
\eng{Student A: You should wash your hands with soap before eating and after using the latrine.}
\eng{Student B: Okay. Must we clean the latrines too?}
\eng{Student A: Yes, we have to keep them clean every day. It is very important.}
\eng{Student B: Thank you for the advice. I will follow these rules.}
\eng{Student A: You are welcome. Let's go to class now.}
\end{corrige}

\begin{attention}[Les 5 erreurs qui coûtent des points au Bac]
\begin{enumerate}
\item \textbf{Garder l'inversion dans une question indirecte.} Ne dis pas \eng{Could you tell me how can we prevent cholera?} mais \eng{Could you tell me how we can prevent cholera?} Après \eng{Could you tell me...}, ordre sujet + verbe.
\item \textbf{Oublier l'auxiliaire dans les questions directes.} Ne dis pas \eng{How cholera spreads?} mais \eng{How does cholera spread?} Toute question au présent simple (verbe plein) exige \eng{do/does} (et le verbe principal perd son -s avec \eng{does}).
\item \textbf{Confondre les faux amis du lexique médical.}
\begin{itemize}
\item \eng{a disease} = une maladie (et non « une défaise »).
\item \eng{to prevent} = empêcher, prévenir (action avant) ; \eng{to warn} = avertir, prévenir (dire avant).
\item \eng{sensible} (anglais) = raisonnable, sensé ; \eng{sensitive} (anglais) = sensible (qui ressent fort).
\item \eng{actually} = en fait, en réalité ; \eng{currently} = actuellement.
\end{itemize}
\item \textbf{Tromper de préposition (collocations verbes/adjectifs).} On dit \eng{vaccinated against} (contre), \eng{protect ... from} (de), \eng{suffer from} (de), \eng{listen to} (à), \eng{responsible for} (de). Apprends chaque mot avec sa préposition.
\item \textbf{Calquer l'ordre des mots du français et mal orthographier le lexique.}
\begin{itemize}
\item Adjectif avant nom : \eng{communicable diseases} (pas \eng{diseases communicable}).
\item Orthographe : \eng{hygiene} (un seul \eng{g}), \eng{diarrhoea} (orthographe britannique), \eng{environment} (\eng{n} avant \eng{m}), \eng{mosquito} (pas \eng{mosquitoe}), \eng{vaccination} (double \eng{c}, \eng{tion}).
\end{itemize}
\end{enumerate}
\end{attention}

\newpage

\coursec{Exercices}

\courssub{Je vérifie mes connaissances}

\begin{exercice}[QCM : les bons réflexes de communication]
Entoure la bonne réponse (une seule par question).
\begin{enumerate}
\item Pour demander poliment une information à l'infirmière, on dit :
\begin{itemize}
\item[a)] \eng{Tell me now how cholera spreads!}
\item[b)] \eng{Could you tell me how cholera spreads?}
\item[c)] \eng{How cholera spreads, yes or no?}
\end{itemize}
\item Pour montrer que l'on suit la conversation et relancer, on dit :
\begin{itemize}
\item[a)] \eng{I see. And what about hand washing?}
\item[b)] \eng{Stop talking, please.}
\item[c)] \eng{I don't care.}
\end{itemize}
\item Pour conclure poliment une interview, on dit :
\begin{itemize}
\item[a)] \eng{Goodbye, it's finished.}
\item[b)] \eng{Thank you very much for your time.}
\item[c)] \eng{I am going now.}
\end{itemize}
\item \eng{“You should boil drinking water.”} Ici, \eng{should} exprime :
\begin{itemize}
\item[a)] une obligation absolue
\item[b)] un conseil
\item[c)] une interdiction
\end{itemize}
\end{enumerate}
\end{exercice}

\begin{exercice}[Vrai ou faux, à justifier]
Réponds par \eng{True} ou \eng{False} et justifie ta réponse en une phrase en anglais.
\begin{enumerate}
\item \eng{Cholera spreads through mosquito bites.}
\item \eng{Washing hands with soap helps prevent communicable diseases.}
\item \eng{“Could you tell me...?” is more polite than a direct question.}
\item \eng{In English, the auxiliary “do” is not used in indirect questions.}
\end{enumerate}
\end{exercice}

\begin{exercice}[Vocabulaire : le bon mot]
Complète chaque phrase avec le mot qui convient, choisi dans la liste : \eng{outbreak — vaccine — hygiene — symptoms — contaminated}.
\begin{enumerate}
\item Cholera is caused by drinking \eng{\ldots\ldots} water.
\item A high fever and vomiting are common \eng{\ldots\ldots} of malaria.
\item There was a cholera \eng{\ldots\ldots} in the neighbourhood last month.
\item Good personal \eng{\ldots\ldots} protects the whole community.
\item The health centre gives every child a \eng{\ldots\ldots} against measles.
\end{enumerate}
\end{exercice}

\begin{exercice}[Grammaire : les questions indirectes]
Complète chaque question indirecte avec la forme correcte du verbe entre parenthèses.
\begin{enumerate}
\item \eng{Could you tell me how malaria \ldots\ldots?} (to spread)
\item \eng{Do you know when the vaccination campaign \ldots\ldots?} (to start)
\item \eng{Can you explain why people \ldots\ldots boil the water?} (must)
\item \eng{I'd like to know where the health centre \ldots\ldots.} (to be)
\item \eng{Could you tell me what the symptoms of typhoid \ldots\ldots?} (to be)
\end{enumerate}
\end{exercice}

\courssub{Je m'entraîne}

\begin{exercice}[Traduction]
Traduis les phrases suivantes en anglais.
\begin{enumerate}
\item Pourriez-vous me dire comment le choléra se transmet ?
\item Il faut se laver les mains avec du savon avant chaque repas.
\item Je vois. Et qu'en est-il de la vaccination dans les écoles ?
\item Merci beaucoup pour votre temps, Madame l'infirmière.
\item L'eau contaminée peut provoquer de graves maladies.
\end{enumerate}
\end{exercice}

\begin{exercice}[Des questions directes aux questions indirectes polies]
Transforme chaque question directe en question indirecte commençant par \eng{Could you tell me...?} ou \eng{Do you know...?}.
\begin{enumerate}
\item \eng{How does malaria spread?}
\item \eng{Where can we get vaccinated?}
\item \eng{When does the hygiene campaign begin?}
\item \eng{Why must we boil drinking water?}
\item \eng{What are the symptoms of typhoid fever?}
\end{enumerate}
\end{exercice}

\begin{exercice}[Dialogue à trous]
Complète le dialogue suivant avec les formules fonctionnelles de la liste : \eng{Could you tell me — I see — What about — Thank you for your time — Do you know}.

\textbf{Journalist:} Good morning, Madam. \eng{\ldots\ldots} how cholera spreads in schools?

\textbf{Nurse:} It spreads mainly through contaminated food and water.

\textbf{Journalist:} \eng{\ldots\ldots}. \eng{\ldots\ldots} hand washing? Is it really effective?

\textbf{Nurse:} Yes, it is the simplest and cheapest protection.

\textbf{Journalist:} \eng{\ldots\ldots} if the school has enough clean water points?

\textbf{Nurse:} We have three, and we plan to build two more.

\textbf{Journalist:} That is very interesting. \eng{\ldots\ldots}, Madam.

\textbf{Nurse:} You are welcome.
\end{exercice}

\begin{exercice}[Appariement : répliques et fonctions]
Associe chaque réplique (1--5) à sa fonction de communication (a--e).

\begin{center}
\begin{tabularx}{\linewidth}{|X|X|}
\hline
\rowcolor{popblueL} \textbf{Réplique} & \textbf{Fonction} \\
\hline
1. \eng{Excuse me, could you tell me how to protect my family from cholera?} & a) conclure poliment \\
\hline
2. \eng{I see. That is very useful information.} & b) demander une information \\
\hline
3. \eng{And what about vaccination? Is it free?} & c) donner une information \\
\hline
4. \eng{You should wash your hands with soap and boil drinking water.} & d) montrer que l'on suit \\
\hline
5. \eng{Thank you very much for your time.} & e) relancer la conversation \\
\hline
\end{tabularx}
\end{center}
\end{exercice}

\courssub{Je me prépare au Bac}

\begin{exercice}[Compréhension de l'écrit et langue — type Bac]
Lis le texte suivant, puis réponds aux questions en anglais, avec tes propres mots.

\begin{textebox}[Health Day at Government High School, Bafoussam]
Last Friday, Government High School Bafoussam organised its first “Health Day”. The event was prepared by the school nurse, Mrs Ngono, with the help of the English club. In the morning, the nurse gave a talk in the assembly hall. She explained that communicable diseases such as cholera, typhoid fever and malaria still kill many people in our region, although most of them can be prevented with simple habits. “Clean hands save lives,” she repeated several times.

After the talk, students visited stands where they learned how to wash their hands properly with soap, how to treat drinking water, and how to keep latrines clean. A team from the district hospital also offered free vaccination against measles. Many parents attended the event and asked questions about hygiene at home.

At the end of the day, the principal announced that hand-washing points would be installed near every classroom before the end of the term. “A healthy student learns better,” he said. The English club recorded interviews with the nurse and the doctors, which will be broadcast on a local radio station next week.
\end{textebox}

\textbf{A. Comprehension questions}
\begin{enumerate}
\item Who organised the “Health Day” at Government High School Bafoussam?
\item Name two communicable diseases mentioned in the text.
\item What did students learn at the stands?
\item What did the district hospital team offer?
\item What did the principal promise at the end of the day?
\end{enumerate}

\textbf{B. Language work}
\begin{enumerate}
\item Find in the text the opposite of: a) \eng{dirty} ; b) \eng{expensive}.
\item Turn into an indirect question: \eng{“How can we prevent cholera?”} (begin with \eng{Could you tell me...?}).
\item Complete with the correct form of the verb: \eng{The principal said that hand-washing points \ldots\ldots (to install) before the end of the term.}
\end{enumerate}
\end{exercice}

\begin{exercice}[Traduction — type Bac]
Traduis en français le passage suivant, extrait du texte ci-dessus :

\begin{textebox}[Excerpt]
“After the talk, students visited stands where they learned how to wash their hands properly with soap, how to treat drinking water, and how to keep latrines clean. A team from the district hospital also offered free vaccination against measles.”
\end{textebox}
\end{exercice}

\begin{exercice}[Expression écrite et orale — type Bac]
\textbf{Sujet A — Composition (150--180 mots).} Your school wants to prevent an outbreak of cholera. Write an article for the school magazine in which you explain how cholera spreads and give at least four pieces of advice to students and staff. Use linking words (\eng{first, then, moreover, finally}) and at least two modal verbs (\eng{should, must, can}).

\medskip

\textbf{Sujet B — Jeu de rôle noté (préparation écrite : un script d'au moins 8 répliques).} You are a journalist working for a local radio station in Douala. Interview the school nurse about preventing typhoid fever in your school. Your script must contain:
\begin{itemize}
\item a polite greeting and an introduction of yourself ;
\item at least four questions, including two indirect questions (\eng{Could you tell me...?}) ;
\item one expression to show you are following (\eng{I see.}) and one to move to a new point (\eng{What about...?}) ;
\item a polite conclusion (\eng{Thank you for your time.}).
\end{itemize}
Tu joueras ensuite ce dialogue en binôme devant la classe.
\end{exercice}

\newpage

\coursec{Corrigés des exercices}

\begin{corrige}[Exercice 1 — QCM : les bons réflexes de communication]
\begin{enumerate}
\item \textbf{b)} \eng{Could you tell me how cholera spreads?} — C'est la formule polie de demande d'information (question indirecte). La réponse a) est un ordre brutal~; la réponse c) est grammaticalement incorrecte (il manque l'auxiliaire : il faudrait \eng{How does cholera spread?}).
\item \textbf{a)} \eng{I see. And what about hand washing?} — \eng{I see} montre que l'on suit~; \eng{What about...?} relance la conversation vers un nouveau point. Les réponses b) et c) sont impolies et ferment l'échange.
\item \textbf{b)} \eng{Thank you very much for your time.} — C'est la formule de conclusion polie d'une interview. \eng{Goodbye, it's finished} et \eng{I am going now} sont abruptes et maladroites.
\item \textbf{b)} un conseil — \eng{should} + base verbale exprime le conseil ou la recommandation. L'obligation absolue s'exprime avec \eng{must}, l'interdiction avec \eng{mustn't}.
\end{enumerate}
\end{corrige}

\begin{corrige}[Exercice 2 — Vrai ou faux, à justifier]
\begin{enumerate}
\item \eng{False. Cholera spreads through contaminated food and water, not through mosquito bites.} (Le paludisme, \eng{malaria}, est transmis par les moustiques~; ne pas confondre les modes de transmission.)
\item \eng{True. Washing hands with soap kills germs and helps prevent communicable diseases like cholera and typhoid fever.}
\item \eng{True. “Could you tell me...?” is an indirect question, and indirect questions are more polite than direct questions.}
\item \eng{True. In indirect questions we use the affirmative word order, so the auxiliary “do” disappears: “Do you know how cholera spreads?”, not “how does cholera spread”.}
\end{enumerate}
\end{corrige}

\begin{corrige}[Exercice 3 — Vocabulaire : le bon mot]
\begin{enumerate}
\item \eng{Cholera is caused by drinking \textbf{contaminated} water.} («~contaminée~» ; adjectif, participe passé de \eng{to contaminate}.)
\item \eng{A high fever and vomiting are common \textbf{symptoms} of malaria.} («~symptômes~» ; pluriel car le verbe est \eng{are}.)
\item \eng{There was a cholera \textbf{outbreak} in the neighbourhood last month.} («~flambée épidémique~» ; attention, \eng{outbreak} est un seul mot.)
\item \eng{Good personal \textbf{hygiene} protects the whole community.} («~hygiène~» ; prononciation : «~haï-djine~».)
\item \eng{The health centre gives every child a \textbf{vaccine} against measles.} («~vaccin~» ; prononciation : «~vac-sine~».)
\end{enumerate}
\end{corrige}

\begin{corrige}[Exercice 4 — Grammaire : les questions indirectes]
Rappel : dans une question indirecte, on retrouve l'ordre sujet + verbe (ordre affirmatif), sans inversion ni auxiliaire \eng{do}.
\begin{enumerate}
\item \eng{Could you tell me how malaria \textbf{spreads}?} — 3\textsuperscript{e} personne du singulier au présent simple : \eng{malaria} est singulier, d'où le \eng{-s}.
\item \eng{Do you know when the vaccination campaign \textbf{starts}?} — présent simple, 3\textsuperscript{e} personne du singulier.
\item \eng{Can you explain why people \textbf{must} boil the water?} — \eng{must} est un modal invariable, suivi de la base verbale \eng{boil}.
\item \eng{I'd like to know where the health centre \textbf{is}.} — pas d'inversion : on dit \eng{where the health centre is}, jamais \eng{where is the health centre} après \eng{I'd like to know}.
\item \eng{Could you tell me what the symptoms of typhoid \textbf{are}?} — le sujet \eng{the symptoms} est pluriel, d'où \eng{are}.
\end{enumerate}
\end{corrige}

\begin{corrige}[Exercice 5 — Traduction]
\begin{enumerate}
\item \eng{Could you tell me how cholera spreads?} (ou \eng{...how cholera is transmitted?}) — question indirecte : ordre affirmatif, pas de \eng{does}.
\item \eng{We must wash our hands with soap before every meal.} (ou \eng{You must wash your hands...}) — «~il faut~» se rend par \eng{must} ou \eng{have to}~; en anglais, on précise le possesseur : \eng{our hands / your hands}.
\item \eng{I see. And what about vaccination in schools?} — «~qu'en est-il de...?~» se dit \eng{What about...?}~; pas d'article devant \eng{vaccination} (sens général).
\item \eng{Thank you very much for your time, Madam Nurse.} (ou simplement \eng{...for your time, Madam.}) — on ne traduit pas «~Madame l'infirmière~» mot à mot.
\item \eng{Contaminated water can cause serious diseases.} — \eng{can} + base verbale~; «~provoquer~» se dit \eng{cause}, pas \eng{provoke} (faux ami : \eng{provoke} = «~provoquer quelqu'un~»).
\end{enumerate}
\end{corrige}

\begin{corrige}[Exercice 6 — Des questions directes aux questions indirectes polies]
\begin{enumerate}
\item \eng{Could you tell me how malaria spreads?} — disparition de \eng{does} et retour du \eng{-s} de la 3\textsuperscript{e} personne sur \eng{spread}.
\item \eng{Do you know where we can get vaccinated?} — inversion supprimée : \eng{we can}, pas \eng{can we}.
\item \eng{Could you tell me when the hygiene campaign begins?} — \eng{does} disparaît, \eng{begin} devient \eng{begins}.
\item \eng{Do you know why we must boil drinking water?} — le modal \eng{must} se place après le sujet \eng{we}.
\item \eng{Could you tell me what the symptoms of typhoid fever are?} — ordre affirmatif : sujet \eng{the symptoms} + verbe \eng{are}.
\end{enumerate}
\textbf{Méthode :} mot interrogatif + sujet + verbe conjugué~; jamais d'auxiliaire \eng{do/does/did} dans la subordonnée.
\end{corrige}

\begin{corrige}[Exercice 7 — Dialogue à trous]
\begin{enumerate}
\item \eng{\textbf{Could you tell me} how cholera spreads in schools?} — demande polie d'information pour ouvrir l'interview.
\item \eng{\textbf{I see}. \textbf{What about} hand washing? Is it really effective?} — \eng{I see} montre que le journaliste suit~; \eng{What about...?} introduit un nouveau point.
\item \eng{\textbf{Do you know} if the school has enough clean water points?} — question indirecte fermée (réponse oui/non) introduite par \eng{if}.
\item \eng{That is very interesting. \textbf{Thank you for your time}, Madam.} — conclusion polie de l'interview.
\end{enumerate}
Dialogue complet vérifié : les cinq formules de la liste sont utilisées chacune une fois, dans l'ordre logique de l'échange (demander — suivre — relancer — demander — conclure).
\end{corrige}

\begin{corrige}[Exercice 8 — Appariement : répliques et fonctions]
\begin{center}
\begin{tabularx}{\linewidth}{|X|X|}
\hline
\rowcolor{popblueL} \textbf{Réplique} & \textbf{Fonction} \\
\hline
1. \eng{Excuse me, could you tell me how to protect my family from cholera?} & \textbf{b)} demander une information \\
\hline
2. \eng{I see. That is very useful information.} & \textbf{d)} montrer que l'on suit \\
\hline
3. \eng{And what about vaccination? Is it free?} & \textbf{e)} relancer la conversation \\
\hline
4. \eng{You should wash your hands with soap and boil drinking water.} & \textbf{c)} donner une information \\
\hline
5. \eng{Thank you very much for your time.} & \textbf{a)} conclure poliment \\
\hline
\end{tabularx}
\end{center}
Justifications : la réplique 1 contient la formule de demande polie \eng{Could you tell me...?}~; la 2 contient \eng{I see}, marqueur d'écoute active~; la 3 ouvre un nouveau sujet avec \eng{What about...?}~; la 4 donne des conseils avec \eng{should}~; la 5 remercie, ce qui clôt l'échange.
\end{corrige}

\begin{corrige}[Exercice 9 — Compréhension de l'écrit et langue — type Bac]
\textbf{A. Comprehension questions} (réponses attendues, avec ses propres mots)
\begin{enumerate}
\item \eng{The “Health Day” was organised by the school nurse, Mrs Ngono, with the help of the English club.}
\item \eng{The text mentions cholera, typhoid fever, malaria and measles.} (Deux suffisent, par exemple \eng{cholera and malaria}.)
\item \eng{At the stands, students learned how to wash their hands properly with soap, how to treat drinking water and how to keep latrines clean.}
\item \eng{The district hospital team offered free vaccination against measles.}
\item \eng{The principal promised that hand-washing points would be installed near every classroom before the end of the term.}
\end{enumerate}
\textbf{B. Language work}
\begin{enumerate}
\item a) L'opposé de \eng{dirty} est \eng{clean} («~Clean hands save lives~», \eng{keep latrines clean}). b) L'opposé de \eng{expensive} est \eng{free} (\eng{free vaccination}).
\item \eng{Could you tell me how we can prevent cholera?} — ordre affirmatif : \eng{we can}, sans inversion.
\item \eng{The principal said that hand-washing points \textbf{would be installed} before the end of the term.} — discours rapporté au passé : \eng{will be installed} recule à \eng{would be installed} (passif : auxiliaire \eng{be} + participe passé \eng{installed}).
\end{enumerate}
\textbf{Barème indicatif (sur 20) :} A. 5 $\times$ 2 = 10 points (réponse correcte 1 pt, anglais correct 1 pt) ; B. 1) 2 pts, 2) 4 pts, 3) 4 pts.
\textbf{Points de vigilance :} répondre par des phrases complètes, pas par un seul mot~; ne pas recopier de longs passages du texte~; respecter la concordance des temps dans le discours rapporté.
\end{corrige}

\begin{corrige}[Exercice 10 — Traduction — type Bac]
\textbf{Proposition de traduction :}

«~Après l'exposé, les élèves ont visité des stands où ils ont appris comment se laver les mains correctement avec du savon, comment traiter l'eau de boisson et comment garder les latrines propres. Une équipe de l'hôpital de district a également offert la vaccination gratuite contre la rougeole.~»

\textbf{Points de vigilance :}
\begin{itemize}
\item \eng{the talk} : «~l'exposé / la causerie~» (et non «~la conversation~»).
\item \eng{how to wash...} : se traduit par «~comment se laver...~» (infinitif français).
\item \eng{drinking water} : «~l'eau de boisson~» / «~l'eau potable~».
\item \eng{free vaccination} : «~la vaccination gratuite~»~; \eng{free} signifie ici «~gratuit~», pas «~libre~».
\item \eng{measles} : «~la rougeole~» (nom singulier en français, malgré le \eng{-s} anglais).
\item \eng{district hospital} : «~hôpital de district~».
\end{itemize}
\textbf{Barème indicatif (sur 5) :} 1 pt par segment correctement traduit (5 segments), moins 0,5 pt par contresens ou faute de français grave.
\end{corrige}

\begin{corrige}[Exercice 11 — Expression écrite et orale — type Bac]
\textbf{Sujet A — Proposition de rédaction modèle (environ 165 mots) :}

\begin{textebox}[Model article — Stop Cholera Before It Starts!]
Cholera is a dangerous communicable disease, but we can prevent it. It spreads mainly through contaminated water and food: when people drink dirty water or eat food touched with unwashed hands, the germs enter the body and cause vomiting and diarrhoea.

How can we protect our school? First, everyone must wash their hands with soap after using the latrines and before every meal. Then, we should drink only boiled or treated water; students should never drink from unknown sources. Moreover, the school kitchen must keep food covered and cook it well. Finally, anyone who feels sick should go immediately to the school nurse instead of staying in class.

Clean hands save lives. If every student and every teacher follows these simple rules, our school will stay healthy and no outbreak will stop our studies. Let us act together, today!
\end{textebox}

\textbf{Barème indicatif (sur 10) :} respect du sujet et du format article 2 pts~; contenu (mode de transmission + au moins 4 conseils) 3 pts~; connecteurs et modaux exigés 2 pts~; correction grammaticale et vocabulaire 2 pts~; longueur respectée 1 pt.

\textbf{Points de vigilance :} \eng{informations} et \eng{advices} n'existent pas (\eng{advice} est indénombrable)~; \eng{must} et \eng{should} sont suivis de la base verbale sans \eng{to}~; ne pas écrire \eng{wash your hands} sans possesseur.

\medskip
\textbf{Sujet B — Script modèle (jeu de rôle) :}

\begin{textebox}[Model interview script]
\textbf{Journalist:} Good morning, Madam. My name is Anne Biya and I work for Douala City Radio. Thank you for welcoming me.

\textbf{Nurse:} Good morning. You are welcome.

\textbf{Journalist:} Could you tell me what typhoid fever is exactly?

\textbf{Nurse:} It is a serious disease caused by germs in contaminated food and water.

\textbf{Journalist:} I see. And how does it spread in a school like ours?

\textbf{Nurse:} Mainly when students drink untreated water or buy food from dirty stalls.

\textbf{Journalist:} Do you know if our school water points are safe?

\textbf{Nurse:} Yes, they are tested every month, and the water is chlorinated.

\textbf{Journalist:} What about vaccination? Could you tell me where students can get vaccinated?

\textbf{Nurse:} They can get vaccinated free of charge at the district hospital.

\textbf{Journalist:} That is very useful. What final advice would you give to our students?

\textbf{Nurse:} Wash your hands with soap, drink safe water and avoid street food.

\textbf{Journalist:} Thank you very much for your time, Madam.

\textbf{Nurse:} You are welcome.
\end{textebox}

\textbf{Barème indicatif (sur 10) :} salutation et présentation 1 pt~; 4 questions dont 2 indirectes correctes 3 pts~; marqueurs \eng{I see} et \eng{What about...?} 2 pts~; conclusion polie 1 pt~; correction de l'anglais et prononciation à l'oral 3 pts.

\textbf{Points de vigilance :} dans \eng{Could you tell me what typhoid fever is?}, pas d'inversion~; articuler clairement et faire monter la voix à la fin des questions fermées (\eng{Do you know if...?})~; ne pas lire le script les yeux baissés pendant le jeu de rôle.
\end{corrige}

\newpage

\coursec{Fiche bilan}

\begin{popfigure}
\begin{tikzpicture}[mindmap, grow cyclic, every node/.style={concept, font=\small},
  root concept/.append style={concept color=popblue, font=\bfseries\small, minimum size=2.6cm},
  level 1 concept/.append style={sibling angle=60, level distance=3.6cm, minimum size=2.2cm}]
\node[root concept]{Preventing communicable diseases}
  child[concept color=poporange]{ node[align=center]{Key vocabulary\\ health \& hygiene} }
  child[concept color=popgreen]{ node[align=center]{Asking for\\ information} }
  child[concept color=poppurple]{ node[align=center]{Giving\\ information} }
  child[concept color=poppink]{ node[align=center]{Keeping the\\ talk going} }
  child[concept color=popgold]{ node[align=center]{Modal verbs\\ must, should} }
  child[concept color=popblue!70]{ node[align=center]{Pronunciation\\ \& intonation} };
\end{tikzpicture}
\legende{Carte mentale de la leçon : échanger des informations sur la prévention des maladies transmissibles.}
\end{popfigure}

\begin{bilan}
\textbf{1. Lexique clé (à connaître par cœur)}
\begin{itemize}
  \item \eng{a communicable disease} : une maladie transmissible ; \eng{an outbreak} : une épidémie ; \eng{to spread} : se propager ; \eng{germs} : les microbes.
  \item \eng{to wash one's hands} : se laver les mains ; \eng{to cover one's mouth} : se couvrir la bouche ; \eng{to sneeze} : éternuer ; \eng{a cough} : une toux.
  \item \eng{a vaccine / to get vaccinated} : un vaccin / se faire vacciner ; \eng{a symptom} : un symptôme ; \eng{a fever} : de la fièvre.
  \item \eng{the school nurse / the infirmary} : l'infirmière scolaire / l'infirmerie ; \eng{a health worker} : un agent de santé ; \eng{a health campaign} : une campagne de santé.
  \item \eng{to avoid crowded places} : éviter les endroits bondés ; \eng{to disinfect} : désinfecter ; \eng{safe drinking water} : l'eau potable.
\end{itemize}

\textbf{2. Formules fonctionnelles de l'échange d'informations}
\begin{center}
\begin{tabularx}{\linewidth}{|X|X|}\hline
\rowcolor{popblueL} Fonction & Formules anglaises \\ \hline
Demander une information & \eng{Could you tell me how cholera spreads?} ; \eng{What should we do to prevent malaria?} \\ \hline
Donner une information & \eng{Well, the best way is to...} ; \eng{According to the nurse, we must...} \\ \hline
Relancer le dialogue & \eng{Really? And what about...?} ; \eng{Could you explain that again?} ; \eng{Anything else?} \\ \hline
Vérifier la compréhension & \eng{Do you mean that...?} ; \eng{So, if I understand correctly,...} \\ \hline
Conclure & \eng{Thank you very much for the information.} ; \eng{That's very useful, thanks.} \\ \hline
\end{tabularx}
\end{center}

\textbf{3. Grammaire à connaître par cœur}
\begin{itemize}
  \item Les modaux de conseil et d'obligation : \eng{should / shouldn't} (conseil), \eng{must / mustn't} (obligation / interdiction), \eng{have to} (obligation extérieure). \eng{We should wash our hands regularly.} « Nous devrions nous laver les mains régulièrement. »
  \item L'impératif pour les consignes d'hygiène : \eng{Cover your mouth when you cough.} « Couvre-toi la bouche quand tu tousses. » Négatif : \eng{Don't touch your face.} « Ne te touche pas le visage. »
  \item Les questions indirectes : \eng{Could you tell me how it spreads?} (jamais \eng{how does it spread} après \eng{Could you tell me}).
  \item Le présent simple pour les vérités générales : \eng{Malaria is transmitted by mosquitoes.} « Le paludisme est transmis par les moustiques. »
\end{itemize}

\textbf{4. Méthodes express}
\begin{itemize}
  \item Structure d'un échange : saluer $\rightarrow$ poser une question $\rightarrow$ écouter et relancer $\rightarrow$ reformuler $\rightarrow$ remercier et conclure.
  \item Pour relancer : reprendre un mot-clé de la réponse + \eng{what about...?} ou \eng{why?}.
  \item À l'oral : monter le ton à la fin des questions en \eng{yes/no}, le baisser pour les questions en \eng{wh-}.
\end{itemize}
\end{bilan}

\begin{aretenir}[Checklist avant le Bac]
\begin{itemize}
  \item $\square$ Je connais au moins 15 mots de vocabulaire sur les maladies transmissibles et l'hygiène.
  \item $\square$ Je sais demander une information poliment avec \eng{Could you tell me...?}.
  \item $\square$ Je sais donner un conseil avec \eng{should} et une obligation avec \eng{must}.
  \item $\square$ Je sais formuler une consigne d'hygiène à l'impératif (\eng{Wash your hands!}).
  \item $\square$ Je maîtrise l'ordre des mots dans les questions indirectes.
  \item $\square$ Je sais relancer un dialogue avec \eng{What about...?} et \eng{Could you explain that again?}.
  \item $\square$ Je sais vérifier ma compréhension avec \eng{Do you mean that...?}.
  \item $\square$ Je sais conclure poliment un échange (\eng{Thank you for the information.}).
  \item $\square$ Je fais la différence de prononciation entre \eng{disease} (« di-ziz ») et \eng{decease}, et je prononce \eng{cough} « kof ».
  \item $\square$ Je place correctement l'intonation : montante pour les questions fermées, descendante pour les questions en \eng{wh-}.
  \item $\square$ Je peux mener un dialogue de 8 à 10 répliques à l'infirmerie ou dans la communauté sans hésitation majeure.
  \item $\square$ J'évite les faux amis : \eng{to catch a disease} (attraper), \eng{sensitive} (sensible) $\neq$ \eng{sensible} (raisonnable).
\end{itemize}
\end{aretenir}

\findecours

\end{document}
